% ============ 第 V 部分: \index{Gelfand 对}Gelfand 对 (Gelfand Pairs) ============
% 原书 V.1 - V.8 (原书 pp.231-310)
% 编号: V.1.1 (节级) / V.3.1.1 (子节级) / V.8.1.1 (习题)

在本章中, $G$ 总是一个局部紧群, $K$ 是 $G$ 的紧子群. 我们固定 $G$ 上的左 Haar 测度 $\mu = \mu_G$ 和 $K$ 上的规范化 Haar 测度 $\mu_K$.

\section{不变与双不变函数}

\begin{definition}[左/右/双不变函数]
\label{def:V.1.1}
若对每个 $x \in K$ 有 $L_xf = f$ (相应地 $R_xf = f$, 相应地 $L_xf = R_xf = f$), 则称 $G$ 上的函数 $f$ 被 $K$ 左不变 (left invariant) (相应地右不变 (right invariant), 相应地双不变 (bi-invariant)).
若 $F(G)$ 是 $G$ 上的函数空间 (例如 $\mathcal{C}_c(G)$), 记 $F(K\backslash G)$ (相应地 $F(G/K)$, 相应地 $F(K\backslash G/K)$) 为其左不变 (相应地右不变, 相应地双不变) 函数的子空间.
\end{definition}

设 $\Delta_G$ 是 $G$ 的模函数. 由于 $K$ 紧, 有 $\Delta_G|_K = 1$, 所以可以使用习题 \ref{ex:I.5.3.5} 的结果. 特别地:

\begin{proposition}[平均化算子]
\label{prop:V.1.2}
设 $f \in \mathcal{C}(G)$, 定义两个函数 $f^K \colon G \to \CC$ 和 ${}^Kf \colon G \to \CC$ 为
\[
f^K(x) = \int_K f(xk)\,\dd k
\]
且
\[
{}^Kf(x) = \int_K f(kx)\,\dd k.
\]
则 $f^K$ 右不变, ${}^Kf$ 左不变.
\end{proposition}

\begin{proposition}[商测度的存在性]
\label{prop:V.1.3}
存在 $G/K$ (相应地 $K\backslash G$) 上唯一的正则 Borel 测度 $\mu_{G/K}$ (相应地 $\mu_{K\backslash G}$) 使得对每个 $f \in \mathcal{C}_c(G)$,
\[
\int_G f(x)\,\dd x = \int_{G/K}f^K(x)\,\dd\mu_{G/K}(x)
\]
(相应地
\[
\int_G f(x)\,\dd x = \int_{K\backslash G}{}^Kf(x)\,\dd\mu_{K\backslash G}(x)).
\]
\end{proposition}

\begin{definition}[双平均化]
\label{def:V.1.4}
若 $f$ 是 $G$ 上的连续函数, 记
\[
{}^Kf^K = {}^K(f^K) = ({}^Kf)^K.
\]
换句话说, 这是由
\[
{}^Kf^K(x) = \int_{K \times K}f(kxk')\,\dd k\,\dd k'
\]
定义的 $G$ 上的连续函数.
\end{definition}

注意 ${}^Kf^K$ 显然是双不变函数.

\begin{proposition}[不变性的刻画]
\label{prop:V.1.5}
设 $f \in \mathcal{C}(G)$. 则 $f$ 左不变 (相应地右不变, 相应地双不变) 当且仅当 $f = {}^Kf$ (相应地 $f = f^K$, 相应地 $f = {}^Kf^K$).
\end{proposition}

\begin{proof}
这由命题 \ref{prop:V.1.2} 和 $K$ 上的测度规范化这一事实立即得出.
\end{proof}

\begin{lemma}[双平均化的积分]
\label{lem:V.1.6}
对每个 $f \in \mathcal{C}_c(G)$,
\[
\int_G {}^Kf^K(x)\,\dd x = \int_G f(x)\,\dd x.
\]
\end{lemma}

\begin{proof}
有
\[
\int_G {}^Kf^K(x)\,\dd x = \int_{G \times K^2}f(kxk')\,\dd x\,\dd k\,\dd k' = \int_G f(x)\,\dd x,
\]
因为对所有 $k, k' \in K$,
\[
\int_G f(kxk')\,\dd x = \Delta(k')\int_G f(x)\,\dd x = \int_G f(x)\,\dd x
\]
(由命题 \ref{prop:I.2.8}).
\end{proof}

\begin{proposition}[投影算子]
\label{prop:V.1.7}
设 $(\pi, V)$ 是 $G$ 的酉表示, $P_K \colon V \to V$ 是 $V^K$ 上的正交投影. 则对每个 $v \in V$,
\[
P_K(v) = \int_K \pi(k)(v)\,\dd k.
\]
而且, 若 $f \in \mathcal{C}_c(G)$, $v \in V$, 则 $\pi(f)(P_K(v)) = \pi(f^K)(v)$ 且 $P_K(\pi(f)(v)) = \pi({}^Kf)(v)$. 特别地:
\begin{enumerate}
\item[(i)] 若 $f \in \mathcal{C}_c(G)$, $v \in V^K$, 则 $\pi(f)(v) = \pi(f^K)(v)$.
\item[(ii)] 若 $f \in \mathcal{C}_c(K\backslash G)$, $v \in V$, 则 $\pi(f)(v) \in V^K$.
\end{enumerate}
(回忆 $\pi \colon L^1(G) \to \End(V)$ 在定理 \ref{thm:I.4.2.6} 中定义.)
\end{proposition}

\begin{proof}
设 $v \in V$. 积分 $w := \int_K\pi(k)(v)\,\dd k$ 的存在性由习题 \ref{ex:I.5.6.2} 得出. 若 $x \in K$, 则
\[
\pi(x)(w) = \int_K\pi(xk)(v)\,\dd k = \int_K\pi(k)(v)\,\dd k = w,
\]
所以 $w \in V^K$. 又, 若 $w' \in V^K$, 则
\[
\langle w, w' \rangle = \int_K\langle \pi(k)(v), w' \rangle\,\dd k = \int_K\langle v, \pi(k)^{-1}(w') \rangle\,\dd k = \int_K\langle v, w' \rangle\,\dd k = \langle v, w' \rangle.
\]
所以 $w$ 是 $v$ 在 $V^K$ 上的正交投影.
现在我们证明最后一个结论. 设 $f \in \mathcal{C}_c(G)$, $v \in V$. 则:
\begin{align*}
\pi(f^K)(v) &= \int_G f^K(x)\pi(x)(v)\,\dd x = \int_{G \times K}f(xk)\pi(x)(v)\,\dd x\,\dd k \\
&= \int_{G \times K}f(x)\pi(x)\pi(k)^{-1}(v)\,\dd x\,\dd k \\
&= \int_{G \times K}f(x)\pi(x)\pi(k)(v)\,\dd x\,\dd k \quad \text{($K$ 幺模)} \\
&= \pi(f)(P_K(v)).
\end{align*}
另一方面:
\begin{align*}
P_K(\pi(f)(v)) &= \int_{K \times G}f(x)\pi(kx)(v)\,\dd k\,\dd x \\
&= \int_{K \times G}f(k^{-1}x)\pi(x)(v)\,\dd k\,\dd x \\
&= \int_{K \times G}f(kx)\pi(x)(v)\,\dd k\,\dd x \\
&= \pi({}^Kf)(v).
\end{align*}
\end{proof}

同样的证明给出:

\begin{proposition}[卷积与平均化]
\label{prop:V.1.8}
设 $f, g \in \mathcal{C}_c(G)$. 则
\[
{}^K(f * g) = ({}^Kf) * g \quad \text{且} \quad (f * g)^K = f * (g^K).
\]
特别地, 若 $f$ 和 $g$ 双不变, 则 $f * g$ 也双不变, 所以 $\mathcal{C}_c(K\backslash G/K)$ 对卷积乘积是 $\mathcal{C}_c(G)$ 的子代数.
\end{proposition}

\begin{remark}[$L^p$ 空间的情形]
\label{rem:V.1.9}
设 $L^p(K\backslash G/K)$ 是 $L^p(G)$ 中双不变函数的子空间. 则若 $1 \leq p < +\infty$, $f \in L^1(K\backslash G/K)$, $g \in L^p(K\backslash G/K)$, 它们的卷积乘积 $f * g$ 在 $L^p(K\backslash G/K)$ 中. 这在定义 $f * g$ 的公式 (见命题 \ref{prop:I.4.1.3}) 上显然; 事实上,
\[
f * g(x) = \int_G f(y)g(y^{-1}x)\,\dd x = \int_G f(xy^{-1})g(y)\,\dd y
\]
(第一个公式表明 $f * g$ 右不变, 第二个表明 $f * g$ 左不变).
特别地, $L^1(G)$ 的子空间 $L^1(K\backslash G/K)$ 是子代数, 且若 $G$ 紧, 对 $L^2$ 空间有类似结果.
\end{remark}

\begin{remark}[用测度代数理解]
\label{rem:V.1.10}
如果我们把卷积乘积和表示 $\pi$ 延拓到 $G$ 上 Radon 测度的空间 $\mathcal{M}(G)$ 上, 这一切更容易记住. (见注记 \ref{rem:I.4.1.6}.) 可以通过把 $\mu_K$ 等同于 Radon 测度 $\mathcal{C}_c(G) \to \CC$, $f \longmapsto \int_K f(x)\,\dd\mu_K(x)$ 把它看成 $\mathcal{M}(G)$ 的元素. 则 $\mu_K * \mu_K = \mu_K$, $f^K = f * \mu_K$, ${}^Kf = \mu_K * f$, $P_K = \pi(\mu_K)$, 所以例如命题 \ref{prop:V.1.7} 的最后部分只是由 $\pi$ 是 $*$-同态得出.
\end{remark}

\section{Gelfand 对的定义}

\begin{definition}[Gelfand 对 (Gelfand pair)]
\label{def:V.2.1}
若代数 $\mathcal{C}_c(K\backslash G/K)$ 对卷积乘积交换, 则称 $(G, K)$ 为 Gelfand 对.
\end{definition}

\begin{remark}[定义的等价形式]
\label{rem:V.2.2}
若 $p \in [1, +\infty)$, $f \in L^p(K\backslash G/K)$, $g \in \mathcal{C}_c(G)$, 则
\begin{align*}
\|{}^Kg^K - {}^Kf^K\|_p^p &= \int_G\Bigl|f(x) - \int_{K \times K}g(kxk')\,\dd k\,\dd k'\Bigr|^p\,\dd x \\
&= \int_G\Bigl|\int_{K \times K}(f(kxk') - g(kxk'))\,\dd k\,\dd k'\Bigr|^p\,\dd x.
\end{align*}
所以, 由 Minkowski 不等式 (见习题 \ref{ex:I.5.6.7}),
\[
\|{}^Kf^K - {}^Kg^K\|_p \leq \int_{K \times K}\|L_kR_{k'}f - L_kR_{k'}g\|_p\,\dd k\,\dd k' = \|f - g\|_p.
\]
由于 $\mathcal{C}_c(G)$ 在 $L^p(G)$ 中稠密, $L^p(K\backslash G/K)$ 的每个函数都可以由 $\mathcal{C}_c(G)$ 的元素逼近, 从而由上面的计算, 可以由 $\mathcal{C}_c(K\backslash G/K)$ 的元素逼近. 换句话说, 空间 $\mathcal{C}_c(K\backslash G/K)$ 在 $L^p(K\backslash G/K)$ 中稠密. 所以在 Gelfand 对的定义中, 可以把条件``$\mathcal{C}_c(K\backslash G/K)$ 对卷积乘积交换''换成条件``$L^1(K\backslash G/K)$ 对卷积乘积交换'' (或对紧 $G$, 换成``$L^2(K\backslash G/K)$ 对卷积乘积交换'').
\end{remark}

\begin{example}[Gelfand 对的例子]
\label{exa:V.2.3}
若 $G$ 交换, 则 $(G, \{1\})$ 是 Gelfand 对.
以下是其他例子 (但我们还不证明它们是 Gelfand 对):
\begin{itemize}
\item $(\SO(n+1), \SO(n))$, 其中 $\SO(n)$ 通过把 $x \in \SO(n)$ 送到 $(n+1) \times (n+1)$ 矩阵 $\begin{pmatrix} x & 0 \\ 0 & 1 \end{pmatrix}$ 等同于 $\SO(n+1)$ 的子群;
\item $(S_{n+m}, S_n \times S_m)$;
\item $(\GL_n(\Q_p), \GL_n(\Z_p))$.
\end{itemize}
\end{example}

\begin{proposition}[Gelfand 对中的群幺模]
\label{prop:V.2.4}
设 $(G, K)$ 是 Gelfand 对. 则 $G$ 幺模.
\end{proposition}

\begin{proof}
由命题 \ref{prop:I.2.12}, 对每个 $f \in \mathcal{C}_c(G)$,
\[
\int_G f(x)\,\dd x = \int_G\Delta(x)^{-1}f(x^{-1})\,\dd x.
\]
所以只需证明对每个 $f \in \mathcal{C}_c(G)$ 有 $\int_G f(x)\,\dd x = \int_G f(x^{-1})\,\dd x$. 首先注意
\[
\int_G {}^Kf^K(x)\,\dd x = \int_G f(x)\,\dd x
\]
且
\[
\int_G {}^Kf^K(x^{-1})\,\dd x = \int_G f(x^{-1})\,\dd x,
\]
由引理 \ref{lem:V.1.6}. 所以只需证明对每个 $f \in \mathcal{C}_c(K\backslash G/K)$ 有 $\int_G f(x)\,\dd x = \int_G f(x^{-1})\,\dd x$. 固定 $f \in \mathcal{C}_c(K\backslash G/K)$. 可以找到 $g \in \mathcal{C}_c(K\backslash G/K)$ 使得 $g$ 在 $(\operatorname{supp}f) \cup (\operatorname{supp}f)^{-1}$ 上等于 $1$ (因为 $\operatorname{supp}f = K(\operatorname{supp}f)K$). 则
\[
f * g(1) = \int_G f(y)g(y^{-1})\,\dd y = \int_{\operatorname{supp}f}f(y)\,\dd y = \int_G f(y)\,\dd y
\]
且
\[
g * f(1) = \int_G g(y)f(y^{-1})\,\dd y = \int_{(\operatorname{supp}f)^{-1}}f(y^{-1})\,\dd y = \int_G f(y^{-1})\,\dd y.
\]
但 $f * g = g * f$, 因为 $(G, K)$ 是 Gelfand 对, 所以这蕴含所要的结果.
\end{proof}

下面的判据将帮助我们找到更多 Gelfand 对.

\begin{proposition}[对合判据]
\label{prop:V.2.5}
假设存在连续自同构 $\theta \colon G \to G$ 使得:
\begin{enumerate}
\item[(a)] $\theta^2 = \mathrm{id}_G$ (即 $\theta$ 是对合);
\item[(b)] 对每个 $x \in G$, $\theta(x) \in Kx^{-1}K$.
\end{enumerate}
则 $(G, K)$ 是 Gelfand 对.
\end{proposition}

\begin{proof}
考虑线性泛函 $\mathcal{C}_c(G) \to \CC$, $f \longmapsto \int_G f(\theta(x))\,\dd x$. 这是 $\mathcal{C}_c(G)$ 上的左不变正线性泛函, 所以由定理 \ref{thm:I.2.7} 中的唯一性, 存在 $c \in \R_{>0}$ 使得对每个 $f \in \mathcal{C}_c(G)$,
\[
\int_G f(\theta(x))\,\dd x = c\int_G f(x)\,\dd x.
\]
由于 $\theta^2 = \mathrm{id}_G$, 必有 $c^2 = 1$, 所以 $c = 1$.
设 $f, g \in \mathcal{C}_c(G)$. 一方面, 对每个 $x \in G$,
\begin{align*}
(f \circ \theta) * (g \circ \theta)(x) &= \int_G f(\theta(y))g(\theta(y)^{-1}\theta(x))\,\dd y \\
&= \int_G f(y)g(y^{-1}\theta(x))\,\dd y \\
&= (f * g) \circ \theta(x)
\end{align*}
(第二个等式由本证明的第一段得出). 另一方面, 对每个 $x \in G$,
\begin{align*}
(g * f)(x^{-1}) &= \int_G g(x^{-1}y)f(y)\,\dd y = \int_G f'(y^{-1})g'(yx)\,\dd y = (f' * g')(x),
\end{align*}
其中 $f'(z) = f(z^{-1})$, $g'(z) = g(z^{-1})$. (我们用 $G$ 幺模这一事实做代换 $y \longmapsto y^{-1}$.)
假设 $f$ 和 $g$ 双不变. 则由条件 (b), $f(\theta(x)) = f(x^{-1})$, $g(\theta(x)) = g(x^{-1})$, 且对 $g * f$ 有类似的等式, 因为 $g * f$ 也双不变, 所以对每个 $x \in G$,
\[
(f * g)(\theta(x)) = ((f \circ \theta) * (g \circ \theta))(x) = (f' * g')(x) = (g * f)(x^{-1}) = (g * f)(\theta(x)).
\]
由于 $\theta$ 是自同构, 这蕴含 $f * g = g * f$.
\end{proof}

\begin{example}[对合判据的应用]
\label{exa:V.2.6}
\begin{enumerate}
\item[(1)] 若 $G$ 交换, 则取 $\theta \colon x \longmapsto x^{-1}$, 所以对任何紧子群 $K$, 特别地对 $K = \{1\}$, $(G, K)$ 是 Gelfand 对.
\item[(2)] 若 $G$ 紧, 则 $(G \times G, \{(x, x) \mid x \in G\})$ 是 Gelfand 对. 事实上, 只需把上面的命题应用于 $\theta(x, y) = (y, x)$. 对每个 $(x, y) \in G \times G$, 有 $\theta(x, y) = (x, x)(x^{-1}, y^{-1})(y, y)$.
\end{enumerate}
\end{example}

\section{Gelfand 对与表示}
在本节中, 我们给出 $(G, K)$ 成为 Gelfand 对的两个表示论判据, 一个在一般情况下有效, 一个对紧 $G$ 有效.

\subsection{\texorpdfstring{Gelfand 对与 $K$-固定向量}{Gelfand 对与 K-固定向量}}

\begin{theorem}[$K$-固定向量判据]
\label{thm:V.3.1.1}
$(G, K)$ 是 Gelfand 对当且仅当: 对 $G$ 的每个不可约酉表示 $(\pi, V)$, 有 $\dim(V^K) \leq 1$.
\end{theorem}

我们需要下面这个 Gelfand-Raikov 定理的变体.

\begin{lemma}[积分分离]
\label{lem:V.3.1.2}
设 $f \in \mathcal{C}_c(G)$. 若 $f \neq 0$, 则存在 $\phi \in E(P_1)$ (见第 III.3 节) 使得 $\int_G f(x)\phi(x)\,\dd x \neq 0$.
\end{lemma}

\begin{proof}
假设对每个 $\phi \in E(P_1)$ 有 $\int_G f\phi\,\dd\mu = 0$. 由定理 \ref{thm:III.4.1}, 对每个正定型函数 $\phi$ 有 $\int_G f\phi\,\dd\mu = 0$. 由定理 \ref{thm:III.2.5}, 对 $G$ 的每个酉表示 $(\pi, V)$ 和任何 $v \in V$, 有 $\langle \pi(f)(v), v \rangle = 0$. 把它应用于 $G$ 的左正则表示, 得到对每个 $g \in L^2(G)$ 有 $\langle f * g, g \rangle_{L^2(G)} = 0$. 像定理 \ref{thm:III.5.1} 的证明那样, 看到这蕴含对所有 $g_1, g_2 \in L^2(G)$ 有 $\langle f * g_1, g_2 \rangle_{L^2(G)} = 0$. 又像该定理的证明那样, 看到对所有 $g_1, g_2 \in L^2(G)$ 有 $\langle f * g_1, g_2 \rangle_{L^2(G)} = \langle f, g_2 * \tilde{g}_1 \rangle_{L^2(G)}$, 其中 $\tilde{g}_1(x) = \overline{g_1(x^{-1})}$. 所以对所有 $g_1, g_2 \in L^2(G)$ 有 $\langle f, g_1 * g_2 \rangle_{L^2(G)} = 0$. 把它应用于 $g_1 = f$ 和 $g_2 = \psi_U$ (其中 $(\psi_U)_{U \in \mathcal{U}}$ 是近似单位元), 最终得到 $\langle f, f \rangle_{L^2(G)} = 0$, 所以 $f = 0$.
\end{proof}

我们还需要 Schur 引理的下述变体.

\begin{lemma}[交换代数的表示]
\label{lem:V.3.1.3}
设 $A$ 是交换 Banach $*$-代数, $u \colon A \to \End(V)$ 是 $A$ 在非零希尔伯特空间 $V$ 上的表示. 假设被所有 $u(x)$ ($x \in A$) 固定的 $V$ 的闭子空间只有 $\{0\}$ 和 $V$. 则 $\dim V = 1$.
\end{lemma}

\begin{proof}
由假设, 子集 $u(A)$ 满足推论 \ref{cor:II.4.4} 的假设, 所以它在 $\End(V)$ 中的中心化子等于 $\CC\mathrm{id}_V$. 但由于 $A$ 交换, $u(A)$ 的每个元素都在 $u(A)$ 的中心化子中, 所以这蕴含 $\Im(u) \subset \CC\mathrm{id}_V$. 特别地, $V$ 的每个子空间都被 $u(A)$ 的所有元素不变, 所以 $V$ 没有非平凡闭子空间, 这只有 $\dim V \leq 1$ 才可能.
\end{proof}

\begin{lemma}[$V^K$ 上的作用]
\label{lem:V.3.1.4}
设 $(\pi, V)$ 是 $G$ 的酉表示. 则对每个 $f \in L^1(K\backslash G/K)$, $\pi(f)$ 把 $V^K$ 送到自身. 若此外 $\pi$ 不可约, 则被所有 $\pi(f)$ ($f \in L^1(K\backslash G/K)$) 稳定的 $V^K$ 的闭子空间只有 $\{0\}$ 和 $V^K$.
\end{lemma}

\begin{proof}
由命题 \ref{prop:V.1.7}, 对每个 $f \in \mathcal{C}_c(K\backslash G/K)$ 和每个 $v \in V^K$, 有 $\pi(f)(v) \in V^K$. 第一个结论由 $\mathcal{C}_c(K\backslash G/K)$ 在 $L^1(K\backslash G/K)$ 中稠密得出.
为证明第二个结论, 只需证明对每个 $v \in V^K - \{0\}$, 空间 $\{\pi(f)(v) \mid f \in \mathcal{C}_c(K\backslash G/K)\}$ 在 $V^K$ 中稠密. 设 $w \in V^K$, $\varepsilon > 0$. 由于 $V$ 不可约, 空间 $\{\pi(f)(v) \mid f \in L^1(G)\}$ 在 $V$ 中稠密. 由于 $\mathcal{C}_c(G)$ 在 $L^1(G)$ 中稠密, 存在 $f \in \mathcal{C}_c(G)$ 使得 $\|\pi(f)(v) - w\| \leq \varepsilon$. 再次由命题 \ref{prop:V.1.7}, 有 $\pi(f)(v) = \pi(f^K)(v)$, 从而 $\pi({}^Kf^K)(v) = P_K(\pi(f)(v))$, 其中 $P_K$ 是 $V$ 在 $V^K$ 上的正交投影. 由于 $w \in V^K$, 得到 $\|\pi({}^Kf^K)(v) - w\| = \|P_K(\pi(f)(v)) - w\| \leq \|\pi(f)(v) - w\| \leq \varepsilon$.
\end{proof}

\begin{proof}[定理 \ref{thm:V.3.1.1} 的证明]
假设 $(G, K)$ 是 Gelfand 对. 设 $(\pi, V)$ 是 $G$ 的不可约酉表示. 由引理 \ref{lem:V.3.1.4}, $\pi$ 定义从 $L^1(K\backslash G/K)$ 到 $\End(V^K)$ 的 $*$-同态, 且被 $L^1(K\backslash G/K)$ 的所有元素稳定的 $V^K$ 的闭子空间只有 $\{0\}$ 和 $V^K$. 由于 $L^1(K\backslash G/K)$ 交换, 引理 \ref{lem:V.3.1.3} 蕴含 $\dim(V^K) \leq 1$.
我们证明其逆. 假设对 $G$ 的每个不可约酉表示 $(\pi, V)$ 有 $\dim(V^K) \leq 1$. 设非零的 $f \in \mathcal{C}_c(K\backslash G/K)$. 由引理 \ref{lem:V.3.1.2}, 存在 $\phi \in E(P_1)$ 使得 $\int_G f\phi\,\dd\mu \neq 0$. 设 $(\pi, V)$ 是 $G$ 的循环酉表示, $v \in V$ 是循环向量, 满足对每个 $x \in G$ 有 $\phi(x) = \langle \pi(x)(v), v \rangle$ (见定理 \ref{thm:III.2.5}). 则
\[
\int_G f(x)\phi(x)\,\dd x = \int_G f(x)\langle \pi(x)(v), v \rangle\,\dd x = \langle \pi(f)(v), v \rangle,
\]
所以 $\pi(f)(v) \neq 0$. 由定理 \ref{thm:III.3.2}, 表示 $(\pi, V)$ 不可约. 由引理 \ref{lem:V.3.1.4}, $V$ 的自同态 $\pi(f)$ 保持 $V^K$, 且由命题 \ref{prop:V.1.7}, 若 $w$ 是 $v$ 在 $V^K$ 上的正交投影, 则 $\pi(f)(w) = \pi(f)(v) \neq 0$. 特别地, $V$ 的子空间 $V^K$ 非零, 所以由假设 $\dim(V^K) = 1$. 所以 $\End(V^K) = \CC$, 这意味着我们找到了 $*$-同态 $u \colon \mathcal{C}_c(K\backslash G/K) \to \CC$ (把 $g$ 送到 $\pi(g)|_{V^K}$) 使得 $u(f) \neq 0$.
现在设 $f_1, f_2 \in \mathcal{C}_c(K\backslash G/K)$. 由于 $\CC$ 交换, 对每个代数态射 $u \colon \mathcal{C}_c(K\backslash G/K) \to \CC$ 有 $u(f_1 * f_2 - f_2 * f_1) = u(f_1)u(f_2) - u(f_2)u(f_1) = 0$. 由上一段, 这蕴含 $f_1 * f_2 - f_2 * f_1 = 0$, 我们完成了.
\end{proof}

\subsection{Gelfand 对与\index{无重数表示}无重数表示}

\begin{definition}[无重数表示 (multiplicity-free representation)]
\label{def:V.3.2.1}
设 $(\pi, V)$ 是 $G$ 的酉表示, 假设可以把 $V$ 写成 $V = \bigoplus_{i \in I}V_i$, 其中 $V_i$ 是 $V$ 的闭 $G$-不变子空间, 作为 $V$ 的表示不可约.\footnote{若 $G$ 紧这总成立, 见定理 \ref{thm:IV.2.1}.} 则称 $(\pi, V)$ 无重数 (multiplicity-free), 如果对 $G$ 的每个不可约酉表示 $W$, 使得 $V_i$ 与 $W$ 等价的 $i \in I$ 的集合的基数 $\leq 1$.
\end{definition}

注意群 $G$ 通过左平移作用于齐性空间 $G/K$ 上, 所以若 $x \in G$, $f$ 是 $G/K$ 上的函数, 可以定义 $L_xf$ 为 $L_xf(y) = f(x^{-1}y)$.

\begin{definition}[\index{拟正则表示}拟正则表示 (quasiregular representation)]
\label{def:V.3.2.2}
$G$ 在 $L^2(G/K)$ 上的拟正则表示是由 $x \cdot f = L_xf$ (对每个 $x \in G$ 和每个 $f \in L^2(G/K)$) 定义的表示.
\end{definition}

\begin{proposition}[拟正则表示是酉表示]
\label{prop:V.3.2.3}
上面的定义有意义, 且给出 $G$ 的酉表示.
\end{proposition}

\begin{proof}
由 $G/K$ 上测度的定义, 对每个 $f \in \mathcal{C}_c(G/K)$ 和每个 $x \in G$ 有 $\int_{G/K}f\,\dd\mu_{G/K} = \int_{G/K}L_xf\,\dd\mu_{G/K}$. 由于 $\mathcal{C}_c(G/K)$ 在 $L^2(G/K)$ 中稠密, 这蕴含算子 $L_x$ 保持 $L^2(G/K)$ 且是等距. 由命题 \ref{prop:I.3.1.10}, 只需证明对每个 $f \in L^2(G/K)$, 映射 $G \to L^2(G/K)$, $x \longmapsto L_xf$ 连续. 像命题 \ref{prop:I.3.1.13} 的证明那样, 只需对 $f \in \mathcal{C}_c(G/K)$ 证明, 此时它由命题 \ref{prop:I.1.12} 得出.
\end{proof}

\begin{remark}[另一种拟正则表示]
\label{rem:V.3.1}
若让 $G$ 通过右正则表示作用于 $L^2(G)$, 则 $L^2(G/K)$ 是 $L^2(G)$ 中 $K$-不变向量的空间. 拟正则表示正是左正则表示在 $L^2(G/K)$ 上的限制.
我们也可以定义 $L^2(K\backslash G)$ 上的拟正则表示 (这是 $K$ 通过左正则表示作用时 $L^2(G)$ 中 $K$-不变向量的空间, 它通过右正则表示得到 $G$ 的作用). 我们得到的表示酉等价于 $L^2(G/K)$ 上的拟正则表示.
\end{remark}

\begin{theorem}[紧群情形]
\label{thm:V.3.2.4}
假设 $G$ 紧. 则 $(G, K)$ 是 Gelfand 对当且仅当 $G$ 在 $L^2(G/K)$ 上的拟正则表示无重数.
又, 若 $(G, K)$ 是 Gelfand 对, 则有 $G$-等变同构
\[
L^2(G/K) \simeq \bigoplus_{\substack{(\pi, V) \in \widehat{G}\\ V^K \neq 0}}V.
\]
\end{theorem}

\begin{proof}
首先注意, 若 $K$ 通过右正则表示作用, $L^2(G/K)$ 是 $L^2(G)$ 中 $K$-不变向量的空间. Peter-Weyl 定理 (推论 \ref{cor:IV.4.3}) 说, 作为 $G \times G$ 的表示, $L^2(G)$ 同构于 $\bigoplus_{(\pi, V) \in \widehat{G}}\End(V) = \bigoplus_{(\pi, V) \in \widehat{G}}V \otimes_{\CC}V^*$ 的完备化. 所以 $L^2(G/K)$ 作为 $G$ 的表示同构于
\[
\bigoplus_{\substack{(\pi, V) \in \widehat{G}\\ V^K \neq 0}}(V^*)^{\dim(V)}
\]
的完备化.
注意对每个 $(\pi, V) \in \widehat{G}$, 表示 $V^*$ 也不可约; 这例如由推论 \ref{cor:IV.5.8}(iii) 得出, 因为 $\chi_{V^*} = \overline{\chi_V}$, 所以 $\|\chi_{V^*}\|_2 = \|\chi_V\|_2$. 所以表示 $L^2(G/K)$ 无重数当且仅当: 对 $G$ 的每个不可约酉表示 $(\pi, V)$, 有 $V^K = 0$ 或 $\dim(V^K) = 1$. 所以定理的第一个结论由定理 \ref{thm:V.3.1.1} 得出.
现在我们证明第二个结论. 我们已经看到
\[
L^2(G/K) \simeq \bigoplus_{\substack{(\pi, V) \in \widehat{G}\\ V^K \neq 0}}V^*,
\]
所以只需证明: 若 $(\pi, V)$ 是 $G$ 的有限维表示, 则 $V^K \neq 0$ 当且仅当 $(V^*)^K \neq 0$. 把定理 \ref{thm:IV.5.7} 应用于表示 $V$ 和 $V^*$ 到 $K$ 的限制, 得到
\[
\dim(V^K) = \int_K\chi_V(k)\,\dd k
\]
且
\[
\dim((V^*)^K) = \int_K\chi_{V^*}(k)\,\dd k = \int_K\chi_V(k)\,\dd k = \dim(V^K).
\]
\end{proof}

\section{\index{球函数}球函数}
在本节中, 我们假设 $(G, K)$ 是 Gelfand 对.

\begin{definition}[球函数 (spherical function)]
\label{def:V.4.1}
设 $\phi \in \mathcal{C}(K\backslash G/K)$. 若线性泛函 $\chi_{\phi} \colon \mathcal{C}_c(K\backslash G/K) \to \CC$, $f \longmapsto f * \phi(1) = \int_G f(x)\phi(x^{-1})\,\dd x$ 是乘法泛函 (其中 $\mathcal{C}_c(K\backslash G/K)$ 上的乘法是卷积乘积), 则称 $\phi$ 为球函数.
换句话说, 若 $\phi \neq 0$ 且对所有 $f, g \in \mathcal{C}_c(K\backslash G/K)$ 有 $\chi_{\phi}(f * g) = \chi_{\phi}(f)\chi_{\phi}(g)$, 则称函数 $\phi$ 球.
\end{definition}

\begin{example}[交换群的情形]
\label{exa:V.4.2}
若 $G$ 交换且 $K = \{1\}$, 则每个连续群态射 $\phi \colon G \to \CC^{\times}$ 都是球函数. 事实上, 设 $f, g \in \mathcal{C}_c(G)$. 则:
\begin{align*}
\int_G (f * g)(x)\phi(x^{-1})\,\dd x &= \int_{G \times G}f(y)g(y^{-1}x)\phi(x^{-1})\,\dd x\,\dd y \\
&= \int_{G \times G}f(y)g(z)\phi(z^{-1}y^{-1})\,\dd y\,\dd z \\
&= \Bigl(\int_G f(y)\phi(y^{-1})\,\dd y\Bigr)\Bigl(\int_G g(z)\phi(z^{-1})\,\dd z\Bigr).
\end{align*}
这些实际上是该情形下的唯一球函数. (这由下一个命题立即得出.)
\end{example}

\begin{proposition}[球函数的刻画]
\label{prop:V.4.3}
设 $\phi \in \mathcal{C}(K\backslash G/K)$. 下列条件等价:
\begin{enumerate}
\item[(i)] 函数 $\phi$ 球.
\item[(ii)] 函数 $\phi$ 非零, 且对所有 $x, y \in G$,
\[
\int_K\phi(xky)\,\dd k = \phi(x)\phi(y).
\]
\item[(iii)] 有:
\begin{enumerate}
\item[(a)] $\phi(1) = 1$;
\item[(b)] 对每个 $f \in \mathcal{C}_c(K\backslash G/K)$, 存在 $\chi(f) \in \CC$ 使得 $f * \phi = \chi(f)\phi$.
\end{enumerate}
\end{enumerate}
\end{proposition}

\begin{proof}
可以把 $\chi_{\phi}$ 用同样的公式延拓到 $\mathcal{C}_c(G)$ 上, 即 $\chi_{\phi}(f) = \int_G f(x)\phi(x^{-1})\,\dd x$. 注意对所有 $f, g \in \mathcal{C}_c(G)$,
\[
\chi_{\phi}(f * g) = \int_{G \times G}f(y)g(y^{-1}x)\phi(x^{-1})\,\dd x\,\dd y = \int_{G \times G}f(y)g(z)\phi(z^{-1}y^{-1})\,\dd y\,\dd z,
\]
所以
\[
\chi_{\phi}(f * g) - \chi_{\phi}(f)\chi_{\phi}(g) = \int_{G \times G}f(y)g(z)(\phi(z^{-1}y^{-1}) - \phi(z^{-1})\phi(y^{-1}))\,\dd y\,\dd z.
\]
设 $f, g \in \mathcal{C}_c(G)$. 把上面的计算应用于 $f' := {}^Kf^K$ 和 $g' := {}^Kg^K$, 并用 $\phi$ 的双不变性 (以及 $K$ 上测度规范化这一事实), 得到 $\chi_{\phi}(f' * g') - \chi_{\phi}(f')\chi_{\phi}(g')$ 等于
\begin{align*}
&\int_{G^2 \times K^4}f(k_1xk_2)g(k_3yk_4)(\phi(y^{-1}x^{-1}) - \phi(y^{-1})\phi(x^{-1}))\,\dd x\,\dd y\,\dd k_1\,\dd k_2\,\dd k_3\,\dd k_4 \\
&= \int_{G^2 \times K^2}f(x)g(y)(\phi(y^{-1}k_3k_2x^{-1}) - \phi(x^{-1})\phi(y^{-1}))\,\dd x\,\dd y\,\dd k_2\,\dd k_3 \\
&= \int_{G^2}f(x)g(y)\Bigl(\int_K\phi(y^{-1}kx^{-1})\,\dd k - \phi(x^{-1})\phi(y^{-1})\Bigr)\,\dd x\,\dd y.
\end{align*}
这证明 $\chi_{\phi}$ 乘法当且仅当对所有 $x, y \in G$ 有 $\int_G\phi(y^{-1}kx^{-1})\,\dd k = \phi(y^{-1})\phi(x^{-1})$. 由于 $\chi_{\phi} \neq 0$ 当且仅当 $\phi \neq 0$, 这证明 (i) 和 (ii) 等价.
假设 $\phi$ 满足 (iii) 的条件 (a) 和 (b). 则对每个 $f \in \mathcal{C}_c(K\backslash G/K)$,
\[
\chi_{\phi}(f) = f * \phi(1) = \chi(f).
\]
由于 $f \longmapsto \chi(f)$ 乘法 (由卷积乘积的结合性), 这蕴含 $\phi$ 球.
最后, 假设 $\phi$ 球. 要证明满足 (iii) 的条件 (a) 和 (b). 设 $f \in \mathcal{C}_c(K\backslash G/K)$. 则对每个 $x \in G$,
\begin{align*}
f * \phi(x) &= \int_G f(y)\phi(y^{-1}x)\,\dd y \\
&= \int_{G \times K}f(y)\phi(y^{-1}kx)\,\dd y\,\dd k \quad \text{($f$ 的左不变性)} \\
&= \int_G f(y)\phi(y)^{-1}\phi(x)\,\dd y \quad \text{(由 (ii))} \\
&= \chi_{\phi}(f)\phi(x).
\end{align*}
这证明条件 (b). 选取满足 $\chi_{\phi}(f) \neq 0$ 的 $f \in \mathcal{C}_c(K\backslash G/K)$, 把上面的等式应用于 $x = 1$, 得到 $\chi_{\phi}(f) = \chi_{\phi}(f)\phi(1)$, 所以 $\phi(1) = 1$.
\end{proof}

回忆 $L^1(G)$ 对卷积乘积和对合 $f^*(x) = \overline{f(x^{-1})}$\footnote{由于 $(G, K)$ 是 Gelfand 对, 由命题 \ref{prop:V.2.4} 群 $G$ 自动幺模, 所以不需要因子 $\Delta(x)^{-1}$.} 是 Banach $*$-代数. 我们已经看到 $L^1(K\backslash G/K)$ 是 $L^1(G)$ 的交换 Banach 子代数, 且显然它也被对合保持. 所以自然要问 $L^1(K\backslash G/K)$ 的谱是什么.
若 $\phi \in \mathcal{C}_b(K\backslash G/K)$ (注意有界性条件), 则积分 $\int_G f(x)\phi(x^{-1})\,\dd x$ 对每个 $f \in L^1(G)$ 收敛, 所以可以把 $\mathcal{C}_c(K\backslash G/K)$ 上的线性泛函 $\chi_{\phi}$ 延拓成 $L^1(K\backslash G/K)$ 上的有界线性泛函, 仍记为 $\chi_{\phi}$.

\begin{theorem}[球函数与谱]
\label{thm:V.4.4}
映射 $\phi \longmapsto \chi_{\phi}$ 把有界球函数的集合等同于 $\sigma(L^1(K\backslash G/K))$.
\end{theorem}

\begin{example}[交换群的情形]
\label{exa:V.4.5}
若 $G$ 交换且 $K = \{1\}$, 有界球函数就是有界连续群态射 $G \to \CC^{\times}$, 即连续群态射 $G \to S^1$, 也就是 $G$ 的不可约酉表示. 所以我们得到典则双射 $\widehat{G} \xrightarrow{\sim} \sigma(L^1(G))$.
特别地, 此时 $L^1(G)$ 上的每个乘法泛函都是 $*$-同态, 即 Banach $*$-代数 $L^1(G)$ 对称. 这重现了习题 \ref{ex:II.5.4}(c) 的结果.
若 $G$ 紧, 我们将在 (定理 \ref{thm:V.7.1}) 中看到仍然真的每个球函数都定义 $L^1(K\backslash G/K)$ 的 $*$-同态, 即 $L^1(K\backslash G/K)$ 对称. 但一般情形这不是真的.
\end{example}

\begin{proof}[定理 \ref{thm:V.4.4} 的证明]
若 $\phi$ 是有界球函数, 则 $\chi_{\phi}$ 在 $\mathcal{C}_c(K\backslash G/K)$ 上乘法, 从而因 $\mathcal{C}_c(K\backslash G/K)$ 在 $L^1(K\backslash G/K)$ 中稠密也在 $L^1(K\backslash G/K)$ 上乘法.
反之, 设 $\chi \colon L^1(K\backslash G/K) \to \CC$ 是乘法泛函. 由推论 \ref{cor:II.2.6}, 线性泛函 $\chi$ 连续且范数 $\leq 1$.
由注记 \ref{rem:V.2.2}, 线性算子 $\mathcal{C}_c(G) \to \mathcal{C}_c(K\backslash G/K)$, $f \longmapsto {}^Kf^K$ 减小 $L^1$ 范数, 所以它延拓成连续线性算子 $L^1(G) \to L^1(K\backslash G/K)$, 仍记为 $f \longmapsto {}^Kf^K$. 则 $f \longmapsto \chi({}^Kf^K)$ 是 $L^1(G)$ 上的连续线性泛函, 其范数等于 $\chi$ 的范数, 所以存在唯一的 $\phi \in L^{\infty}(G)$ 使得 $\|\phi\|_{\infty} = \|\chi\|_{\mathrm{op}}$ 且对每个 $f \in L^1(G)$,
\[
\int_G f(x)\phi(x^{-1})\,\dd x = \chi({}^Kf^K).
\]
特别地, 对所有 $k, k' \in K$ 和每个 $f \in L^1(G)$,
\begin{align*}
\int_G f(x)\phi(kx^{-1}k')\,\dd x &= \int_G f((k')^{-1}xk^{-1})\phi(x^{-1})\,\dd x \\
&= \chi({}^K(L_{k'}R_{k^{-1}}f)^K) \\
&= \chi({}^Kf^K) \\
&= \int_G f(x)\phi(x^{-1})\,\dd x.
\end{align*}
所以 $\phi$ 双不变.
设 $f, g \in L^1(K\backslash G/K)$. 有
\begin{align*}
\chi(f * g) &= \int_G (f * g)(x)\phi(x^{-1})\,\dd x \\
&= \int_{G \times G}f(y)g(y^{-1}x)\phi(x^{-1})\,\dd x\,\dd y \\
&= \int_{G \times G}f(y)\phi(y^{-1}z)g(z^{-1})\,\dd y\,\dd z \\
&= \int_G (f * \phi)(z)g(z^{-1})\,\dd z.
\end{align*}
由于 $\chi(f * g) = \chi(f)\chi(g) = \chi(f)\int_G \phi(z)g(z^{-1})\,\dd z$, 这蕴含
\[
\int_G ((f * \phi) - \chi(f)\phi)(z)g(z^{-1})\,\dd z = 0.
\]
所以对每个 $f \in L^1(K\backslash G/K)$, 有 $f * \phi = \chi(f)\phi$. 选取满足 $\chi(f) \neq 0$ 的 $f \in \mathcal{C}_c(K\backslash G/K)$. 则 $\chi(f) = f * \phi(1) = \chi(f)\phi(1)$, 所以 $\phi(1) = 1$. 又, 函数 $f * \phi$ 连续, 因为它左一致连续 (注意对每个 $x \in G$, $\|L_x(f * \phi) - f * \phi\|_{\infty} = \|(L_xf - f) * \phi\|_{\infty} \leq \|L_xf - f\|_1\|\phi\|_{\infty}$, 并应用命题 \ref{prop:I.3.1.13}). 所以 $\phi$ 局部几乎处处等于某个双不变连续有界函数, 且这个连续有界函数由命题 \ref{prop:V.4.3} 是球函数.
最后, 设 $\phi'$ 是另一个有界球函数, 使得对每个 $f \in L^1(K\backslash G/K)$,
\[
\int_G f(x)\phi'(x^{-1})\,\dd x = \int_G f(x)\phi(x^{-1})\,\dd x.
\]
我们上面已经看到, 对每个 $f \in L^1(G)$,
\[
\int_G f(x)\phi(x^{-1})\,\dd x = \chi({}^Kf^K) = \int_G {}^Kf^K(x)\phi(x^{-1})\,\dd x,
\]
对 $\phi'$ 有类似的等式. 所以
\[
\int_G f(x)\phi'(x^{-1})\,\dd x = \int_G f(x)\phi(x^{-1})\,\dd x
\]
对每个 $f \in L^1(G)$ 成立, 这蕴含 $\phi' = \phi$.
\end{proof}

\section{正定型球函数}
对第一个结果, 我们不需要假设 $(G, K)$ 是 Gelfand 对.

\begin{proposition}[循环向量与双不变性]
\label{prop:V.5.1}
设 $\phi$ 是 $G$ 上的正定型函数, $(\pi_{\phi}, V_{\phi})$ 和 $v_{\phi} \in V_{\phi}$ 是与 $\phi$ 关联的酉表示和循环向量. (见第 III.2 节.)
则 $v_{\phi} \in V_{\phi}^K$ 当且仅当 $\phi$ 双不变.
\end{proposition}

\begin{proof}
对所有 $k, k' \in K$ 和 $x \in G$,
\[
\phi(kxk') = \langle \pi_{\phi}(kxk')(v_{\phi}), v_{\phi} \rangle = \langle \pi_{\phi}(x)(\pi_{\phi}(k')(v_{\phi})), \pi_{\phi}(k^{-1})(v_{\phi}) \rangle.
\]
所以, 若 $v_{\phi} \in V_{\phi}^K$, 得到 $\phi(kxk') = \phi(x)$. 反之, 假设 $\phi$ 双不变. 在上面的等式中取 $k' = 1$, 看到对每个 $k \in K$ 和每个 $x \in G$,
\[
\phi(x) = \langle \pi_{\phi}(x)(v_{\phi}), v_{\phi} \rangle = \phi(k^{-1}x) = \langle \pi_{\phi}(x)(v_{\phi}), \pi_{\phi}(k)(v_{\phi}) \rangle.
\]
由于 $v_{\phi}$ 是循环向量, $\{\pi_{\phi}(x)(v_{\phi}) \mid x \in G\}$ 的张成在 $V_{\phi}$ 中稠密, 所以这蕴含对每个 $k \in K$ 有 $\pi_{\phi}(k)(v_{\phi}) = v_{\phi}$.
\end{proof}

\begin{theorem}[球函数与\index{端点}端点]
\label{thm:V.5.2}
再次假设 $(G, K)$ 是 Gelfand 对. 设 $\phi$ 是 $G$ 上的连续双不变函数.
若 $\phi$ 是规范化正定型函数 (即 $\phi \in P_1$), 则 $\phi$ 球当且仅当 $\phi \in E(P_1)$, 即当且仅当表示 $(\pi_{\phi}, V_{\phi})$ 不可约.
\end{theorem}

\begin{proof}
记 $(\pi, V)$ 和 $v$ 为 $(\pi_{\phi}, V_{\phi})$ 和 $v_{\phi}$. 由于 $\phi$ 双不变, 由命题 \ref{prop:V.5.1} 知道 $v \in V^K$. 首先假设 $\phi \in E(P_1)$, 即 $\pi$ 不可约. 由定理 \ref{thm:V.3.1.1}, 有 $\dim(V^K) = 1$. 设 $f \in \mathcal{C}_c(K\backslash G/K)$. 则对每个 $x \in G$,
\begin{align*}
f * \phi(x) &= \int_G f(y)\langle \pi(y^{-1}x)(v), v \rangle\,\dd y \\
&= \int_G f(y)\langle \pi(x)(v), \pi(y)v \rangle\,\dd y \\
&= \langle \pi(x)(v), \pi(f)(v) \rangle.
\end{align*}
由于 $\pi(f)(v) \in V^K$, 可以把 $\pi(f)(v)$ 写成 $\pi(f)(v) = \chi(f)v$, 其中 $\chi(f) \in \CC$, 得到对每个 $x \in G$,
\[
f * \phi(x) = \chi(f)\langle \pi(x)(v), v \rangle = \chi(f)\phi(x).
\]
由命题 \ref{prop:V.4.3}, 这蕴含 $\phi$ 球.
反之, 假设 $\phi$ 球. 则再次由命题 \ref{prop:V.4.3}, 存在映射 $\chi \colon \mathcal{C}_c(K\backslash G/K) \to \CC$ 使得对每个 $f \in \mathcal{C}_c(K\backslash G/K)$ 有 $f * \phi = \chi(f)\phi$. 换句话说, 对每个 $f \in \mathcal{C}_c(K\backslash G/K)$ 和每个 $x \in G$,
\begin{align*}
f * \phi(x) &= \int_G f(y)\langle \pi(y^{-1}x)(v), v \rangle\,\dd y = \langle \pi(x)(v), \pi(f)(v) \rangle \\
&= \chi(f)\phi(x) \\
&= \chi(f)\langle \pi(x)(v), v \rangle.
\end{align*}
由于 $v$ 是循环向量, 这蕴含 $\pi(f)(v) = \chi(f)v \in V^K$. 但我们在引理 \ref{lem:V.3.1.4} 的证明中看到空间 $\{\pi(f)(v) \mid f \in \mathcal{C}_c(K\backslash G/K)\}$ 在 $V^K$ 中稠密 (若 $v$ 循环), 所以 $\dim(V^K) = 1$. 由引理 \ref{lem:V.5.3}, 这蕴含 $(\pi, V)$ 不可约.
\end{proof}

\begin{lemma}[$K$-固定循环向量与不可约性]
\label{lem:V.5.3}
我们不假设 $(G, K)$ 是 Gelfand 对. 设 $(\pi, V)$ 是 $G$ 的酉表示, 假设 $V^K$ 中存在循环向量. 若 $\dim(V^K) \leq 1$, 则 $(\pi, V)$ 不可约.
\end{lemma}

\begin{proof}
只需证明 $\End_G(V) = \CC\mathrm{id}_V$. 事实上, 若 $V$ 有闭 $G$-不变子空间 $W$ 满足 $W \neq \{0\}, V$, 则到 $W$ 上的正交投影是 $V$ 的 $G$-等变自同态 (由引理 \ref{lem:I.3.4.3}), 它不是 $\mathrm{id}_V$ 的倍数.
所以设 $T \in \End_G(V)$. 则由命题 \ref{prop:V.1.7}, 算子 $T$ 与到 $V^K$ 上的正交投影交换, 所以它保持 $V^K$. 选取循环向量 $v \in V^K$. 由于 $\dim(V^K) = 1$, 有 $T(v) = \lambda v$, 其中 $\lambda \in \CC$. 由于 $T$ 是 $G$-等变的, 得到对每个 $x \in G$ 有 $T(\pi(x)(v)) = \lambda\pi(x)(v)$. 由于 $v$ 循环, 这蕴含 $T = \lambda\mathrm{id}_V$.
\end{proof}

\begin{corollary}[球函数与不可约表示的对应]
\label{cor:V.5.4}
假设 $(G, K)$ 是 Gelfand 对. 则 $\phi \longmapsto (\pi_{\phi}, V_{\phi})$ 诱导从 $E(P_1)$ 中球函数的集合到满足 $V^K \neq \{0\}$ 的 $G$ 的不可约酉表示 $(\pi, V)$ 的酉等价类的集合的双射.
\end{corollary}

\begin{proof}
唯一不立即由命题 \ref{prop:V.5.1} 和定理 \ref{thm:V.5.2} 得出的结论是: 若 $E(P_1)$ 中两个球函数给出酉等价的表示, 则它们必相等. 设 $\phi_1, \phi_2 \in E(P_1)$ 球, 假设存在等距 $G$-等变同构 $T \colon V_{\phi_1} \to V_{\phi_2}$. 由命题 \ref{prop:V.5.1}, 向量 $v_{\phi_1}$ 和 $v_{\phi_2}$ 是 $K$-不变的. 又, 由于 $(G, K)$ 是 Gelfand 对, 空间 $V_{\phi_1}^K$ 和 $V_{\phi_2}^K$ 的维数都 $\leq 1$, 从而因它们包含非零向量而为 $1$. 但 $T$ 是 $G$-等变的, 所以 $T(V_{\phi_1}^K) \subset V_{\phi_2}^K$, 这蕴含对某个 $\lambda \in \CC$ 有 $T(v_{\phi_1}) = \lambda v_{\phi_2}$. 由于 $\|v_{\phi_1}\| = \|v_{\phi_2}\| = 1$, 必有 $|\lambda| = 1$. 所以对每个 $x \in G$,
\begin{align*}
\phi_2(x) &= \langle \pi_{\phi_2}(x)(v_{\phi_2}), v_{\phi_2} \rangle \\
&= \langle \pi_{\phi_2}(x)(\lambda^{-1}T(v_{\phi_1})), \lambda^{-1}T(v_{\phi_1}) \rangle \\
&= \langle T(\pi_{\phi_1}(x)(v_{\phi_1})), T(v_{\phi_1}) \rangle \\
&= \langle \pi_{\phi_1}(x)(v_{\phi_1}), v_{\phi_1} \rangle.
\end{align*}
\end{proof}

\section{\index{对偶空间}对偶空间与\index{球傅里叶变换}球傅里叶变换}
在本节中, 我们假设 $(G, K)$ 是 Gelfand 对. 我们将不加证明地陈述关于 (球)傅里叶变换的一些结果. 在下一节中, 若 $G$ 紧, 我们将给出这些结果的某种版本的证明.

\begin{definition}[对偶空间 (dual space)]
\label{def:V.6.1}
$(G, K)$ 的对偶空间是 $E(P_1)$ 中球函数的集合 $\mathcal{Z}$, 赋予由嵌入 $E(P_1) \subset L^{\infty}(G) \simeq \Hom(L^1(G), \CC)$ 给出的弱*拓扑.
\end{definition}

\begin{example}[交换群的情形]
\label{exa:V.6.2}
若 $G$ 交换且 $K = \{1\}$, 则 $\mathcal{Z} = \widehat{G}$, 即 $G$ 的对偶群. (见习题 \ref{ex:I.5.4.1}.)
\end{example}

\begin{proposition}[$\mathcal{Z}$ 局部紧]
\label{prop:V.6.3}
空间 $\mathcal{Z}$ 局部紧, 且其拓扑与 $G$ 的紧子集上的收敛拓扑一致.
\end{proposition}

\begin{proof}
对第一个结论, 首先注意 $P_0 = \{\psi \text{ 正定型} \mid \psi(1) \leq 1\}$ 弱*紧, 因为它在 $L^{\infty}(G)$ 的闭单位球中弱*闭. 由定理 \ref{thm:V.4.4} 的证明, 子集 $P_0 \cap \mathcal{C}(K\backslash G/K)$ 是满足如下条件的 $\phi \in P_0$ 的集合: 对每个 $f \in L^1(G)$, $\int_G f(x)\phi(x^{-1})\,\dd x = \int_G {}^Kf^K(x)\phi(x^{-1})\,\dd x$. 这些都是弱*闭条件, 所以 $P_0 \cap \mathcal{C}(K\backslash G/K)$ 在 $P_0$ 中弱*闭, 从而弱*紧. 最后, 由定理 \ref{thm:V.5.2}, 集合 $\mathcal{Z} \cup \{0\}$ 是 $P_0 \cap \mathcal{C}(K\backslash G/K)$ 中满足如下条件的 $\phi$ 的集合: 对所有 $f, g \in L^1(K\backslash G/K)$ 有 $\int_G (f * g)(x)\phi(x^{-1})\,\dd x = \bigl(\int_G f(x)\phi(x^{-1})\,\dd x\bigr)\bigl(\int_G g(x)\phi(x^{-1})\,\dd x\bigr)$. 这是弱*闭条件, 所以 $\mathcal{Z} \cup \{0\}$ 弱*紧, 且 $\mathcal{Z}$ 局部紧. 注意这也证明 $\mathcal{Z} \cup \{0\}$ 是 $\mathcal{Z}$ 的 Alexandroff 紧化.
第二个结论由 Raikov 定理 (定理 \ref{thm:III.4.3}) 立即得出.
\end{proof}

\begin{definition}[球傅里叶变换 (spherical Fourier transform)]
\label{def:V.6.4}
设 $f \in L^1(K\backslash G/K)$. $f$ 的 (球)傅里叶变换是函数 $\hat{f} \colon \mathcal{Z} \to \CC$:
\[
\hat{f}(\phi) = \int_G f(x)\phi(x^{-1})\,\dd x = \chi_{\phi}(f).
\]
\end{definition}

\begin{proposition}[傅里叶变换的性质]
\label{prop:V.6.5}
傅里叶变换有下列性质:
\begin{enumerate}
\item[(i)] 对每个 $f \in L^1(K\backslash G/K)$, 函数 $\hat{f} \in \mathcal{C}_0(\mathcal{Z})$, 且 $\|\hat{f}\|_{\infty} \leq \|f\|_1$.
\item[(ii)] 映射 $L^1(K\backslash G/K) \to \mathcal{C}_0(\mathcal{Z})$, $f \longmapsto \hat{f}$ 是 $\CC$-线性的且有稠密像.
\item[(iii)] 对所有 $f, g \in L^1(K\backslash G/K)$, $\widehat{f * g} = \hat{f}\hat{g}$.
\item[(iv)] 对每个 $f \in L^1(K\backslash G/K)$, $\widehat{f^*} = \overline{\hat{f}}$.
\end{enumerate}
\end{proposition}

\begin{proof}
(i). $\hat{f}$ 的连续性由弱*拓扑的定义立即得出. 事实上, 可以用同样的公式把 $\hat{f}$ 延拓成整个 $L^{\infty}(G)$ 上的连续线性泛函. 但我们在命题 \ref{prop:V.6.3} 的证明中看到 $\mathcal{Z} \cup \{0\}$ 是 $\mathcal{Z}$ 的 Alexandroff 紧化, 所以这蕴含 $\hat{f} \in \mathcal{C}(\mathcal{Z})$ 在无穷远处消失. 不等式 $\|\hat{f}\|_{\infty} \leq \|f\|_1$ 由对每个 $\phi \in \mathcal{Z}$ 有 $\|\phi\|_{\infty} = 1$ 这一事实得出.
(iii) 和 (iv). 这只是在表达对每个 $\phi \in \mathcal{Z}$, $\chi_{\phi}$ 是从 $L^1(K\backslash G/K)$ 到 $\CC$ 的 $*$-同态这一事实.
(ii). 线性显然. 第二个结论由 Stone-Weierstrass 定理得出: 事实上, 球傅里叶变换的像是由 (iii) 给出的 $\mathcal{C}_0(\mathcal{Z})$ 的 $\CC$-子代数, 它由 (iv) 被复共轭稳定, 它分离点 (因为由定理 \ref{thm:V.4.4}, 映射 $\mathcal{Z} \to \sigma(L^1(K\backslash G/K))$, $\phi \longmapsto (f \longmapsto \hat{f}(\phi))$ 单射), 且它不在任何地方消失 (对每个 $\phi \in \mathcal{Z}$, 映射 $f \longmapsto \hat{f}(\phi)$ 是 $L^1(K\backslash G/K)$ 上的乘法泛函, 所以它非零).
\end{proof}

\begin{theorem}[Fourier 反演]\footnote{见 {[}25{]} 定理 6.4.5.}
\label{thm:V.6.6}
设 $V^1(K\backslash G/K)$ 是 $G$ 上正定型双不变函数的复线性组合的 $L^1$ 函数的空间.
则存在 $\mathcal{Z}$ 上唯一的测度 $\nu$, 称为 \index{Plancherel 测度}Plancherel 测度, 使得对每个 $f \in V^1(K\backslash G/K)$, $\hat{f} \in L^1(\mathcal{Z}, \nu)$, 且对每个 $x \in G$,
\[
f(x) = \int_{\mathcal{Z}}\phi(x)\hat{f}(\phi)\,\dd\nu.
\]
\end{theorem}

\begin{theorem}[Plancherel 公式]\footnote{见 {[}25{]} 定理 6.4.6.}
\label{thm:V.6.7}
对每个 $f \in \mathcal{C}_c(K\backslash G/K)$, $\hat{f} \in L^2(\mathcal{Z}, \nu)$, 且
\[
\int_G |f(x)|^2\,\dd x = \int_{\mathcal{Z}}|\hat{f}(\phi)|^2\,\dd\nu(\phi).
\]
特别地, 映射 $f \longmapsto \hat{f}$ 延拓成等距 $L^2(K\backslash G/K) \to L^2(\mathcal{Z}, \nu)$, 且它是同构.
\end{theorem}

\begin{remark}[Pontrjagin 对偶]
\label{rem:V.6.8}
若 $G$ 交换且 $K = \{1\}$, 则 $\mathcal{Z} = \widehat{G}$ 是局部紧群, 测度 $\nu$ 是 $\widehat{G}$ 上的 Haar 测度, 且 Pontrjagin 对偶定理说典则映射 $G \to \widehat{\widehat{G}}$, $x \longmapsto (\phi \longmapsto \phi(x))$ 是拓扑群的同构. (见例如 {[}11{]} 定理 4.22 和 4.32, 或 {[}25{]} 定理 5.5.1 和 5.7.1.)
但一般地, Plancherel 测度 $\nu$ 可以支撑在 $\mathcal{Z}$ 的真子集上.
\end{remark}

\section{紧群的情形}
在本节中, 我们假设 $(G, K)$ 是 Gelfand 对, 且 $G$ 紧. 还假设 $G$ 上的 Haar 测度规范化.

\begin{theorem}[紧 Gelfand 对]
\label{thm:V.7.1}
\begin{enumerate}
\item[(i)] $(G, K)$ 的对偶空间 $\mathcal{Z}$ 离散, 且是 $L^2(G)$ 的正交子集.
\item[(ii)] $G$ 上的每个球函数正定型 (从而由定理 \ref{thm:V.5.2} 在 $E(P_1)$ 中). 换句话说, 集合 $\mathcal{Z}$ 通过 $\phi \longmapsto (\pi_{\phi}, V_{\phi})$ 与满足 $\dim(V_{\phi}^K) = 1$ 的 $G$ 的不可约酉表示的等价类的集合典则双射对应.
\item[(iii)] 对每个 $\phi \in \mathcal{Z}$,
\[
\phi(x) = \int_K\chi_{\pi_{\phi}}(xk)\,\dd k
\]
对 $x \in G$, 且
\[
\int_G|\phi(x)|^2 = \frac{1}{\dim V_{\phi}}.
\]
\item[(iv)] 若 $f \in L^2(K\backslash G/K)$, $(\pi, V) \in \widehat{G}$, 则 $V^K = \{0\}$ 时 $f * \chi_{\pi} = 0$, 否则 $f * \chi_{\pi}$ 是由推论 \ref{cor:V.5.4} 对应 $\pi$ 的 $\mathcal{Z}$ 的元素 $\phi_{\pi}$ 的倍数.
\end{enumerate}
\end{theorem}

\begin{proof}
设 $\phi, \phi' \in \mathcal{Z}$ 满足 $\phi \neq \phi'$. 由推论 \ref{cor:V.5.4} (和命题 \ref{prop:IV.2.6}) 知道表示 $V_{\phi}$ 和 $V_{\phi'}$ 酉且不等价. 又知道 (由表示的构造) $\phi$ 和 $\phi'$ 分别是 $V_{\phi}$ 和 $V_{\phi'}$ 的矩阵系数. 由 Schur 正交性 (定理 \ref{thm:IV.3.8}), 这蕴含 $\langle \phi, \phi' \rangle_{L^2(G)} = 0$.
我们证明 $\mathcal{Z}$ 离散. 设 $\phi \in \mathcal{Z}$, 考虑 $U = \{\phi' \in \mathcal{Z} \mid \|\phi - \phi'\|_{\infty} < \|\phi\|_2\}$. 这在 $G$ 的紧子集上的收敛拓扑中开 (因为 $G$ 紧), 从而由 Raikov 定理 (定理 \ref{thm:III.4.3}) 是 $\mathcal{Z}$ 的开子集. 又, 若 $\phi' \in U$, 则
\begin{align*}
|\langle \phi, \phi' \rangle_{L^2(G)}| &= |\langle \phi, \phi \rangle_{L^2(G)} - \langle \phi, \phi' - \phi \rangle_{L^2(G)}| \\
&\geq \|\phi\|_2^2 - \|\phi\|_2\|\phi - \phi'\|_2 \\
&\geq \|\phi\|_2^2 - \|\phi\|_2\|\phi - \phi'\|_{\infty} \\
&> 0,
\end{align*}
所以由第一段, $\phi' = \phi$. 这意味着 $U = \{\phi\}$, 即 $\phi$ 是 $\mathcal{Z}$ 的孤立点.
设 $(\pi, V)$ 是 $G$ 的不可约酉表示, $f \in L^2(K\backslash G/K)$. 要计算 $f * \chi_{\pi}$. 设 $(v_1, \ldots, v_d)$ 是标准正交基; 则对每个 $x \in G$ 有 $\chi_{\pi}(x) = \sum_{i=1}^d\langle \pi(x)(e_i), e_i \rangle$. 所以对每个 $x \in G$,
\[
f * \chi_{\pi}(x) = \int_G f(y)\sum_{i=1}^d\langle \pi(y^{-1}x)(e_i), e_i \rangle = \sum_{i=1}^d\langle \pi(x)(e_i), \pi(f)(e_i) \rangle.
\]
设 $P_K \in \End(V)$ 是 $V^K$ 上的正交投影. 由于 $f$ 双不变, 由命题 \ref{prop:V.1.7} 有 $\pi(f) = P_K \circ \pi(f) \circ P_K$. 首先假设 $V^K = \{0\}$. 则上面的公式给出 $f * \chi_{\pi} = 0$. 现在假设 $V^K \neq \{0\}$. 则由推论 \ref{cor:V.5.4}, 存在唯一正定型球函数 $\phi_{\pi}$, 其关联表示是 $(\pi, V)$, 且有酉循环向量 $v \in V^K$ 使得 $\phi_{\pi}(x) = \langle \pi(x)(v), v \rangle$. 可以选择标准正交基使得 $v_1 = v$. 则对 $i \geq 2$ 有 $P_K(v_i) = 0$, $P_K(v_1) = v_1$, 且对每个 $x \in G$,
\[
f * \chi_{\pi}(x) = \sum_{i=1}^d\langle \pi(x)(v_i), P_K(\pi(f)(P_K(v_i))) \rangle = \langle \pi(x)(v_1), P_K(\pi(f)(v_1)) \rangle.
\]
由于 $V^K$ 是 $1$ 维的, 向量 $P_K(\pi(f)(v_1))$ 是 $v_1$ 的倍数, 所以 $f * \chi_{\pi}$ 是 $\phi_{\pi}$ 的倍数. 这证明 (iv). 又注意对每个 $x \in G$,
\begin{align*}
\int_K\chi_{\pi}(kx)\,\dd k &= \int_K\sum_{i=1}^d\langle \pi(kx)(v_i), v_i \rangle\,\dd k \\
&= \sum_{i=1}^d\Bigl\langle \pi(x)(v_i), \int_K\pi(k^{-1})(v_i)\,\dd k\Bigr\rangle \\
&= \sum_{i=1}^d\langle \pi(x)(v_i), P_K(v_i) \rangle \quad \text{(由命题 \ref{prop:V.1.7})} \\
&= \langle \pi(x)(v_1), v_1 \rangle \\
&= \phi_{\pi}(x),
\end{align*}
这给出 (iii) 的第一部分. (iii) 的第二部分包含在命题 \ref{thm:IV.3.8}(ii) 中.
现在考虑 $G$ 上的球函数 $\phi$. 由命题 \ref{prop:IV.7.1} (即 Fourier 反演公式), 在 $L^2(G)$ 中有等式
\[
\phi = \sum_{\pi \in \widehat{G}}\dim(\pi)\phi * \chi_{\pi}.
\]
由上面的计算, 只有具有非零 $K$-不变向量的 $\pi \in \widehat{G}$ 出现在上面的和中, 且此时 $\phi * \chi_{\pi}$ 是上一段记为 $\phi_{\pi}$ 的函数的倍数. 换句话说, 再次用推论 \ref{cor:V.5.4}, 得到
\[
\phi = \sum_{\psi \in \mathcal{Z}}c_{\psi}\psi,
\]
其中 $c_{\psi} \in \CC$. 若记 $\chi_{\phi}$ (相应地 $\chi_{\psi}$) 为 $L^1(K\backslash G/K)$ 上的线性泛函 $f \longmapsto f * \phi(1)$ (相应地 $f \longmapsto f * \psi(1)$), 知道它是乘法的 (对 $\chi_{\psi}$, 这用定理 \ref{thm:V.5.2}). 又, 由于 $\phi = \sum_{\psi \in \mathcal{Z}}c_{\psi}\psi$, 有 $\chi_{\phi} = \sum_{\psi \in \mathcal{Z}}c_{\psi}\chi_{\psi}$. 设 $\psi, \psi' \in \mathcal{Z}$ 满足 $\psi \neq \psi'$. 则 $\psi * \psi' = \psi' * \psi$ 既是 $\psi$ 又是 $\psi'$ 的倍数 (由命题 \ref{prop:V.4.3}), 所以 $\psi * \psi' = 0$. 特别地, $\chi_{\psi}(\psi') = \chi_{\psi'}(\psi) = 0$. 这蕴含对每个 $\psi \in \mathcal{Z}$ 有 $\chi_{\phi}(\psi) = c_{\psi}\chi_{\psi}(\psi)$; 又注意
\[
\chi_{\psi}(\psi) = \int_G \psi(x)\psi(x^{-1})\,\dd x = \int_G \psi(x)\overline{\psi(x)}\,\dd x > 0.
\]
所以, 若 $\psi, \psi' \in \mathcal{Z}$ 且 $\psi \neq \psi'$, 则
\[
0 = \chi_{\phi}(\psi * \psi') = \chi_{\phi}(\psi)\chi_{\phi}(\psi') = c_{\psi}c_{\psi'}\chi_{\psi}(\psi)\chi_{\psi'}(\psi'),
\]
所以 $c_{\psi}c_{\psi'} = 0$. 所以 $c_{\psi}$ 中至多一个非零, 即存在 $\psi \in \mathcal{Z}$ 使得 $\phi = c_{\psi}\psi$. 由于 $\phi(1) = 1 = \psi(1)$, 必有 $c_{\psi} = 1$, 所以最终看到 $\phi = \psi$ 正定型. 这完成 (ii) 的证明.
\end{proof}

\begin{corollary}[紧 Gelfand 对的调和分析]
\label{cor:V.7.2}
\begin{enumerate}
\item[(i)] 有 $G$-等变同构
\[
L^2(G/K) \simeq \bigoplus_{\phi \in \mathcal{Z}}V_{\phi}.
\]
\item[(ii)] 族 $((\dim V_{\phi})^{1/2}\phi)_{\phi \in \mathcal{Z}}$ 是 $L^2(K\backslash G/K)$ 的 Hilbert 基.
\item[(iii)] 对每个 $f \in L^2(K\backslash G/K)$,
\[
f = \sum_{\phi \in \mathcal{Z}}\dim(V_{\phi})\hat{f}(\phi)\phi
\]
(Fourier 反演公式) 且
\[
\|f\|_{L^2(G)}^2 = \sum_{\phi \in \mathcal{Z}}\dim(V_{\phi})|\hat{f}(\phi)|^2
\]
(Parseval 公式).
\end{enumerate}
\end{corollary}

\begin{proof}
(i) 只是定理 \ref{thm:V.3.2.4} 最后一个结论的重述.
对 (ii), 我们已经知道族 $(\sqrt{\dim(V_{\phi})}\phi)_{\phi \in \mathcal{Z}}$ 在 $L^2(G)$ 中标准正交. 又, 若 $f \in L^2(K\backslash G/K)$, 由命题 \ref{prop:IV.7.1} 有
\[
f = \sum_{(\pi, V) \in \widehat{G}}\dim(V)f * \chi_{\pi},
\]
所以 $f$ 在 $\operatorname{Span}(\mathcal{Z})$ 的闭包中 (由定理的 (iv)), 这意味着 $\operatorname{Span}(\mathcal{Z})$ 在 $L^2(K\backslash G/K)$ 中稠密.
(iii) 的第二个公式由第一个公式和 (ii) 得出. 为证明第一个公式, 只需证明对每个 $f \in L^2(K\backslash G/K)$ 和每个 $\phi \in \mathcal{Z}$ 有 $f * \chi_{\pi_{\phi}} = \hat{f}(\phi)\phi$. 由于我们已经知道 $f * \chi_{\pi_{\phi}}$ 是 $\phi$ 的倍数, 只需验证 $f * \chi_{\pi_{\phi}}(1) = \hat{f}(\phi)$. 由定理的 (iii), 对每个 $x \in G$ 有 $\phi(x) = \int_K\chi_{\pi_{\phi}}(kx)\,\dd k$. 所以:
\begin{align*}
f * \chi_{\pi_{\phi}}(1) &= \int_G f(x)\chi_{\pi_{\phi}}(x^{-1})\,\dd x \\
&= \int_{G \times K}f(k^{-1}x)\chi_{\pi_{\phi}}(x)\,\dd x\,\dd k \quad \text{($f$ 左不变且 } \operatorname{vol}(K) = 1\text{)} \\
&= \int_{G \times K}f(x)\chi_{\pi_{\phi}}(kx)\,\dd x\,\dd k \\
&= \int_G f(x)\phi(x)\,\dd x \\
&= \int_G f(x)\overline{\phi(x^{-1})}\,\dd x \\
&= \hat{f}(\phi).
\end{align*}
\end{proof}

\begin{remark}[Plancherel 测度]
\label{rem:V.7.3}
推论特别地说 $\mathcal{Z}$ 上的 Plancherel 测度 $\nu$ 由 $|\nu(\{\phi\})| = \dim(V_{\phi})$ 给出.
\end{remark}

\section{习题}

\subsection{\texorpdfstring{Gelfand 对 $(\SO(n), \SO(n-1))$}{Gelfand 对}}

这一系列习题的材料是经典的, 但这里的叙述深受 {[}22{]} 的第 2.3.2 节和 {[}25{]} 的第 7.3 节的影响.

\begin{exercise}[调和多项式]
\label{ex:V.8.1.1}
固定正整数 $n$. 对每个 $m \in \Z_{\geq 0}$, 记 $V_m(\R^n)$ 为 $\R^n$ 上 $m$ 次齐次复值多项式函数的向量空间. 定义 $\Or(n)$ 在 $V_m(\R^n)$ 上的作用为 $(x \cdot f)(v) = f(x^{-1}v)$ ($x \in \Or(n)$, $f \in V_m(\R^n)$, $v \in \R^n$; 换句话说 $x \cdot f = L_xf$).
对 $i \in \{1, \ldots, n\}$, 记 $\partial_{x_i}$ 为 $\mathcal{C}^{\infty}(\R^n)$ (从 $\R^n$ 到 $\CC$ 的光滑函数空间) 的自同态 $f \longmapsto \frac{\partial f}{\partial x_i}$, 令 $\Delta = \sum_{i=1}^n(\partial_{x_i})^2$ (这称为 Laplace 算子 (Laplacian operator)).
$\R^n$ 上 $m$ 次调和多项式 (harmonic polynomial) 的空间是
\[
\mathcal{H}_m(\R^n) = \{f \in V_m(\R^n) \mid \Delta(f) = 0\}.
\]
\begin{enumerate}
\item[(a)] 计算 $\dim(V_m(\R^n))$.
\item[(b)] 证明 $\Or(n)$ 在 $V_m(\R^n)$ 上的作用是连续表示.
\item[(c)] 证明对每个 $x \in \Or(n)$ 和每个 $f \in \mathcal{C}^{\infty}(\R^n)$, $\Delta(L_xf) = L_x(\Delta(f))$. (用习题 \ref{ex:V.8.1.2}(c) 可以帮你解决此题.)
\item[(d)] 证明 $V_m(\R^n)$ 的子空间 $\mathcal{H}_m(\R^n)$ 是 $\Or(n)$-不变的.
\end{enumerate}
\end{exercise}

\begin{solution}
(a). 对每个 $i \in \{1, \ldots, n\}$, 记 $x_i \in V_1(\R^n)$ 为函数 $(z_1, \ldots, z_n) \longmapsto z_i$. 则 $\{x_1^{i_1}\cdots x_n^{i_n} \mid i_1, \ldots, i_n \in \Z_{\geq 0}, i_1 + \cdots + i_n = m\}$ 是 $V_m(\R^n)$ 的基. 所以
\[
\dim(V_m(\R^n)) = |\{(i_1, \ldots, i_n) \in (\Z_{\geq 0})^n \mid i_1 + \cdots + i_n = m\}|.
\]
这也等于
\[
|\{(j_1, \ldots, j_n) \in (\Z_{\geq 1})^n \mid j_1 + \cdots + j_n = m + n\}|
\]
(取 $j_r = i_r + 1$). 在上面的集合中选择 $(j_1, \ldots, j_n)$ 等价于选择数 $j_1, j_1 + j_2, \ldots, j_1 + \cdots + j_{n-1}$, 它们构成 $\{1, \ldots, n + m - 1\}$ 的基数为 $n - 1$ 的子集. 所以得到
\[
\dim(V_m(\R^n)) = \binom{n + m - 1}{n - 1} = \binom{n + m - 1}{m}.
\]
(b). 若用 (a) 的 $V_m(\R^n)$ 的基, $x \in \Or(n)$ 的作用由系数是 $x$ 的元素的多元多项式的矩阵给出. 所以对每个 $f \in V_m(\R^n)$, 映射 $\Or(n) \to V_m(\R^n)$, $x \longmapsto x \cdot f$ 连续. 由于 $V_m(\R^n)$ 有限维, 这蕴含作用连续.
(c). 为避免计算, 用习题 \ref{ex:V.8.1.2}(c). 注意 $\Delta = \partial_{x_1^2 + \cdots + x_n^2}$. 所以由习题 \ref{ex:V.8.1.2}(c), 对每个 $f \in \mathcal{C}^{\infty}(\R^n)$ 和每个 $x \in G$,
\[
\Delta(x \cdot f) = x \cdot (\partial_g f),
\]
其中 $g = L_{x^T}(x_1^2 + \cdots + x_n^2)$. 所以我们只需证明 $x_1^2 + \cdots + x_n^2 \in V_2(\R^n)$ 被 $\Or(n)$ 的所有元素不变, 这直接由 $\Or(n)$ 的定义得出.
(d). 问题 (c) 蕴含 $\Delta \colon V_m(\R^n) \to V_{m-2}(\R^n)$ 是 $\Or(n)$-等变的, 且 $\mathcal{H}_m(\R^n)$ 是它的核.
\end{solution}

\begin{exercise}[调和多项式的内积]
\label{ex:V.8.1.2}
我们保留习题 \ref{ex:V.8.1.1} 的记号. 对 $i \in \{1, \ldots, n\}$, 记 $x_i$ 为 $\R^n$ 上的第 $i$ 个坐标函数.
\begin{enumerate}
\item[(a)] 证明映射 $x_i \longmapsto \partial_{x_i}$ 延拓成从 $\bigoplus_{m \geq 0}V_m(\R^n)$ ($\R^n$ 上复值多项式函数的代数) 到 $\End(\mathcal{C}^{\infty}(\R^n))$ 的唯一 $\CC$-代数态射. 记此态射为 $f \longmapsto \partial_f$.
\end{enumerate}
若 $f, g \in V_m(\R^n)$, 令 $\langle f, g \rangle = \partial_g(f)$. (注意 $g$ 仍是 $\R^n$ 上的多项式函数.)
\begin{enumerate}
\item[(a)] 证明 $\langle \cdot, \cdot \rangle$ 是 $V_m(\R^n)$ 上的内积. (提示: 能找到正交基吗?)
\item[(b)] 证明对每个 $f \in V_m(\R^n)$ 和每个 $y \in \Or(n)$, 在 $\End(\mathcal{C}^{\infty}(\R^n))$ 中有 $\partial_f \circ L_y = L_y \circ \partial_{L_{y^T}f}$.
\item[(c)] 证明习题 \ref{ex:V.8.1.1} 中定义的 $\Or(n)$ 在 $V_m(\R^n)$ 上的连续表示对内积 $\langle \cdot, \cdot \rangle$ 酉.
\item[(d)] 若 $m \leq 1$, 证明 $V_m(\R^n) = \mathcal{H}_m(\R^n)$.
\item[(e)] 若 $m \geq 2$, 证明 $\mathcal{H}_m(\R^n)^{\perp} = |x|^2V_{m-2}(\R^n)$, 其中 $|x|^2$ 是函数 $\sum_{i=1}^n x_i^2 \in V_2(\R^n)$.
\item[(f)] 证明
\[
V_m(\R^n) = \bigoplus_{k=0}^{\lfloor m/2 \rfloor}|x|^{2k}\mathcal{H}_{m-2k}(\R^n),
\]
且这诱导 $\Or(n)$-等变同构
\[
V_m(\R^n) = \bigoplus_{k=0}^{\lfloor m/2 \rfloor}\mathcal{H}_{m-2k}(\R^n).
\]
\item[(g)] 若 $S \subset \R^n$ 是单位球面, 证明映射 $\bigoplus_{m \geq 0}\mathcal{H}_m(\R^n) \to \mathcal{C}(S)$, $f \longmapsto f|_S$ 单射.
\item[(h)] 证明对每个 $f \in V_m(\R^n)$, 存在唯一的 $g \in \bigoplus_{k=0}^{\lfloor m/2 \rfloor}\mathcal{H}_{m-2k}(\R^n)$ 使得 $f|_S = g|_S$.
\end{enumerate}
\end{exercise}

\begin{solution}
(a). 注意 $\bigoplus_{m \geq 0}V_m(\R^n)$ 同构于多项式代数 $\CC[x_1, \ldots, x_n]$. 所以只需验证对所有 $i, j \in \{1, \ldots, n\}$, $\partial_{x_i}$ 和 $\partial_{x_j}$ 交换. 但这是 $\mathcal{C}^2$ 函数的偏导数的众所周知的性质.
(b). 首先, 由定义显然 $\langle \cdot, \cdot \rangle$ 对第一个变量线性、对第二个变量反线性. 我们在习题 \ref{ex:V.8.1.1}(a) 的基中计算此形式的矩阵.
设 $f = x_1^{i_1}\cdots x_n^{i_n}$, 其中 $i_1, \ldots, i_n \in \Z_{\geq 0}$, $i_1 + \cdots + i_n = m$. 若 $r \in \{1, \ldots, n\}$, $a \in \Z_{\geq 0}$, 有
\[
\partial_{x_r}^a f = \begin{cases} 0 & \text{若 } a > i_r, \\ i_r(i_r - 1)\cdots(i_r - a + 1)x_r^{i_r - a}\prod_{s \neq r}x_s^{i_s} & \text{若 } a \leq i_r. \end{cases}
\]
设 $g = x_1^{j_1}\cdots x_n^{j_n}$, 其中 $j_1, \ldots, j_n \in \Z_{\geq 0}$, $j_1 + \cdots + j_n = m$. 由于 $i_1 + \cdots + i_n = j_1 + \cdots + j_n$, 要么存在 $r \in \{1, \ldots, n\}$ 使得 $j_r > i_r$, 要么对所有 $r \in \{1, \ldots, n\}$ 有 $i_r = j_r$. 第一种情形 $\langle f, g \rangle = \partial_g f = 0$. 第二种情形,
\[
\langle f, g \rangle = \partial_g f = i_1!i_2!\cdots i_n!.
\]
所以 $\langle \cdot, \cdot \rangle$ 在 $V_m(\R^n)$ 的基 $\{x_1^{i_1}\cdots x_n^{i_n} \mid i_1, \ldots, i_n \in \Z_{\geq 0}, i_1 + \cdots + i_n = m\}$ 中的矩阵是对角矩阵, 对角线元素为正实数, 特别地 Hermitian 正定. 这蕴含 $\langle \cdot, \cdot \rangle$ 是内积.
(c). 此结论实际上对每个 $y \in \GL_n(\R)$ 成立, 我们将证明这一点.
首先注意命题中的恒等式对代数 $V(\R^n) := \bigoplus_{m \geq 0}V_m(\R^n)$ 中的 $f$ 有意义, 且对 $f$ 线性. 又, 若它对 $f, g \in V(\R^n)$ 成立, 则对 $y \in \GL_n(\R)$,
\begin{align*}
\partial_{fg} \circ L_y &= \partial_f \circ \partial_g \circ L_y \\
&= \partial_f \circ L_y \circ \partial_{L_{y^T}g} \\
&= L_y \circ \partial_{L_{y^T}f} \circ \partial_{L_{y^T}g} \\
&= L_y \circ \partial_{L_{y^T}fL_{y^T}g} \\
&= L_y \circ \partial_{L_{y^T}(fg)},
\end{align*}
即恒等式对 $fg$ 也成立. 总之, 只需对函数 $x_1, \ldots, x_n$ 证明.
设 $i \in \{1, \ldots, n\}$, $y \in \GL_n(\R)$, 写 $y^{-1} = a = (a_{ij}) \in \GL_n(\R)$. 则对每个 $(z_1, \ldots, z_n) \in \R^n$,
\[
(L_{y^T}x_i)(z_1, \ldots, z_n) = x_i(a^T(z_1, \ldots, z_n)) = \sum_{j=1}^n a_{ji}z_j.
\]
换句话说,
\[
L_{y^T}x_i = \sum_{j=1}^n a_{ji}x_j,
\]
所以
\[
\partial_{L_{y^T}x_i} = \sum_{j=1}^n a_{ji}\partial_{x_j}.
\]
设 $\phi \in \mathcal{C}^{\infty}(\R^n)$, $z = (z_1, \ldots, z_n) \in \R^n$. 有
\[
\partial_{L_{y^T}x_i}\phi = \sum_{j=1}^n a_{ji}\frac{\partial\phi}{\partial x_j},
\]
所以
\[
L_y\partial_{L_{y^T}x_i}\phi(z) = \sum_{j=1}^n a_{ji}\frac{\partial\phi}{\partial x_j}(az).
\]
另一方面, $L_y\phi(z) = \phi(az)$, 其中
\[
az = \Bigl(\sum_{j=1}^n a_{rj}z_r\Bigr)_{1 \leq r \leq n},
\]
所以
\[
(\partial_{x_i}L_y\phi)(z) = \sum_{r=1}^n a_{ri}\partial_{x_r}\phi(az).
\]
我们看到 $L_y\partial_{L_{y^T}x_i}\phi(z)$ 和 $(\partial_{x_i}L_y\phi)(z)$ 确实得到相同的结果.
(d). 设 $f, g \in V_m(\R^n)$, $y \in \Or(n)$. 由 (c),
\[
\langle L_yf, L_yg \rangle = \partial_{L_yg}L_yf = L_y\partial_{L_{y^T}L_yg}f = L_y\partial_g f.
\]
由于 $\partial_g f$ 是常值函数, 有
\[
L_y\partial_g f = \partial_g f = \langle f, g \rangle,
\]
这正是我们想要的.
(e). 若 $m \leq 1$, 则 $\Delta = 0$ 在 $V_m(\R^n)$ 上, 所以 $\mathcal{H}_m(\R^n) = \Ker(\Delta) = V_m(\R^n)$.
(f). 注意 $\Delta = \partial_{|x|^2}$. 所以若 $f \in V_m(\R^n)$, $g \in V_{m-2}(\R^n)$, 有
\[
\langle f, |x|^2g \rangle = \partial_{|x|^2g}f = \partial_g(\partial_{|x|^2}f) = \langle \Delta f, g \rangle.
\]
换句话说, 映射 $V_{m-2}(\R^n) \to V_m(\R^n)$, $g \longmapsto |x|^2g$ 是 $\Delta \colon V_m(\R^n) \to V_{m-2}(\R^n)$ 的伴随, 所以它的像是 $\Ker\Delta = \mathcal{H}_m(\R^n)$ 的正交补.
(g). 第一个公式只是由 (f) 做简单归纳得出. 对第二个公式, 注意对每个 $k \in \{0, \ldots, \lfloor\frac{m}{2}\rfloor\}$, 单射线性变换 $V_{m-2}(\R^n) \to V_m(\R^n)$, $g \longmapsto |x|^{2k}g$ 是 $\Or(n)$-等变的, 因为函数 $|x|^{2k}$ 被 $\Or(n)$ 不变.
(h). 我们在 $\R^n$ 上用极坐标: $\R^n$ 的点 $z$ 可以写成 $z = rs$, 其中 $r \in \R_{\geq 0}$, $s \in S$, 且 $z \neq 0$ 时 $r$ 和 $s$ 唯一确定.
设 $f \in \mathcal{H}_m(\R^n)$, $g \in \mathcal{H}_p(\R^n)$. 则若 $r \in \R_{\geq 0}$, $s \in S$, 有 $f(rs) = r^mf(s)$, $g(rs) = r^pg(s)$. 由 Green 第二公式,
\[
\int_B (f\Delta g - g\Delta f)\,\dd\lambda = c\int_S\Bigl(f\frac{\partial g}{\partial r} - g\frac{\partial f}{\partial r}\Bigr)\,\dd\mu,
\]
其中 $B$ 是闭单位球, $\lambda$ 是 $\R^n$ 上的 Lebesgue 测度, $c$ 是正常数, $\mu$ 是习题 \ref{ex:I.5.3.5}(f) 中定义的 $S$ 上的测度. 由于 $f$ 和 $g$ 在 $\Delta$ 的核中, 这给出
\[
0 = c(p - m)\int_S f(s)\overline{g(s)}\,\dd\mu(s).
\]
若 $m \neq p$, 得到 $\int_S f(s)\overline{g(s)}\,\dd\mu(s) = 0$. 所以对 $m \neq p$, 子空间 $\mathcal{H}_m(\R^n)|_S$ 和 $\mathcal{H}_p(\R^n)|_S$ 对 $L^2(S, \mu)$ 的内积正交. 特别地, 若 $f \in \bigoplus_{m \geq 0}\mathcal{H}_m(\R^n)$ 满足 $f|_S = 0$, 则把 $f$ 写成 $f = \sum_{m \geq 0}f_m$ (其中 $f_m \in \mathcal{H}_m(\R^n)$), 必有对每个 $m \geq 0$ 有 $f_m|_S = 0$. 但 $f_m$ 是 $m$ 次齐次的, 所以对每个 $r \in \R_{\geq 0}$ 和 $s \in S$ 有 $f_m(rs) = r^mf_m(s)$, 因此 $f_m|_S = 0$ 蕴含 $f_m = 0$.
(i). 存在性由 (g) 的第一个恒等式得出 (因为在 $S$ 上 $|x|^2 = 1$), 唯一性由 (h) 得出.
\end{solution}

\begin{exercise}[$(\SO(n), \SO(n-1))$ 是 Gelfand 对]
\label{ex:V.8.1.3}
设 $n \geq 2$, 通过映射 $x \longmapsto \begin{pmatrix} 1 & 0 \\ 0 & x \end{pmatrix}$ 把 $\Or(n-1)$ 嵌入 $\Or(n)$. 设 $G = \SO(n)$, $K$ 是 $\SO(n-1)$ 在 $G$ 中通过刚定义的嵌入的像.
\begin{enumerate}
\item[(a)] 设 $A$ 是 $\SO(n)$ 中形如
\[
\begin{pmatrix} \cos\theta & 0 & \sin\theta \\ 0 & I_{n-2} & 0 \\ -\sin\theta & 0 & \cos\theta \end{pmatrix}
\]
($\theta \in \R$) 的矩阵的子集.
证明 $A$ 是 $G$ 的子群且 $G = KAK$.
\item[(b)] 证明 $(G, K)$ 是 Gelfand 对. (你可能想用 $G$ 的对合 $\theta(x) = JxJ$, 其中 $J$ 是对角线系数为 $-1, 1, \ldots, 1$ 的对角矩阵.)
\end{enumerate}
\end{exercise}

\begin{solution}
(a). 对每个 $\theta \in \R$, 记 $A_{\theta} = \begin{pmatrix} \cos\theta & 0 & \sin\theta \\ 0 & I_{n-2} & 0 \\ -\sin\theta & 0 & \cos\theta \end{pmatrix}$. 有 $A = \{A_{\theta} \mid \theta \in \R\}$. 容易验证 $A_{\theta}A_{\theta'} = A_{\theta + \theta'}$, 所以 $A$ 是 $G$ 的子群.
设 $v_0 = (1, 0, \ldots, 0) \in S$. 则 $\Or(n)$ 在 $\R^n$ 上的作用保持 $S$ (且 $\Or(n)$ 在 $S$ 上传递作用), $K$ 是 $v_0$ 在 $\Or(n)$ 中的稳定化子. 设 $x \in \Or(n)$, 写 $z = x \cdot v_0 = (z_1, \ldots, z_n)$. 可以找到 $y \in K$ 使得 $y \cdot z = (z_1, 0, \ldots, 0, c)$, 其中 $c^2 = z_2^2 + \cdots + z_n^2$, 然后可以找到 $a \in A$ 使得 $a \cdot (y \cdot z) = (1, 0, \ldots, 0) = v_0$. 则 $(ayx) \cdot v_0 = v_0$, 所以 $ayx \in K$, 且 $x \in Ka^{-1}y^{-1} \subset KAK$.
(b). 想对 $\theta$ 应用命题 \ref{prop:V.2.5}, 其中 $\theta$ 把 $x \in G$ 送到 $JxJ$, $J = \begin{pmatrix} -1 & 0 \\ 0 & I_{n-1} \end{pmatrix}$. 注意 $J^2 = I_n$, 所以 $J = J^{-1}$, 因此 $\theta$ 是群态射且是对合. 显然也有 $\theta(K) = K$. 需要验证对每个 $x \in G$ 有 $\theta(x) \in Kx^{-1}K$. 设 $x \in G$, 写 $x = kak'$, 其中 $k, k' \in K$, $a \in A$. 则 $\theta(x) = \theta(k)\theta(a)\theta(k')$, $\theta(k), \theta(k') \in K$, 且若 $a = A_{\theta}$, 则 $\theta(a) = A_{-\theta} = a^{-1}$. 所以 $\theta(x) = \theta(k)k'x^{-1}k\theta(k') \in Kx^{-1}K$.
\end{solution}

\begin{exercise}[调和表示]
\label{ex:V.8.1.4}
我们用习题 \ref{ex:V.8.1.1} 和 \ref{ex:V.8.1.2} 的记号, 以及习题 \ref{ex:V.8.1.3} 中定义的嵌入 $\Or(n-1) \subset \Or(n)$.
\begin{enumerate}
\item[(a)] 证明有 $\Or(n-1)$-等变同构
\[
V_m(\R^n) \simeq \bigoplus_{k=0}^m V_{m-k}(\R^{n-1}).
\]
\item[(b)] 证明有 $\Or(n-1)$-等变同构
\[
\mathcal{H}_m(\R^n) \simeq \bigoplus_{k=0}^m\mathcal{H}_{m-k}(\R^{n-1}).
\]
\item[(c)] 若 $m \geq 2$, 证明 $\mathcal{H}_m(\R^2)$ 是 $\Or(2)$ 的不可约表示, 但作为 $\SO(2)$ 的表示不可约.
\end{enumerate}
从现在起, 我们假设 $n \geq 3$.
\begin{enumerate}
\item[(a)] 若 $m \geq 1$, 证明 $\mathcal{H}_m(\R^n)^{\SO(n)} = \{0\}$.
\item[(b)] 证明对每个 $m \geq 0$, 空间 $\mathcal{H}_m(\R^n)^{\SO(n-1)}$ 是 $1$ 维的.
\item[(c)] 设 $S \subset \R^n$ 是单位球面, $v_0 = (1, 0, \ldots, 0) \in S$. 证明映射 $\SO(n) \to S$, $x \longmapsto x \cdot v_0$ 诱导同胚 $\SO(n)/\SO(n-1) \xrightarrow{\sim} S$.
\item[(d)] 证明习题 \ref{ex:I.5.3.5}(f) 中定义的 $S$ 上的测度 $\mu$ (用 $\SO(n)$ 和 $\SO(n-1)$ 上的规范化 Haar 测度) 由 $\mu(E) = c\lambda(\{tx \mid t \in [0,1], x \in E\})$ 给出, 其中 $\lambda$ 是 $\R^n$ 上的 Lebesgue 测度, $c^{-1}$ 是单位球 (对 $\lambda$) 的体积.
\item[(e)] 由上一问, 有 $\SO(n)$ 在 $L^2(S)$ 上的拟正则表示, 且它保持连续函数的子空间. 若 $V \subset \mathcal{C}(S)$ 是非零有限维 $\SO(n)$-稳定子空间, 证明 $V^{\SO(n-1)} \neq \{0\}$. (提示: 从满足 $f(v_0) \neq 0$ 的函数 $f \in V$ 开始.)
\item[(f)] 证明 $\SO(n)$ 的表示 $\mathcal{H}_m(\R^n)$ 不可约.
\item[(g)] 证明 $\SO(n)$ 的表示 $\mathcal{H}_m(\R^n)$ 和 $\mathcal{H}_{m'}(\R^n)$ 在 $m \neq m'$ 时不等价. (提示: 比较维数.)
\item[(h)] 若 $m \geq 2$, 证明 $\mathcal{H}_m(\R^n)$ 由函数 $(z_1, \ldots, z_n) \longmapsto (a_1z_1 + \cdots + a_nz_n)^m$ ($a_1, \ldots, a_n \in \CC$ 满足 $a_1^2 + \cdots + a_n^2 = 0$) 张成.
\end{enumerate}
\end{exercise}

\begin{solution}
(a). 若 $f \in V_m(\R^n)$, 则可以把 $f$ 写成 $f = \sum_{k=0}^m x_1^kf_k$, 其中 $f_k \in V_{m-k}(\R^{n-1})$ 唯一确定. 由于 $\Or(n-1) \subset \Or(n)$ 平凡作用于 $x_1$, 这给出 $\Or(n-1)$-等变同构 $V_m(\R^n) \simeq \bigoplus_{k=0}^m V_{m-k}(\R^{n-1})$.
(b). 在此证明中, 约定 $m < 0$ 时 $V_m(\R^n) = 0$. 固定 $n$ 和 $m$. 由习题 \ref{ex:V.8.1.2}(f) 和 (g), 有 $\Or(n)$-等变同构 $V_m(\R^n) \simeq \mathcal{H}_m(\R^n) \oplus V_{m-2}(\R^n)$. 用 (a), 由此推出 $\Or(n-1)$-等变同构
\[
V_m(\R^n) \simeq \mathcal{H}_m(\R^n) \oplus \bigoplus_{k=0}^{m-2}V_{m-2-k}(\R^{n-1}).
\]
另一方面, 把 (a) 应用于 $V_m(\R^n)$ 给出 $\Or(n-1)$-等变同构 $V_m(\R^n) \simeq \bigoplus_{k=0}^m V_{m-k}(\R^{n-1})$. 对每个和项用习题 \ref{ex:V.8.1.2}(f) 或 (g), 得到 $\Or(n-1)$-等变同构
\begin{align*}
V_m(\R^n) &\simeq \bigoplus_{k=0}^m(\mathcal{H}_{m-k}(\R^{n-1}) \oplus V_{m-2-k}(\R^{n-1})) \\
&= \Bigl(\bigoplus_{k=0}^m\mathcal{H}_{m-k}(\R^{n-1})\Bigr) \oplus \Bigl(\bigoplus_{k=0}^{m-2}V_{m-2-k}(\R^{n-1})\Bigr).
\end{align*}
定义 $\Or(n-1)$ 的表示 $V_1, V_2, V_3$ 为 $V_1 = \mathcal{H}_m(\R^n)$, $V_2 = \bigoplus_{k=0}^m\mathcal{H}_{m-k}(\R^{n-1})$, $V_3 = \bigoplus_{k=0}^{m-2}V_{m-2-k}(\R^{n-1})$. 我们刚看到 $V_1 \oplus V_3 \simeq V_2 \oplus V_3$, 所以 $\chi_{V_1} + \chi_{V_3} = \chi_{V_2} + \chi_{V_3}$, 从而 $\chi_{V_1} = \chi_{V_2}$. 由推论 \ref{cor:IV.5.10}, 这蕴含 $V_1$ 和 $V_2$ 等价.
(c). 注意 $\SO(2)$ 是交换群 (它同构于 $S^1$), 所以它的不可约表示都是 $1$ 维的. 另一方面, 由习题 \ref{ex:V.8.1.2}(f) 和习题 \ref{ex:V.8.1.1}(a),
\[
\dim\mathcal{H}_m(\R^2) = \dim V_m(\R^2) - \dim V_{m-2}(\R^2) = \binom{m+1}{m} - \binom{m-1}{m-2} = 2 > 1
\]
若 $m \geq 2$, 所以 $\mathcal{H}_m(\R^n)$ 不能是 $\SO(2)$ 的不可约表示.
若 $m = 2$, $\mathcal{H}_m(\R^2)$ 的基由函数 $x_1^2 - x_2^2$ 和 $x_1x_2$ 给出, 它们都张成被 $\Or(2)$ 的作用稳定的直线, 所以 $\mathcal{H}_2(\R^2)$ 不是 $\Or(2)$ 的不可约表示.
假设 $m \geq 3$. 由于 $\dim\mathcal{H}_m(\R^2) = 2$, 若 $\mathcal{H}_m(\R^2)$ 不是 $\Or(2)$ 的不可约表示, 必有非零的 $f \in \mathcal{H}_m(\R^2)$ 使得对每个 $x \in \Or(2)$ 有 $L_xf \in \CC f$. 按通常方式把 $\R^2$ 等同于复平面 $\CC$. 则对 $z \in \CC$, $f(z)$ 可以写成 $f(z) = \sum_{r=0}^m a_rz^r\bar{z}^{m-r}$, 其中 $a_0, \ldots, a_m \in \CC$. $\SO(2)$ 的作用变成 $S^1$ 在 $\CC$ 上通过乘法的作用, $\begin{pmatrix} 1 & 0 \\ 0 & -1 \end{pmatrix}$ 的作用对应于复共轭. 由对 $f$ 的假设, 对每个 $u \in S^1$, 函数 $f(uz) = \sum_{r=0}^m a_ru^{2r-m}z^r\bar{z}^{m-r}$ 是 $f$ 的倍数. 这只有在存在 $r \in \{0, \ldots, m\}$ 使得对 $s \neq r$ 有 $a_s = 0$ 才可能. 所以可以假设 $f(z) = z^r\bar{z}^{m-r}$. 函数 $\bar{f}(z) = z^{m-r}\bar{z}^r$ 也是 $f$ 的倍数, 所以必有 $m = 2r$, $f = |x|^m$. 则 $\Delta f = m(m-1)|x|^{m-2}$, 与 $\Delta f = 0$ 矛盾.
(d). 设 $f \in V_m(\R^n)^{\SO(n)}$. 由于 $f$ 被 $\SO(n)$ 不变, 它在以 $0$ 为中心的球面上是常数, 所以对每个 $z \in \R^n$ 有 $f(z) = f(\|z\|v_0)$. 由于 $f$ 是 $m$ 次齐次的, 得到对每个 $z \in \R^n$ 有 $f(z) = \|z\|^mf(v_0)$. 所以 $f$ 是多项式当且仅当 $m$ 偶. 又, 容易验证 $\Delta f = m(m-1)|x|^{-2}f$, 所以 $f \in \mathcal{H}_m(\R^n)$ 当且仅当 $f = 0$ 或 $m = 0$.
(e). 由 (b), 有 $\SO(n-1)$-等变同构 $\mathcal{H}_m(\R^n) \simeq \bigoplus_{k=0}^m\mathcal{H}_{m-k}(\R^{n-1})$. 所以由 (d),
\[
\mathcal{H}_m(\R^n)^{\SO(n-1)} \simeq \mathcal{H}_0(\R^{n-1})^{\SO(n-1)} \simeq \CC.
\]
(f). 记映射 $\SO(n) \to S$, $x \longmapsto x \cdot v_0$ 为 $\phi$. 首先, 此映射显然连续, 且它满射, 因为 $\SO(n)$ 在 $S$ 上传递作用. (若 $v_1, v'_1 \in S$, 要找 $x \in \SO(n)$ 使得 $x \cdot v_1 = v'_1$. 可以把 $v_1$ 和 $v'_1$ 完备成 $\R^n$ 的两个标准正交基 $(v_1, \ldots, v_n)$ 和 $(v'_1, \ldots, v'_n)$. 这两个基之间的换基矩阵在 $\Or(n)$ 中. 若它在 $\SO(n)$ 中, 我们完成了. 否则, $(v_1, \ldots, v_n)$ 和 $(v'_1, \ldots, v'_{n-1}, -v'_n)$ 之间的换基矩阵将在 $\SO(n)$ 中.)
又, $v_0$ 在 $\SO(n)$ 中的稳定化子是 $\SO(n)$ 中第一列为 $\begin{pmatrix} 1 \\ 0 \\ \vdots \\ 0 \end{pmatrix}$ 的元素构成的子群, 容易看到这就是 $\SO(n-1)$. 所以 $\phi$ 诱导连续双射 $\SO(n)/\SO(n-1) \to S$. 由于 $\SO(n)/\SO(n-1)$ 紧, 此映射是同胚.
(g). 定义 $S$ 上的正则 Borel 测度 $\nu$:
\[
\nu(E) = c\lambda(\{tx \mid t \in [0,1], x \in E\}).
\]
定义线性泛函 $I \colon \mathcal{C}(\SO(n)) \to \CC$:
\[
I(f) = \int_S f^{\SO(n-1)}(s)\,\dd\nu(s),
\]
其中
\[
f^{\SO(n-1)}(x) = \int_{\SO(n-1)}f(xy)\,\dd y
\]
(用 $\SO(n-1)$ 上的规范化 Haar 测度), 对每个 $s \in \SO(n)$; 函数 $f^{\SO(n-1)}$ 被 $\SO(n-1)$ 右不变, 所以可以由 (f) 等同于 $S$ 上的函数.
这是 $\mathcal{C}(\SO(n))$ 上的正泛函, 所以它来自 $\SO(n)$ 上的正则 Borel 测度, 设为 $\rho$. 要证明 $\rho$ 是 $\SO(n)$ 上的规范化 Haar 测度.
注意若 $f$ 是常值函数 $1$, 则 $I(f) = 1$. 所以为证明 $\rho$ 是规范化 Haar 测度, 只需证明它左不变. 设 $f \in \mathcal{C}(\SO(n))$, $y \in \SO(n)$. 则由定义立即得到 $(L_yf)^{\SO(n-1)} = L_y(f^{\SO(n-1)})$, 所以只需证明 $S$ 上的测度 $\nu$ 被 $\SO(n)$ 的作用左不变. 但这由 Lebesgue 测度 $\lambda$ 被 $\SO(n)$ 的作用左不变 (这可以用换元公式看到) 这一事实立即得出.
(h). 设 $V$ 如题设. 由于 $\SO(n)$ 在 $S$ 上传递作用且 $V$ 被 $\SO(n)$ 稳定, 可以找到 $f \in V$ 使得 $f(v_0) \neq 0$. 设 $(f_1, \ldots, f_r)$ 是 $V$ 的基. 由于 $V$ 被 $\SO(n)$ 稳定且 $\SO(n)$ 在 $V$ 上的作用连续, 可以找到连续函数 $c_1, \ldots, c_r \colon \SO(n) \to \CC$ 使得对每个 $x \in \SO(n)$ 和每个 $s \in S$, $f(x \cdot s) = \sum_{i=1}^r c_i(x)f_i(s)$. 定义 $\tilde{f} \colon S \to \CC$:
\[
\tilde{f}(s) = \int_{\SO(n-1)}f(x \cdot s)\,\dd x.
\]
则 $\tilde{f} = \sum_{i=1}^r\bigl(\int_{\SO(n-1)}c_i(x)\,\dd x\bigr)f_i$, 所以 $\tilde{f} \in V$. 又, $\tilde{f}$ 由构造是 $\SO(n-1)$-不变的. 最后, 由于对每个 $x \in \SO(n-1)$ 有 $x \cdot v_0 = v_0$, 有 $\tilde{f}(v_0) = f(v_0) \neq 0$, 所以 $\tilde{f} \neq 0$.
(i). 由习题 \ref{ex:V.8.1.2}(h), 从 $\R^n$ 到 $S$ 的限制在 $\mathcal{H}_m(\R^n)$ 上单射, 所以 $\mathcal{H}_m(\R^n)$ 作为 $\SO(n)$ 的表示不可约当且仅当 $\mathcal{H}_m = \mathcal{H}_m(\R^n)|_S \subset \mathcal{C}(S)$ 作为 $\SO(n)$ 的表示不可约. 由于 $\SO(n)$ 紧, 若 $\mathcal{H}_m$ 不可约, 则可以把 $\mathcal{H}_m$ 写成 $\mathcal{H}_m = V \oplus V'$, 其中 $V$ 和 $V'$ 是 $\mathcal{H}_m$ 的非零 $\SO(n)$-不变子空间. 由 (h), 这蕴含 $\dim(\mathcal{H}_m^{\SO(n)}) \geq 2$, 与 (d) 矛盾. 所以 $\mathcal{H}_m$ 不可约.
(j). 设 $d_m = \dim\mathcal{H}_m(\R^n)$. 将证明对每个 $m \in \Z_{\geq 0}$ 有 $d_{m+1} > d_m$, 这蕴含 $m \neq m'$ 时 $d_m \neq d_{m'}$, 所以 $\mathcal{H}_m(\R^n)$ 和 $\mathcal{H}_{m'}(\R^n)$ 不等价.
若 $m \leq 1$, 由习题 \ref{ex:V.8.1.1}(a) 和习题 \ref{ex:V.8.1.2}(e), 有 $d_m = \dim V_m(\R^n) = \binom{m+n-1}{m}$. 若 $m \geq 2$, 由习题 \ref{ex:V.8.1.1}(a) 和习题 \ref{ex:V.8.1.2}(f),
\begin{align*}
d_m &= \dim V_m(\R^n) - \dim V_{m-2}(\R^n) \\
&= \binom{m+n-1}{m} - \binom{m+n-3}{m-2} \\
&= \frac{(m+n-3)!}{(m-2)!(n-1)!}\Bigl(\frac{(m+n-1)(m+n-2)}{m(m-1)} - 1\Bigr) \\
&= \frac{(2m+n-2)(m+n-3)!}{m!(n-2)!}.
\end{align*}
特别地, $d_0 = 1$, $d_1 = n$, $d_2 = 2n-1$, 所以 $d_2 > d_1 > d_0$. 设 $m \geq 2$. 则
\begin{align*}
d_{m+1} - d_m &= \frac{(2(m+1)+n-2)(m+1+n-3)!}{(m+1)!(n-2)!} - \frac{(2m+n-2)(m+n-3)!}{m!(n-2)!} \\
&= \frac{(m+n-3)!}{m!(n-2)!}\Bigl(\frac{(2m+n)(m+n-2)}{m+1} - (2m+n-2)\Bigr) \\
&= \frac{(m+n-3)!}{m!(n-2)!}\cdot\frac{(2m+n)(m+n-2) - (2m+n-2)(m+1)}{m+1} \\
&> 0
\end{align*}
(因为 $m + n - 2 \geq m + 1$ 且 $2m + n > 2m + n - 2$).
(k). 设 $a_1, \ldots, a_n \in \CC$, 考虑 $f = (a_1x_1 + \cdots + a_nx_n)^m \in V_m(\R^n)$. 则 $\Delta f = m(m-1)(a_1^2 + \cdots + a_n^2)(a_1x_1 + \cdots + a_nx_n)^{m-2}$, 所以 $f \in \mathcal{H}_m(\R^n)$ 当且仅当 $a_1^2 + \cdots + a_n^2 = 0$. 设
\[
W = \operatorname{Span}\{(a_1x_1 + \cdots + a_nx_n)^m \mid a_1, \ldots, a_n \in \CC, a_1^2 + \cdots + a_n^2 = 0\} \subset \mathcal{H}_m(\R^n).
\]
由于 $W \neq 0$ 且 $\mathcal{H}_m(\R^n)$ 作为 $\SO(n)$ 的表示不可约, 为证明 $W = \mathcal{H}_m(\R^n)$ 只需证明 $W$ 被 $\SO(n)$ 不变. 设 $x \in \SO(n)$, $a_1, \ldots, a_n \in \CC$, 设 $f = (a_1x_1 + \cdots + a_nx_n)^m$. 写 $(b_1, \ldots, b_n) = (a_1, \ldots, a_n)x^T$ (把 $(a_1, \ldots, a_n)$ 看作行向量). 则 $b_1^2 + \cdots + b_n^2 = 0$, 因为 $x \in \Or(n)$, 且 $L_xf = (b_1x_1 + \cdots + b_nx_n)^m$, 所以 $L_xf \in W$.
\end{solution}

\begin{exercise}[球函数]
\label{ex:V.8.1.5}
我们保留习题 \ref{ex:V.8.1.1}-\ref{ex:V.8.1.4} 的记号, 假设 $n \geq 3$.
\begin{enumerate}
\item[(a)] 证明空间 $\sum_{m \geq 0}\mathcal{H}_m(\R^n)|_S$ 在 $L^2(S)$ 中稠密且和是直和.
\item[(b)] 证明 $L^2(S)$ 的子空间 $\mathcal{H}_m(\R^n)|_S$ 和 $\mathcal{H}_{m'}(\R^n)|_S$ 在 $m \neq m'$ 时正交 (对 $L^2$ 内积).
\item[(c)] 证明 $\SO(n)$ 的每个具有非零 $\SO(n-1)$-不变向量的不可约酉表示都同构于某个 $\mathcal{H}_m(\R^n)$.
\end{enumerate}
\end{exercise}

\begin{solution}
(a). 由习题 \ref{ex:V.8.1.2}(i),
\[
\sum_{m \geq 0}\mathcal{H}_m(\R^n)|_S = \sum_{m \geq 0}V_m(\R^n)|_S,
\]
且右端由 Stone-Weierstrass 定理在 $\mathcal{C}(S)$ 中稠密 (从而在 $L^2(S)$ 中稠密). 又, 我们在习题 \ref{ex:V.8.1.2}(h) 的证明中看到空间 $\mathcal{H}_m(\R^n)|_S$ 在 $L^2(S)$ 中两两正交, 所以它们在直和中.
(b). 见 (a).
(c). 设 $V$ 是 $\SO(n)$ 的不可约酉表示, 满足 $V^{\SO(n-1)} \neq 0$. 由定理 \ref{thm:V.3.2.4}, $V$ 是 $L^2(S)$ 的子表示. 但我们已经看到
\[
L^2(S) = \bigoplus_{m \geq 0}\mathcal{H}_m(\R^n)|_S
\]
且所有这些和项都不可约, 所以 $V$ 同构于其中一个.
\end{solution}

\begin{exercise}[\index{Gegenbauer 多项式}Gegenbauer 多项式]
\label{ex:V.8.1.6}
我们保留习题 \ref{ex:V.8.1.1}-\ref{ex:V.8.1.5} 的记号. 若函数 $\phi \in \mathcal{C}(S)$ 被 $\SO(n-1)$ 左不变, 则称它带状 (zonal). (由于 $S = \SO(n)/\SO(n-1)$, 也可以把 $\phi$ 看成 $\SO(n)$ 上的双不变函数.) 假设 $n \geq 3$.
\begin{enumerate}
\item[(a)] 证明 $\phi \in \mathcal{C}(S)$ 带状当且仅当存在连续函数 $f \colon [-1,1] \to \CC$ 使得对每个 $z = (z_1, \ldots, z_n) \in S$, $\phi(z) = f(z_1)$.
\item[(b)] 证明存在 $c \in \R_{>0}$ 使得对每个带状 $\phi \in \mathcal{C}(S)$, 若定义 $f \colon [-1,1] \to \CC$ 如 (a), 则
\[
\int_S\phi(z)\,\dd\mu(z) = c\int_{-1}^1 f(t)(1-t^2)^{(n-3)/2}\,\dd t.
\]
(提示: 可以试着用球坐标, 如 https://en.wikipedia.org/wiki/N-sphere\#Spherical\_coordinates.)
\item[(c)] 设 $m \geq 0$. 若 $t \in S$, 设 $f_t$ 是 $\mathcal{H}_m(\R^n)$ 中满足对每个 $g \in \mathcal{H}_m(\R^n)$ 有 $\langle g, f_t \rangle = g(t)$ 的唯一元素. (注意我们用习题 \ref{ex:V.8.1.2} 的内积.)
证明函数 $Z_m = f_{v_0}|_S$ (其中 $v_0 = (1, 0, \ldots, 0)$) 是\index{带状函数}带状函数.
\item[(d)] 设 $f_m \colon [-1,1] \to \CC$ 是 (a) 中对应于 $Z_m$ 的连续函数. 证明 $f_m$ 是次数 $\leq m$ 的多项式函数.
\item[(e)] 若 $m \neq m'$, 证明 $\int_{-1}^1 f_m(t)f_{m'}(t)(1-t^2)^{(n-3)/2}\,\dd t = 0$.
\item[(f)] 证明 $f_m$ 的次数是 $m$.
\item[(g)] 证明 $x \longmapsto \frac{1}{Z_m(v_0)}Z_m(x \cdot v_0)$ 是 $\SO(n)$ 上的球函数, 且每个球函数都具有此形式.
\end{enumerate}
多项式 $\frac{1}{f_m(1)}f_m$ 称为 Gegenbauer 多项式 (若 $n = 3$ 也称为 \index{Legendre 多项式}Legendre 多项式).
我们现在给出球函数的另一个公式.
\begin{enumerate}
\item[(h)] 考虑函数 $h_m \in V_m(\R^n)$, $h_m(z_1, \ldots, z_n) = (z_1 + iz_2)^m$. 证明 $h_m \in \mathcal{H}_m(\R^n)$.
\item[(i)] 定义函数 $\psi_m \colon S \to \CC$: $\psi_m(z) = \int_{\SO(n-1)}h_m(k \cdot z)\,\dd k$. 证明 $\psi_m$ 被 $\SO(n-1)$ 左不变, $\psi_m \in \mathcal{H}_m(\R^n)|_S$ 且 $\psi_m(v_0) = 1$.
\item[(j)] 证明 $\SO(n)$ 上的每个球函数都具有 $x \longmapsto \psi_m(x \cdot v_0)$ 的形式, 其中 $m \geq 0$ 唯一.
\end{enumerate}
我们可以计算定义 $\psi_m$ 的积分, 得到
\[
\psi_m(\cos\phi, z_2, \ldots, z_n) = \frac{\Gamma(\frac{n-1}{2})}{\sqrt{\pi}\Gamma(\frac{n-2}{2})}\int_0^{\pi}(\cos\phi + i\sin\phi\cos\theta)^m\sin^{n-3}\theta\,\dd\theta.
\]
\end{exercise}

\begin{solution}
(a). 若存在连续函数 $f \colon [-1,1] \to \CC$ 使得 $\phi(z_1, \ldots, z_n) = f(z_1)$, 则 $\phi$ 显然带状.
反之, 假设 $\phi$ 带状, 定义 $f \colon [-1,1] \to \CC$:
\[
f(z_1) = \phi(z_1, 0, \ldots, 0, \sqrt{1-z_1^2}).
\]
则 $f$ 显然连续. 设 $s = (z_1, \ldots, z_n) \in S$. 则存在 $x \in \SO(n-1)$ 使得 $x \cdot s = (z_1, 0, \ldots, 0, \sqrt{1-z_1^2})$ (用 $\SO(n-1)$ 在 $\R^{n-1}$ 的任何球面上传递作用这一事实). 由于 $\phi$ 带状, 有 $\phi(s) = f(z_1)$.
(b). 我们在 $\R^n$ 上用球坐标. 即给定 $(z_1, \ldots, z_n) \in \R^n$, 写
\begin{align*}
z_1 &= r\cos\varphi_1 \\
z_2 &= r\sin\varphi_1\cos\varphi_2 \\
&\vdots \\
z_{n-1} &= r\sin\varphi_1\sin\varphi_2\cdots\sin\varphi_{n-2}\cos\varphi_{n-1} \\
z_n &= r\sin\varphi_1\sin\varphi_2\cdots\sin\varphi_{n-2}\sin\varphi_{n-1},
\end{align*}
其中 $r = \sqrt{z_1^2 + \cdots + z_n^2} \in \R_{\geq 0}$, $\varphi_1, \ldots, \varphi_{n-2} \in [0, \pi]$, $\varphi_{n-1} \in [0, 2\pi)$ ($\varphi_1, \ldots, \varphi_{n-1}$ 一般不唯一确定, 但例如 $z_1, \ldots, z_n$ 都非零时唯一). 若 $dz$ 是 $\R^n$ 上的 Lebesgue 测度, 则
\[
dz = r^{n-1}\sin^{n-2}\varphi_1\sin^{n-3}\varphi_2\cdots\sin\varphi_{n-2}\,\dd r\,\dd\varphi_1\cdots\dd\varphi_{n-1}.
\]
设 $\phi \in \mathcal{C}(S)$ 带状. 差一个正实常数, $\int_S\phi(s)\,\dd\mu(s)$ 等于 $\int_{B-\{0\}}\psi(z)\,\dd z$, 其中 $B$ 是闭单位球, $\psi(z) = \phi(\frac{z}{\|z\|})$ ($z \neq 0$). 这等于 $\int_0^1 r^{n-1}\,\dd r$ (另一个正实常数) 与
\begin{multline*}
\int_0^{\pi}\cdots\int_0^{\pi}\int_0^{2\pi}\phi(\cos\varphi_1, \sin\varphi_1\cos\varphi_2, \ldots, \sin\varphi_1\cdots\sin\varphi_{n-2}\sin\varphi_{n-1}) \\
\times \sin^{n-2}\varphi_1\sin^{n-3}\varphi_2\cdots\sin\varphi_{n-2}\,\dd\varphi_1\cdots\dd\varphi_{n-1}
\end{multline*}
的乘积. 由于 $\phi$ 带状, 上面的大积分等于
\[
\int_0^{\pi}\cdots\int_0^{\pi}\int_0^{2\pi}f(\cos\varphi_1)\sin^{n-2}\varphi_1\sin^{n-3}\varphi_2\cdots\sin\varphi_{n-2}\,\dd\varphi_1\cdots\dd\varphi_{n-1}.
\]
差常数
\[
\int_0^{\pi}\cdots\int_0^{\pi}\int_0^{2\pi}\sin^{n-3}\varphi_2\cdots\sin\varphi_{n-2}\,\dd\varphi_2\cdots\dd\varphi_{n-1}
\]
(它必须为正, 因为它计算的是常值函数 $1$ 的积分, 差一个正常数), 这等于
\[
\int_0^{\pi}f(\cos\varphi_1)\sin^{n-2}\varphi_1\,\dd\varphi_1.
\]
最后用代换 $t = \cos\varphi_1$. 有 $dt = -\sin\varphi_1\,\dd\varphi_1$, 所以
\[
\int_0^{\pi}f(\cos\varphi_1)\sin^{n-2}\varphi_1\,\dd\varphi_1 = \int_1^{-1}f(t)(-\sqrt{1-t^2})^{n-3}\,\dd t = \int_{-1}^1 f(t)(1-t^2)^{(n-3)/2}\,\dd t.
\]
(c). 设 $y \in \SO(n-1)$. 由 $f_{v_0}$ 的定义, 对每个 $g \in \mathcal{H}_m(\R^n)$,
\[
\langle g, L_yf_{v_0} \rangle = \langle L_{y^{-1}}g, f_{v_0} \rangle = L_yg(v_0) = g(y^{-1}v_0) = g(v_0) = \langle g, f_{v_0} \rangle
\]
(因为 $\SO(n-1)$ 是 $v_0$ 在 $\SO(n)$ 中的稳定化子). 所以 $f_{v_0} - L_yf_{v_0}$ 与 $\mathcal{H}_m(\R^n)$ 的每个元素正交, 这蕴含 $f_{v_0} - L_yf_{v_0} = 0$, 即 $f_{v_0} = L_yf_{v_0}$. 所以 $Z_m$ 带状.
(d). 函数 $Z_m$ 是 $f_{v_0}$ 的限制, 而 $f_{v_0} \in \mathcal{H}_m(\R^n)$, 特别地是 $m$ 次多项式函数. 由于 $f_m$ 由 $f_m(t) = Z_m(t, 0, \ldots, 0, \sqrt{1-t^2})$ 定义, 可以把 $f_m$ 写成
\[
f_m(t) = \sum_{k=0}^m c_kt^k(1-t^2)^{(m-k)/2},
\]
其中 $c_0, \ldots, c_m \in \CC$. 但又有 $f_m(t) = Z_m(t, 0, \ldots, 0, -\sqrt{1-t^2})$, 所以
\[
\sum_{k=0}^m c_kt^k(1-t^2)^{(m-k)/2} = \sum_{k=0}^m(-1)^{m-k}c_kt^k(1-t^2)^{(m-k)/2}.
\]
这迫使除非 $m - k$ 偶, $c_k = 0$, 所以 $f(t)$ 确实是 $t$ 的次数 $\leq m$ 的多项式.
(e). 函数 $s \longmapsto Z_m(s)Z_{m'}(s)$ 带状, 所以
\[
\int_S Z_m(s)Z_{m'}(s)\,\dd\mu(s) = c\int_{-1}^1 f_m(t)f_{m'}(t)(1-t^2)^{(n-3)/2}\,\dd t
\]
由 (b). 由习题 \ref{ex:V.8.1.5}(b), 左端为 $0$.
(f). 对每个 $m \geq 0$, 函数 $Z_m$ 由定义非零 (且因为存在满足 $g(v_0) \neq 0$ 的 $f \in \mathcal{H}_m(\R^n)$, 见习题 \ref{ex:V.8.1.4}(h) 的证明), 所以由习题 \ref{ex:V.8.1.2}(h), $f_m \neq 0$. 由 (e), 函数 $(f_m)_{m \geq 0}$ 在 $L^2([-1,1], (1-t^2)^{(n-3)/2}\,\dd t)$ 中正交, 所以它们也线性无关. 固定 $m \geq 0$. 次数 $\leq m$ 的多项式空间 $P_m$ 维数为 $m + 1$ 且包含线性无关族 $(f_0, \ldots, f_m)$, 所以此族是 $P_m$ 的基. 但 $f_0, \ldots, f_{m-1}$ 的次数 $\leq m - 1$, 所以 $f_m$ 的次数必为 $m$.
(g). 由推论 \ref{cor:V.7.2}, 若 $\mathcal{Z}$ 是 $\SO(n)$ 上的球函数的集合, 则
\[
L^2(S) = \bigoplus_{\phi \in \mathcal{Z}}V_{\phi}
\]
且 $\phi$ 生成 $V_{\phi}^{\SO(n-1)}$. 所以用习题 \ref{ex:V.8.1.5}, 看到球函数正好是空间 $(\mathcal{H}_m(\R^n)|_S)^{\SO(n-1)}$ 的生成元中把 $v_0$ 送到 $1$ 的那些.
对每个 $m \geq 0$, 函数 $Z_m \in \mathcal{H}_m(\R^n)|_S$ 被 $\SO(n-3)$ 不变, 所以它生成 $\mathcal{H}_m(\R^n)|_S$ 中 $\SO(n-1)$-不变向量的空间, 且有倍数 $x \longmapsto \frac{1}{Z_m(v_0)}Z_m(x \cdot v_0)$ 是球函数. 由于球函数必把 $v_0$ 送到 $1$, 这个倍数就是球函数.
(h). 这由习题 \ref{ex:V.8.1.4}(k) 得出. 直接证明也很容易.
(i). 这与习题 \ref{ex:V.8.1.4}(h) 的证明中的构造完全相同 (取 $V = \mathcal{H}_m(\R^n)|_S$). 同样的证明表明 $\psi_m \in \mathcal{H}_m(\R^n)|_S$, $\psi_m$ 被 $\SO(n-1)$ 左不变且 $\psi_m(v_0) = f_m(v_0) = 1$.
(j). 与 (g) 的证明思路相同: 在每个 $\mathcal{H}_m(\R^n)|_S$ 中有一个球函数, 它是此空间中把 $v_0$ 送到 $1$ 的唯一 $\SO(n-1)$-不变元素. 由 (i), 函数 $\psi_m$ 满足所有要求的性质.
\end{solution}

\subsection{\texorpdfstring{Gelfand 对 $(S_n, S_r \times S_{n-r})$}{Gelfand 对 (S_n, S_r x S_{n-r})}}

这一系列习题的目标是研究 Gelfand 对 $(S_n, S_r \times S_{n-r})$, 摘录自 {[}7{]} 的第 5.1, 6.1 和 6.2 节. 我们按如下方式把 $S_r \times S_{n-r}$ 嵌入 $S_n$: 若 $\sigma \in S_r$, $\tau \in S_{n-r}$, 则 $\sigma \times \tau \in S_n$ 由
\[
(\sigma \times \tau)(i) = \begin{cases} \sigma(i) & \text{若 } 1 \leq i \leq r, \\ \tau(i - r) + r & \text{若 } r + 1 \leq i \leq n \end{cases}
\]
给出.
若 $E$ 是有限集, 记 $L(E)$ 为函数 $f \colon E \to \CC$ 的空间, 带 $L^2$ 内积 $\langle f, f' \rangle = \sum_{x \in E}f(x)f'(x)$.

\begin{exercise}[距离传递作用]
\label{ex:V.8.2.1}
在此题中, 固定传递作用 (左) 在集合 $E$ 上的有限群 $G$. 设 $x_0 \in E$, $K \subset G$ 是 $x_0$ 的稳定化子.
\begin{enumerate}
\item[(a)] 证明下列条件等价:
  \begin{enumerate}
  \item[(i)] 对所有 $x, y \in E$, 存在 $g \in G$ 使得 $g \cdot (x, y) = (y, x)$.
  \item[(ii)] 对每个 $g \in G$, 有 $g^{-1} \in KgK$.
  \end{enumerate}
\item[(b)] 若 (a) 的条件满足, 证明 $(G, K)$ 是 Gelfand 对.
\end{enumerate}
现在我们假设有度量 $d \colon E \times E \to \R_{\geq 0}$, 且群 $G$ 保距作用. 假设 $G$ 在 $E$ 上的作用是距离传递 (distance-transitive) 的, 即对所有满足 $d(x, y) = d(x', y')$ 的 $(x, y), (x', y') \in E \times E$, 存在 $g \in G$ 使得 $g \cdot (x, y) = (x', y')$.
\begin{enumerate}
\item[(c)] 证明 $(G, K)$ 是 Gelfand 对.
\item[(d)] 证明 $K$ 在 $E$ 上的轨道是球面 $\{x \in E \mid d(x, x_0) = j\}$, 其中 $j \in \R_{\geq 0}$.
\item[(e)] 设 $\Omega_r$ 是 $\{1, \ldots, n\}$ 的基数为 $r$ 的子集的集合. 证明公式 $d(A, B) = r - |A \cap B|$ 定义 $\Omega_r$ 上的度量.
\item[(f)] 证明 $(S_n, S_r \times S_{n-r})$ 是 Gelfand 对.
\end{enumerate}
\end{exercise}

\begin{solution}
(a). 假设 (i) 成立. 设 $g \in G$. 由 (i), 存在 $h \in G$ 使得 $h \cdot (K, gK) = (gK, K)$. 则第一个分量处的等式给出 $g^{-1}h \in K$, 第二个分量处的等式给出 $g^{-1}hg \in g^{-1}K$, 所以 $g \in (g^{-1}h)^{-1}g^{-1}K \subset Kg^{-1}K$.
假设 (ii) 成立. 设 $g, h \in G$. 由 (ii), 可以找到 $k_1, k_2 \in K$ 使得 $g^{-1}h = k_1h^{-1}gk_2$, 然后
\begin{align*}
(gK, hK) &= g \cdot (K, g^{-1}hK) \\
&= g \cdot (K, k_1h^{-1}gK) \\
&= gk_1^{-1} \cdot (K, h^{-1}gK) \\
&= gk_1^{-1}h^{-1} \cdot (hK, gK).
\end{align*}
(b). 这由命题 \ref{prop:V.2.5} 得出, 取 $\theta = \mathrm{id}_G$.
(c). 设 $x, y \in X$. 则 $d(y, x) = d(x, y)$, 所以存在 $g \in G$ 使得 $g \cdot (x, y) = (y, x)$. 换句话说, (a) 的条件 (i) 满足, 所以条件 (ii) 也满足. 由 (b), 这蕴含 $(G, K)$ 是 Gelfand 对.
(d). 对 $j \in \Z_{\geq 0}$ 写 $S_j = \{x \in X \mid d(x, x_0) = j\}$. 由于 $G$ 在 $X$ 上保距作用且 $K$ 固定 $x_0$, 集合 $S_j$ 被 $K$ 稳定. 为证明它们是 $K$ 在 $X$ 上的轨道, 需要证明 $K$ 在每个非空 $S_j$ 上传递作用. 所以设 $j \geq 0$, 假设有 $x, y \in S_j$. 则 $d(x_0, x) = d(x_0, y) = j$, 所以由假设, 存在 $g \in G$ 使得 $g \cdot (x_0, x) = (x_0, y)$. $g \cdot x_0 = x_0$ 这一事实蕴含 $g \in K$, 所以 $x$ 和 $y$ 在同一 $K$-轨道中.
(e). 显然对所有 $A, B \in \Omega_r$ 有 $d(A, B) = d(B, A)$. 设 $A, B, C \in \Omega_r$. 首先, 若 $d(A, B) = 0$, 则 $|A \cap B| = r = |A| = |B|$, 所以 $A \cap B = A$ 且 $A \cap B = B$, 从而 $A = B$. 证明三角不等式. 有
\[
|A \cap B| + |B \cap C| = |(A \cap B) \cup (B \cap C)| + |A \cap B \cap C| \leq |B| + |A \cap C| = r + |A \cap C|,
\]
所以
\[
d(A, C) = r - |A \cap C| \leq 2r - |A \cap B| - |B \cap C| = d(A, B) + d(B, C).
\]
(f). 让 $S_n$ 按 $\sigma \cdot A = \sigma(A)$ 作用于 $\Omega_r$. 这个作用传递, 且 $S_r \times S_{n-r}$ 是 $\{1, \ldots, r\}$ 的稳定化子. 又, 显然 $S_n$ 在 $\Omega_r$ 上保距作用. 所以由 (c), 只需验证作用是距离传递的. 设 $A, B, A', B' \in \Omega_r$ 满足 $d(A, B) = d(A', B')$, 即 $|A \cap B| = |A' \cap B'|$. 选取双射 $\phi \colon A \to A'$ 把 $A$ 的子集 $A \cap B$ 送到 $A' \cap B'$; 这由 $|A \cap B| = |A' \cap B'|$ 可能. 选取双射 $\psi \colon B - (A \cap B) \to B' - (A' \cap B')$; 这也可能, 因为 $|B - (A \cap B)| = r - |A \cap B| = |B' - (A' \cap B')|$. 把 $\phi$ 和 $\psi$ 放在一起给出双射 $A \cup B \xrightarrow{\sim} A' \cup B'$, 它把 $A$ 送到 $A'$, 把 $B$ 送到 $B'$, 且此双射的任何延拓 $\sigma \in S_n$ 将满足 $\sigma \cdot (A, B) = (A', B')$.
\end{solution}

\begin{exercise}[\index{距离正则图}距离正则图]
\label{ex:V.8.2.2}
设 $(\Omega_r, d)$ 是习题 \ref{ex:V.8.2.1}(e) 的有限度量空间. 设 $N$ 是 $\Omega_r$ 的\index{直径}直径, 即
\[
N = \max\{d(A, B) \mid A, B \in \Omega_r\}.
\]
对 $i \in \{0, \ldots, N\}$, 定义线性算子 $\Delta_i \colon L(\Omega_r) \to L(\Omega_r)$:
\[
\Delta_if(A) = \sum_{B \in \Omega_r, d(A,B)=i}f(B)
\]
对每个 $f \in L(\Omega_r)$. 又记 $A_f$ 为 $\End(L(\Omega_r))$ 中由 $\Delta_0, \ldots, \Delta_N$ 生成的子代数.
\begin{enumerate}
\item[(a)] 证明 $N = \min(r, n - r)$.
\item[(b)] 证明存在整数 $b_0, \ldots, b_N, c_0, c_1, \ldots, c_N$ 使得对每个 $i$ 和所有满足 $d(A, B) = i$ 的 $A, B \in \Omega_r$,
\[
|\{C \in \Omega_r \mid d(A, C) = 1 \text{ 且 } d(B, C) = i + 1\}| = b_i
\]
且
\[
|\{C \in \Omega_r \mid d(A, C) = 1 \text{ 且 } d(B, C) = i - 1\}| = c_i.
\]
(当然, $c_0 = 0$ 且 $c_1 = 1$.)\footnote{换句话说, 以 $\Omega_r$ 为顶点集、在满足 $d(A, B) = 1$ 的 $A, B \in \Omega_r$ 之间连边的图是距离正则的, 见 {[}7{]} 的定义 5.1.1.}
\item[(c)] 证明 $c_2, \ldots, c_N > 0$.
\item[(d)] 若 $i \in \{1, \ldots, N\}$, 证明
\[
\Delta_i\Delta_1 = b_{i-1}\Delta_{i-1} + (b_0 - b_i - c_i)\Delta_i + c_{i+1}\Delta_{i+1},
\]
约定 $\Delta_{N+1} = 0$.
\item[(e)] 证明存在多项式 $p_0, \ldots, p_N \in \R[t]$ 使得 $\deg(p_i) = i$ 且 $\Delta_i = p_i(\Delta_1)$.
\item[(f)] 证明 $A_f$ 是 $\End(L(\Omega_r))$ 中由 $\Delta_1$ 生成的子代数.
\item[(g)] 证明 $A_f$ 作为 $\CC$-向量空间由 $\Delta_0, \ldots, \Delta_N$ 张成.
\item[(h)] 证明 $\dim_{\CC}A_f = N + 1$.
\item[(i)] 证明 $L(\Omega_r)$ 的自同态 $\Delta_1$ 自伴.
\item[(j)] 证明有分解成两两正交子空间的直和 $L(\Omega_r) = \bigoplus_{i=0}^N V_i$, 其中 $V_0, \ldots, V_N$ 是 $\Delta_1$ 的特征空间. (提示: 习题 \ref{ex:II.5.2}.)
\end{enumerate}
\end{exercise}

\begin{solution}
(a). 注意对所有 $A, B \in \Omega_r$ 有 $d(A, B) \leq r$, 所以 $N \leq r$. 又, 对所有 $A, B \in \Omega_r$,
\[
d(A, B) = r - |A \cap B| = r - (|A| + |B| - |A \cup B|) = |A \cup B| - r \leq n - r,
\]
所以 $N \leq n - r$, 且 $N \leq \min(r, n - r)$.
现在取 $A = \{1, \ldots, r\}$, $B = \{n - r + 1, \ldots, n\}$. 则 $A, B \in \Omega_r$, 且 $|A \cap B| = \max(0, 2r - n)$. 所以 $N \geq d(A, B) = r - \max(0, 2r - n) = \min(r, n - r)$.
(b). 固定 $i \in \{0, \ldots, N\}$. 对所有满足 $d(A, B) = i$ 的 $A, B \in \Omega_r$, 设
\[
X_i(A, B) = \{C \in \Omega_r \mid d(A, C) = 1 \text{ 且 } d(B, C) = i + 1\}
\]
且
\[
Y_i(A, B) = \{C \in \Omega_r \mid d(A, C) = 1 \text{ 且 } d(B, C) = i - 1\}.
\]
若 $\sigma \in S_n$ 满足 $\sigma(A, B) = (A', B')$, 则 $\sigma$ 诱导双射 $X_i(A, B) \xrightarrow{\sim} X_i(A', B')$ 和 $Y_i(A, B) \xrightarrow{\sim} Y_i(A', B')$, 因为 $S_n$ 在 $\Omega_r$ 上保距作用. 所以结论由我们在习题 \ref{ex:V.8.2.1}(f) 的证明中证明的 $S_n$ 在 $\Omega_r$ 上的作用是距离传递的这一事实得出.
(c). 设 $i \in \{2, \ldots, N\}$. 取 $A = \{1, \ldots, r\}$, $B = \{i + 1, i + 2, \ldots, i + r\}$. (注意 $i + r \leq N + r \leq n$ 由 (a).) 需要证明存在至少一个 $C \in \Omega_r$ 使得 $d(A, C) = 1$ (即 $|A \cap C| = r - 1$) 且 $d(B, C) = i - 1$ (即 $|B \cap C| = r - i + 1$). 这对 $C = \{2, 3, \ldots, r + 1\}$ 成立.
(d). 设 $f \in L(\Omega_r)$, $A \in \Omega_r$. 则
\[
\Delta_i\Delta_1f(A) = \sum_{B \in \Omega_r, d(A,B)=i}\ \sum_{C \in \Omega_r, d(B,C)=1}f(C).
\]
设 $C \in \Omega_r$. 若存在 $B \in \Omega_r$ 使得 $d(A, B) = i$ 且 $d(B, C) = 1$, 则由三角不等式必有 $i - 1 \leq d(A, C) \leq i + 1$.
假设 $d(A, C) = i + 1$. 则
\[
\{B \in \Omega_r \mid d(A, B) = i \text{ 且 } d(B, C) = 1\} = Y_{i+1}(C, A)
\]
(用 (b) 的证明中的记号). 假设 $d(A, C) = i - 1$. 则
\[
\{B \in \Omega_r \mid d(A, B) = i \text{ 且 } d(B, C) = 1\} = X_{i-1}(C, A).
\]
最后, 假设 $d(A, C) = i$. 考虑集合
\begin{multline*}
\{B \in \Omega_r \mid d(A, B) = i \text{ 且 } d(B, C) = 1\} \cup \{B \in \Omega_r \mid d(A, B) = i + 1 \text{ 且 } d(B, C) = 1\} \\
\cup \{B \in \Omega_r \mid d(A, B) = i - 1 \text{ 且 } d(B, C) = 1\}.
\end{multline*}
这个并显然不相交. 要计算第一个集合的基数, 第二个集合是 $X_i(C, A)$, 第三个集合是 $Y_i(C, A)$. 又, 由三角不等式, 这个并正好是
\[
\{B \in \Omega_r \mid d(B, C) = 1\} = X_0(C, C).
\]
所以得到
\[
|\{B \in \Omega_r \mid d(A, B) = i \text{ 且 } d(B, C) = 1\}| + b_i + c_i = b_0.
\]
最后看到
\begin{align*}
\Delta_i\Delta_1f(A) &= c_{i+1}\sum_{C, d(A,C)=i+1}f(C) + b_{i-1}\sum_{C, d(A,C)=i-1}f(C) + (b_0 - b_i - c_i)\sum_{C, d(A,C)=i}f(C) \\
&= c_{i+1}\Delta_{i+1}f(A) + b_{i-1}\Delta_{i-1}f(A) + (b_0 - b_i - c_i)\Delta_if(A).
\end{align*}
(e). 对 $i$ 归纳证明此结论. $i = 0$ 时显然 (注意 $\Delta_0 = \mathrm{id}$, 所以取 $p_0 = 1$), $i = 1$ 时也显然 (取 $p_1(t) = t$). 假设对某个 $i \geq 1$ 已知结果, 证明对 $i + 1$. 由 (c) 和 (d),
\[
\Delta_{i+1} = c_{i+1}^{-1}(\Delta_i\Delta_1 - b_{i-1}\Delta_{i-1} - (b_0 - b_i - c_i)\Delta_i),
\]
所以 $\Delta_{i+1} = p_{i+1}(\Delta_1)$, 其中
\[
p_{i+1}(t) = c_{i+1}^{-1}(tp_i(t) - p_{i-1}(t) - (b_0 - b_i - c_i)p_i(t)).
\]
显然也有 $\deg(p_{i+1}(t)) = i + 1$.
(f). 设 $A'_f$ 是 $\End(L(\Omega_r))$ 中由 $\Delta_1$ 生成的子代数. 则 $A'_f \subset A_f$ 由 $A_f$ 的定义. 由 (e), 有 $\Delta_0, \ldots, \Delta_N \in A'_f$, 所以 $A_f \subset A'_f$.
(g). 对 $i \geq 0$ 归纳证明 $\Delta_1^i \in \operatorname{Span}(\Delta_0, \Delta_1, \ldots, \Delta_i)$. (结论将由 (f) 得出.) $i = 0$ 和 $i = 1$ 时显然. 假设对 $i \geq 1$ 成立, 证明对 $i + 1$. 由 (e), 存在非零 $a \in \R$ 和 $a_0, \ldots, a_i \in \R$ 使得 $\Delta_{i+1} = a\Delta_1^{i+1} + \sum_{j=0}^i a_j\Delta_1^j$. 由于由归纳假设对每个 $j \in \{0, \ldots, i\}$ 有 $\Delta_1^j \in \operatorname{Span}(\Delta_0, \ldots, \Delta_j)$, 推出 $\Delta_1^{i+1} \in \operatorname{Span}(\Delta_0, \ldots, \Delta_{i+1})$.
(h). 知道 $A_f = \operatorname{Span}(\Delta_0, \ldots, \Delta_N)$, 所以只需证明族 $(\Delta_0, \ldots, \Delta_N)$ 线性无关. 设 $c_0, \ldots, c_N \in \CC$. 若 $A, B \in \Omega_r$, 记 $\delta_A$ 为 $\{A\}$ 的特征函数, 则 $d(A, B) = i$ 时 $\Delta_i\delta_A(B) \neq 0$, 且
\[
\sum_{i=0}^N c_i\Delta_i\delta_A(B) = c_{d(A,B)}.
\]
由于存在满足 $d(A, B) = i$ (对每个 $i \in \{0, \ldots, N\}$) 的偶对 $(A, B) \in \Omega_r^2$, 若 $\sum_{i=0}^N c_i\Delta_i = 0$ 则 $c_0 = \cdots = c_N = 0$.
(i). 设 $f, g \in L(\Omega_r)$. 则
\begin{align*}
\langle f, \Delta_1g \rangle &= \sum_{A \in \Omega_r}f(A)(\Delta_1g)(A) \\
&= \sum_{A \in \Omega_r}\ \sum_{B \in \Omega_r, d(A,B)=1}f(A)g(B) \\
&= \sum_{A,B \in \Omega_r, d(A,B)=1}f(A)g(B) \\
&= \sum_{B \in \Omega_r}(\Delta_1f)(B)g(B) \\
&= \langle \Delta_1f, g \rangle.
\end{align*}
(j). 由于 $\Delta_1$ 自伴, 谱定理说它可对角化且其特征空间两两正交. 所以只需证明有 $N + 1$ 个特征空间. 由习题 \ref{ex:II.5.2}(b), 知道 $\End(L(\Omega_r))$ 中由 $\Delta_1$ 生成的子代数, 即 $A_f$ (见 (f)), 约化. 由习题 \ref{ex:II.5.2}(c), 知道 $A_f$ (即 $\Delta_1$) 的特征空间数是 $\dim(A_f)$, 由 (h) 知道 $\dim(A_f) = N + 1$.
\end{solution}

\begin{exercise}[对称幂的分解]
\label{ex:V.8.2.3}
用习题 \ref{ex:V.8.2.2} 的记号. 注意有 $G := S_n$ 在 $\Omega_r$ 上的作用, 且 $\{1, \ldots, r\}$ 的稳定化子是 $K := S_r \times S_{n-r}$. 设 $M^{n-r,r} = L(\Omega_r)$, 看作 $S_n$ 的表示, 通过拟正则表示 (即若 $g \in G$, $f \in M^{n-r,r}$, $A \in \Omega_r$, 有 $(g \cdot f)(A) = f(g^{-1}A)$).
定义 $d \colon M^{n-r,r} \to M^{n-r+1,r-1}$ 和 $d^* \colon M^{n-r+1,r-1} \to M^{n-r,r}$:
\[
(df)(A) = \sum_{B \in \Omega_r, A \subset B}f(B)
\]
且
\[
(d^*f)(B) = \sum_{A \in \Omega_r, A \subset B}f(A).
\]
(若 $r = 0$, 取 $d = 0$ 和 $d^* = 0$.)
又记 $\Delta$ 为习题 \ref{ex:V.8.2.2} 的算子 $\Delta_1$; 即对每个 $f \in M^{n-r,r}$, 函数 $\Delta f \in M^{n-r,r}$ 由
\[
(\Delta f)(A) = \sum_{B \in \Omega_r, d(A,B)=1}f(B)
\]
定义.\footnote{注意函数 $d$, $d^*$ 和 $\Delta$ 对每个 $r$ 都有定义; 记号中不标出 $r$, 它应该由上下文清楚. 这与\index{离散 Laplace 算子}离散 Laplace 算子密切相关. 事实上, $M^{n-r,r}$ 上的离散 Laplace 算子最常见的定义是 $\frac{1}{r}\Delta - \mathrm{id}$.}
最后, 若 $a \in \CC$, $i \in \Z_{\geq 0}$, 记
\[
(a)_i = a(a+1)\cdots(a+i-1).
\]
例如有 $(1)_n = n!$ 和 $\binom{n}{k} = \frac{(n-k+1)_k}{k!}$.
\begin{enumerate}
\item[(a)] 证明 $A \longmapsto \{1, \ldots, n\} - A$ 诱导 $G$-等变同构 $M^{r,n-r} \xrightarrow{\sim} M^{n-r,r}$.
\item[(b)] 证明 $d$ 和 $\Delta$ 是 $G$-等变的.
\item[(c)] 证明 $d^*$ 是 $d$ 的伴随.
\item[(d)] 若 $f \in M^{n-r,r}$, 证明
\[
dd^*f = \Delta f + (n-r)f \quad \text{且} \quad d^*df = \Delta f + rf.
\]
\item[(e)] 设 $f \in M^{n-r,r}$, $1 \leq p \leq q \leq n - r$. 证明
\[
d(d^*)^qf = (d^*)^qdf + q(n - 2r - q + 1)(d^*)^{q-1}f.
\]
若此外 $df = 0$, 证明
\[
d^p(d^*)^qf = (q - p + 1)_p(n - 2r - q + 1)_p(d^*)^{q-p}f.
\]
\end{enumerate}
假设 $0 \leq r \leq n/2$. 若 $r > 0$, 设 $S^{n-r,r} = \Ker(d \colon M^{n-r,r} \to M^{n-r+1,r-1})$; 若 $r = 0$, 设 $S^{n-r,r} = M^{n-r,r}$. 这是 $M^{n-r,r}$ 的 $G$-稳定子空间.
\begin{enumerate}
\item[(f)] 若 $0 \leq m \leq n$ 且 $0 \leq r \leq \min(m, n - m)$, 证明 $(d^*)^{m-r} \colon S^{n-r,r} \to M^{n-m,m}$ 单射. (提示: 计算 $\|(d^*)^{m-r}f\|_2^2$.)
\item[(g)] 在 (f) 的假设下, 证明 $(d^*)^{m-r}(S^{n-r,r})$ 包含于 $\Delta$ 的特征值 $m(n-m) - r(n-r+1)$ 的特征空间中.
\item[(h)] 证明若 $1 \leq m \leq n/2$, 则 $M^{n-m,m}$ 中 $S^{n-m,m}$ 的正交补是 $d^*(M^{n-m+1,m-1})$.
\item[(i)] 证明对满足 $0 \leq r \leq n/2$ 的每个 $r$, $S^{n-r,r} \neq 0$.
\item[(j)] 若 $0 \leq m \leq n$, 证明
\[
M^{n-m,m} = \bigoplus_{r=0}^{\min(m,n-m)}(d^*)^{m-r}(S^{n-r,r}),
\]
其中和项两两正交且正好是 $\Delta$ 的特征空间.
\item[(k)] 证明若 $r > 0$, $\dim_{\CC}(S^{n-r,r}) = \binom{n}{r} - \binom{n}{r-1}$.
\item[(l)] 证明表示 $S^{n-r,r}$ ($0 \leq r \leq n/2$) 不可约且两两不等价. (提示: $M^{m,n-m}$ 有多少个不可约成分?)
\end{enumerate}
\end{exercise}

\begin{solution}
(a). 映射 $A \longmapsto \{1, \ldots, n\} - A$ 是从 $\Omega_r$ 到 $\Omega_{n-r}$ 的双射, 且与 $G$ 在这两个集合上的作用交换. 结论立即得出.
(b). 设 $f \in M^{n-r,r}$, $\sigma \in S_n$. 若 $A \in \Omega_{r-1}$, 则
\begin{align*}
(dL_\sigma f)(A) &= \sum_{B \in \Omega_r, B \supset A}f(\sigma^{-1}(B)) \\
&= \sum_{B' \in \Omega_r, B' \supset \sigma^{-1}(A)}f(B') \\
&= (df)(\sigma^{-1}(A)) \\
&= L_\sigma(df)(A).
\end{align*}
若 $A \in \Omega_r$, 则
\begin{align*}
(\Delta L_\sigma f)(A) &= \sum_{B \in \Omega_r, d(A,B)=1}f(\sigma^{-1}(B)) \\
&= \sum_{B' \in \Omega_r, d(\sigma^{-1}(A),B')=1}f(B') \\
&= \Delta f(\sigma^{-1}(A)) \\
&= L_\sigma(\Delta f)(A).
\end{align*}
(c). 设 $f \in M^{n-r,r}$, $g \in M^{n-r+1,r-1}$. 则
\begin{align*}
\langle df, g \rangle &= \sum_{A \in \Omega_{r-1}}df(A)g(A) \\
&= \sum_{A \in \Omega_{r-1}, B \in \Omega_r, A \subset B}f(B)g(A) \\
&= \sum_{B \in \Omega_r}f(B)d^*g(B).
\end{align*}
(d). 设 $f \in M^{n-r,r}$, $A \in \Omega_r$. 则
\begin{align*}
dd^*f(A) &= \sum_{B \in \Omega_{r+1}, B \supset A}d^*f(B) \\
&= \sum_{B \in \Omega_{r+1}, C \in \Omega_r, C \subset B \supset A}f(C).
\end{align*}
设 $C \in \Omega_r$. 若存在 $B \in \Omega_{r+1}$ 使得 $C \subset B \supset A$, 则 $d(A, C) \leq 1$. 若 $C = A$, 则
\[
|\{B \in \Omega_{r+1} \mid C \subset B \supset A\}| = |\{B \in \Omega_{r+1} \mid A \subset B\}| = n - r.
\]
若 $d(A, C) = 1$, 则同时包含 $A$ 和 $C$ 的 $\Omega_{r+1}$ 的唯一元素是 $A \cup C$. 最后得到
\begin{align*}
dd^*f(A) &= (n - r)f(A) + \sum_{C \in \Omega_r, d(A,C)=1}f(C) \\
&= (n - r)f(A) + \Delta f(A).
\end{align*}
类似地,
\begin{align*}
d^*df(A) &= \sum_{B \in \Omega_{r-1}, B \subset A}df(B) \\
&= \sum_{C \in \Omega_r, B \in \Omega_{r-1}, C \supset B \subset A}f(C).
\end{align*}
设 $C \in \Omega_r$. 若存在 $B \in \Omega_{r-1}$ 使得 $C \supset B \subset A$, 则 $d(A, C) \leq 1$. 若 $A = C$, 则
\[
|\{B \in \Omega_{r-1} \mid C \supset B \subset A\}| = r.
\]
若 $d(A, C) = 1$, 则同时包含于 $A$ 和 $C$ 的唯一 $B \in \Omega_{r-1}$ 是 $B = A \cap C$. 所以得到
\begin{align*}
d^*df(A) &= rf(A) + \sum_{C \in \Omega_r, d(A,C)=1}f(C) \\
&= rf(A) + \Delta f(A).
\end{align*}
(e). 对 $q$ 归纳证明第一个恒等式. 若 $q = 1$, 由 (d),
\[
dd^*f = \Delta f + (n-r)f = \Delta f + rf + (n-2r)f = d^*df + q(n-2r-q+1)(d^*)^{q-1}f
\]
对每个 $f \in M^{n-r,r}$. 现在假设对 $q \in \{1, \ldots, n-r-1\}$, 每个 $s$ 和 $M^{n-s,s}$ 的每个元素已知恒等式, 证明对 $q+1$. 若 $f \in M^{n-r,r}$, 有
\begin{align*}
d(d^*)^{q+1}f &= d(d^*)^q(d^*f) \\
&= (d^*)^qd(d^*f) + q(n-2(r+1)-q+1)(d^*)^{q-1}(d^*f) \quad \text{(对 } d^*f \in M^{n-r-1,r+1} \text{ 用归纳假设)} \\
&= (d^*)^q(d^*df + (n-2r)f) + q(n-2r-q-1)(d^*)^qf \quad \text{(用 } q=1 \text{ 的情形)} \\
&= (d^*)^{q+1}df + (n-2r + q(n-2r-q-1))(d^*)^qf \\
&= (d^*)^{q+1}df + (q+1)(n-2r-(q+1)+1)(d^*)^qf.
\end{align*}
现在对 $p$ 归纳证明第二个恒等式. 若 $p = 1$, 它归结为第一个恒等式 (用 $df = 0$). 假设对某个 $p \in \{1, \ldots, q-1\}$ (和所有 $s$, 所有满足 $df = 0$ 的 $f \in M^{n-s,s}$) 已证明, 证明对 $p+1$. 设 $f \in M^{n-r,r}$ 满足 $df = 0$. 则
\begin{align*}
d^{p+1}(d^*)^qf &= d(d^p(d^*)^qf) \\
&= d((q-p+1)_p(n-2r-q+1)_p(d^*)^{q-p}f) \\
&= (q-p+1)_p(n-2r-q+1)_p(q-p)(n-2r-(q-p)+1)(d^*)^{q-p-1}f \\
&\qquad \text{(用第一个恒等式和 } df = 0\text{)} \\
&= (q-p)_{p+1}(n-2r-q+1)_{p+1}(d^*)^{q-p-1}f.
\end{align*}
(f). 设 $f \in S^{n-r,r}$, $f \neq 0$. 用 (c) 然后用 (e) 的第二个恒等式, 看到
\begin{align*}
\langle (d^*)^{m-r}f, (d^*)^{m-r}f \rangle &= \langle d^{m-r}(d^*)^{m-r}f, f \rangle \\
&= (1)_{m-r}(n-2r-(m-r)+1)_{m-r}\langle f, f \rangle \\
&\neq 0,
\end{align*}
所以 $(d^*)^{m-r}f \neq 0$.
(g). 设 $f \in S^{n-r,r}$. 用 (d) 的第二个公式在 $M^{n-m,m}$ 上计算 $\Delta$, 用 (e) 的第二个公式 (取 $p = 1$), 得到
\begin{align*}
\Delta((d^*)^{m-r}f) &= d(d^*)^{m-r+1}f - (n-m)(d^*)^{m-r}f \\
&= (m-r+1)(n-2r-(m-r+1)+1)(d^*)^{m-r}f - (n-m)(d^*)^{m-r}f \\
&= (m(n-m) - r(n-r+1))(d^*)^{m-r}f.
\end{align*}
(h). 这是 $S^{n-m,m}$ 的定义和 (c) 的直接推论.
(i). 空间 $S^{n-r,r}$ 是 $d \colon M^{n-r,r} \to M^{n-r+1,r-1}$ 的核, 且 $\dim(M^{n-r+1,r-1}) = \binom{n}{r-1} < \binom{n}{r} = \dim(M^{n-r,r})$ (因为 $r \leq n/2$), 所以 $d$ 不能单射.
(j). 子空间 $(d^*)^{m-r}(S^{n-r,r})$ ($0 \leq r \leq \min(m, n-m)$) 由 (g) 包含于 $\Delta$ 的不同特征空间. 它们都非零由 (f) 和 (i). 知道 $\Delta$ 自伴 (由习题 \ref{ex:V.8.2.2}(i)), 所以这些空间两两正交. 又知道 $\Delta \in \End(M^{n-m,m})$ 恰有 $1 + \min(m, n-m)$ 个特征值 (由习题 \ref{ex:V.8.2.2}(a) 和 \ref{ex:V.8.2.2}(j)), 所以这些特征值必是 $m(n-m) - r(n-r+1)$, $0 \leq r \leq \min(m, n-m)$. 还需证明
\[
M^{n-m,m} = \bigoplus_{r=0}^{\min(m,n-m)}(d^*)^{m-r}(S^{n-r,r}).
\]
对 $m$ 归纳证明. $m = 0$ 时显然. 假设对 $m - 1$ (其中 $n/2 \geq m \geq 1$) 已知结果, 证明对 $m$. 由 (h),
\[
M^{n-m,m} = S^{n-m,m} \oplus d^*(M^{n-m+1,m-1}).
\]
由归纳假设,
\[
M^{n-m+1,m-1} = \bigoplus_{r=0}^{m-1}(d^*)^{m-1-r}(S^{n-r,r}).
\]
$m$ 的结果由这两个事实立即得出.
最后处理 $m \geq n/2$ 的情形. 设 $m' = n - m$. 已经看到
\[
M^{n-m',m'} = \bigoplus_{r=0}^{m'}(d^*)^{m'-r}(S^{n-r,r}).
\]
由 (k), 这蕴含 $\dim(M^{n-m',m'}) = \sum_{r=0}^{m'}\dim(S^{n-r,r})$. 又看到
\[
M^{n-m,m} \supset \bigoplus_{r=0}^{m'}(d^*)^{m-r}(S^{n-r,r}),
\]
且再次由 (k), 这蕴含 $\dim(M^{n-m,m}) \geq \sum_{r=0}^{m'}\dim(S^{n-r,r}) = \dim(M^{m',n-m'})$. 但 $\dim(M^{n-m,m}) = \dim(M^{n-m',m'})$ (由 (a)), 所以上面的不等式是等式, 且
\[
M^{n-m,m} = \bigoplus_{r=0}^{m'}(d^*)^{m-r}(S^{n-r,r}).
\]
(k). 由 (j) 和 (f), 映射 $d^* \colon M^{n-r+1,r-1} \to M^{n-r,r}$ 单射. 由 (h), 这蕴含
\[
\dim(S^{n-r,r}) = \dim(M^{n-r,r}) - \dim(M^{n-r+1,r-1}) = \binom{n}{r} - \binom{n}{r-1}.
\]
(l). 设 $m = \lfloor n/2 \rfloor$. 由于映射 $d$ 和 $d^*$ 是 $S_n$-等变的 (对 $d$ 见 (b), $d^*$ 是 $d$ 的伴随, 所以也等变), 子空间 $S^{n-r,r} \subset M^{n-r,r}$ 对每个 $r \leq n/2$ 是 $S_n$-稳定的, (j) 的分解是 $S_n$-子空间的分解. 其次, 由习题 \ref{ex:V.8.2.1}(f) 知道 $(S_n, S_m \times S_{n-m})$ 是 Gelfand 对, 所以由定理 \ref{thm:V.3.2.4}, 对应的拟正则表示, 即 $M^{n-m,r}$, 分解成互不相同的不可约表示的直和. 由推论 \ref{cor:V.7.2}, $M^{n-m,m}$ 中不可约和项的数是 Gelfand 对的球函数数, 它等于 $S_n$ 上双不变函数的空间维数 (因为球函数构成这些双不变函数的基, 由同一推论的 (iii)), 即 $(S_m \times S_{n-m}) \backslash S_n /(S_m \times S_{n-m})$ 的基数, 这也等于 $S_m \times S_{n-m}$ 在 $\Omega_m$ 上的轨道数. 但由习题 \ref{ex:V.8.2.1}(d) 看到, $S_m \times S_{n-m}$ 在 $\Omega_m$ 上的轨道是以 $A_0 := \{1, \ldots, m\}$ 为中心的球面. 这些球面的可能半径是 $0, 1, \ldots, \min(m, n-m) = m$ (由习题 \ref{ex:V.8.2.2}(a)), 且容易看到所有球面非空 (在习题 \ref{ex:V.8.2.2}(f) 的证明中已经用过). 最后得到 $M^{n-m,m}$ 的不可约成分数是 $m + 1$. 由于 (j) 的分解是 $M^{n-m,m}$ 分解成 $m + 1$ 个非零子表示, 它必是其不可约成分的分解, 所以得到结论. (注意 $(d^*)^{m-r}(S^{n-r,r})$ 作为 $S_n$ 的表示与 $S^{n-r,r}$ 等价, 由 (f).)
\end{solution}

\begin{exercise}[置换群 Gelfand 对的球函数]
\label{ex:V.8.2.4}
我们保留习题 \ref{ex:V.8.2.3} 的记号. 若 $m, h \in \{0, \ldots, n\}$, $A \in \Omega_m$, $\max(0, h-m) \leq \ell \leq \min(n-m, h)$, 记 $\sigma_{\ell,h-\ell}(A) \in L(\Omega_h)$ 为集合 $\{C \in \Omega_h \mid |A \cap C| = h - \ell\}$ 的特征函数. 又写 $\sigma_{-1,h+1}(A) = 0$. 固定 $A \in \Omega_m$.
\begin{enumerate}
\item[(a)] 若 $h = m$, 证明对每个 $\ell \in \{0, \ldots, \min(m, n-m)\}$, 函数 $\sigma_{\ell,m-\ell}(A)$ 是球面 $\{C \in \Omega_m \mid d(A, C) = \ell\}$ 的特征函数.
\item[(b)] 证明
\[
d(\sigma_{\ell,h-\ell}(A)) = (n - m - \ell + 1)\sigma_{\ell-1,h-\ell}(A) + (m - h + \ell + 1)\sigma_{\ell,h-\ell-1}(A).
\]
\item[(c)] 若 $k \leq h$ 且 $\max(0, k-m) \leq i \leq \min(k, n-m)$, 证明
\[
\frac{1}{(h-k)!}(d^*)^{h-k}\sigma_{i,k-i}(A) = \sum_{\ell = \max(i,h-m)}^{\min(h-k+i,n-m)}\binom{\ell}{i}\binom{h-\ell}{k-i}\sigma_{\ell,h-\ell}(A).
\]
\end{enumerate}
从现在起, 取 $A = \{1, \ldots, m\}$.
\begin{enumerate}
\item[(a)] 若 $0 \leq h \leq \min(m, n-m)$, 证明 $M^{n-h,h}$ 中 $S_m \times S_{n-m}$-不变向量的空间由函数 $\sigma_{\ell,h-\ell}(A)$, $0 \leq \ell \leq h$, 张成.
\item[(b)] 若 $0 \leq h \leq \min(m, n-m)$, 证明 $S^{n-h,h}$ 中 $S_m \times S_{n-m}$-不变向量的空间由函数
\[
\sum_{\ell=0}^h\frac{(n - m - h + 1)_{h-\ell}}{(-m)_{h-\ell}}\sigma_{\ell,h-\ell}(A)
\]
张成.
\item[(c)] 对 $0 \leq h \leq \min(m, n-m)$, 设 $\phi_h \in M^{n-m,m}$ 是 (d) 的和项 $(d^*)^{m-h}(S^{n-h,h})$ 中包含的唯一球函数. 证明
\[
\phi_h = \sum_{\ell=0}^{\min(m,n-m)}\phi(n, m, h; \ell)\sigma_{\ell,m-\ell}(A),
\]
其中
\[
\phi(n, m, h; \ell) = (-1)^h\binom{n-m}{h}^{-1}\sum_{i=\max(0,\ell-m+h)}^{\min(\ell,h)}\binom{m-\ell}{h-i}\binom{\ell}{i}\frac{(n-m-h+1)_{h-i}}{(-m)_{h-i}}.
\]
\item[(d)] 固定 $h$ 使得 $0 \leq h \leq \min(m, n-m)$. 证明 $\phi_h$ 中 $\sigma_{1,m-1}(A)$ 的系数是 $1 - \frac{h(n-h+1)}{m(n-m)}$. (注记: 有一种用最少计算解决此题的方法.)
\end{enumerate}
\end{exercise}

\begin{solution}
(a). 函数 $\sigma_{\ell,m-\ell}(A)$ 是集合 $\{C \in \Omega_m \mid |A \cap C| = m - \ell\}$ 的特征函数. 由于 $d(A, C) = m - |A \cap C|$, 这个集合正好是以 $A$ 为心、$\ell$ 为半径的球面.
(b). 设 $B \in \Omega_{h-1}$. 有
\[
d(\sigma_{\ell,h-\ell}(A))(B) = \sum_{C \in \Omega_h, C \supset B}\sigma_{\ell,h-\ell}(A)(C) = |\{C \in \Omega_h \mid C \supset B \text{ 且 } |A \cap C| = h - \ell\}|.
\]
若存在至少一个 $C \in \Omega_h$ 使得 $C \supset B$ 且 $|A \cap C| = h - \ell$, 则 $A \cap C \supset A \cap B$, 且这两个集合至多相差一个元素, 所以 $|A \cap B| \in \{h - \ell, h - \ell - 1\}$.
假设 $|A \cap B| = h - \ell$. 则对每个满足 $C \supset B$ 且 $|A \cap C| = h - \ell$ 的 $C \in \Omega_h$, 必有 $A \cap C = A \cap B$. 通过向 $B$ 添加 $\{1, \ldots, n\} - (A \cup B)$ 的一个元素得到每个这样的 $C$, 所以 $C$ 的可能数是
\[
n - |A \cup B| = n - (|A| + |B| - |A \cap B|) = n - (m + h - 1 - (h - \ell)) = n - m - \ell + 1.
\]
假设 $|A \cap B| = h - \ell + 1$. 则对每个满足 $C \supset B$ 且 $|A \cap C| = h - \ell$ 的 $C \in \Omega_h$, $C - B$ 的唯一元素必须是 $A \cap C - A \cap B$ 的元素. 所以 $C$ 的可能数是
\[
|A - A \cap B| = m - h + \ell - 1.
\]
最后得到
\[
d(\sigma_{\ell,h-\ell}(A))(B) = (n - m - \ell + 1)\sigma_{\ell-1,h-\ell}(B) + (m - h + \ell + 1)\sigma_{\ell,h-\ell-1}(B),
\]
如所愿.
(c). 对每个 $i$ 和每个 $D \in \Omega_i$, 记 $\delta_D \in L(\Omega_i)$ 为 $\{D\}$ 的特征函数. 设 $S$ 是集合 $\{C \in \Omega_k \mid |A \cap C| = k - i\}$. 则 $\sigma_{i,k-i}(A)$ 是 $S$ 的特征函数. 若 $C \in S$, 则对每个 $D \in \Omega_h$,
\[
(d^*)^{h-k}\delta_C(D) = \sum_{C \subset D_1 \subset \cdots \subset D_{h-k-1} \subset D_{h-k} = D, D_i \in \Omega_{k+i}}1.
\]
若 $C \not\subset D$, 集合 $\{C \subset D_1 \subset \cdots \subset D_{h-k-1} \subset D_{h-k} = D \mid D_i \in \Omega_{k+i}\}$ 空; 若 $C \subset D$, 这个集合有 $(h-k)!$ 个元素. 所以看到
\[
(d^*)^{h-k}\delta_C = (h-k)!\sum_{C \subset D \in \Omega_h}\delta_D.
\]
所以若 $D \in \Omega_h$, $\frac{1}{(h-k)!}(d^*)^{h-k}\sigma_{i,k-i}(A)$ 中 $\delta_D$ 的系数是集合 $\{C \in S \mid C \subset D\}$ 的基数. 写 $|A \cap D| = h - \ell$, 其中 $0 \leq \ell \leq h$; 注意有
\[
h - \ell = |A \cap D| = |A| + |D| - |A \cup D| \geq m + h - |A \cup D|
\]
且 $h \leq |A \cup D| \leq n$, 所以 $h - m \leq \ell \leq n - m$, 且此范围内的所有非负 $\ell$ 都能出现. 通过从 $A \cap D$ 中移除 $(h - \ell) - (k - i)$ 个元素、从 $D - (A \cap D)$ 中移除 $\ell - i$ 个元素, 得到满足 $C \subset D$ 的 $C \in S$. 这只有在 $h - \ell - (k - i) \geq 0$ 且 $0 \leq \ell - i$ 时才可能, 且有 $\binom{h-\ell}{(h-\ell)-(k-i)}\binom{\ell}{\ell-i} = \binom{h-\ell}{k-i}\binom{\ell}{i}$ 种不同的可能选择. 把所有放在一起, 看到若 $|A \cap D| = h - \ell$, 则 $\frac{1}{(h-k)!}(d^*)^{h-k}\sigma_{i,k-i}(A)$ 中 $\delta_D$ 的系数在 $\max(h-m, i) \leq \ell \leq \min(n-m, h-k+i)$ 时为 $\binom{h-\ell}{k-i}\binom{\ell}{i}$, 否则为 $0$. 这给出结果.
(d). 此结论等价于: $S_m \times S_{n-m}$ 在 $\Omega_h$ 中的轨道是 $S_{\ell} := \{C \in \Omega_h \mid |A \cap C| = h - \ell\}$, $0 \leq \ell \leq h$. 证明这个事实. 由于 $S_m \times S_{n-m}$ 固定 $A$, 集合 $S_{\ell}$ 被 $S_m \times S_{n-m}$ 不变, 所以只需证明 $S_m \times S_{n-m}$ 在这些集合上传递作用. 固定 $\ell \in \{0, \ldots, h\}$, 取 $C, C' \in S_{\ell}$. 由于 $|A \cap C| = |A \cap C'|$, 可以找到 $\sigma \in S_m$ 把 $A \cap C$ 送到 $A \cap C'$. 又, $|\{m+1, \ldots, n\} \cap C| = |\{m+1, \ldots, n\} \cap C'| = \ell$, 所以可以找到 $\tau \in S_{n-m}$ 把 $\{m+1, \ldots, n\} \cap C$ 送到 $\{m+1, \ldots, n\} \cap C'$. 则 $\sigma \times \tau \in S_m \times S_{n-m}$ 把 $C$ 送到 $C'$.
(e). 设 $f \in S^{n-h,h}$ 是 $S_m \times S_{n-m}$-不变向量. 由 (d), 不变性条件等价于可以把 $f$ 写成 $f = \sum_{\ell=0}^h a_{\ell}\sigma_{\ell,h-\ell}(A)$, 其中 $a_0, \ldots, a_h \in \CC$. $f \in S^{n-h,h}$ 这一事实意味着 $df = 0$. 用 (b), 可以把此条件改写成
\begin{align*}
0 &= \sum_{\ell=0}^h a_{\ell}((n - m - \ell + 1)\sigma_{\ell-1,h-\ell}(A) + (m - h + \ell + 1)\sigma_{\ell,h-1-\ell}(A)) \\
&= \sum_{\ell=0}^{h-1}a_{\ell+1}(n - m - \ell)\sigma_{\ell,h-\ell+1}(A) + \sum_{\ell=0}^h a_{\ell}(m - h + \ell + 1)\sigma_{\ell,h-1-\ell}(A).
\end{align*}
由于函数 $\sigma_{\ell,h-1-\ell}(A)$, $0 \leq \ell \leq h - 1$, 线性无关 (因为它们有不相交的支撑), 此等式等价于对每个 $\ell \in \{0, \ldots, h-1\}$,
\[
a_{\ell} = a_{\ell+1}\frac{n-m-\ell}{-m+h-\ell-1}.
\]
对 $\ell$ 做简单的下降归纳表明这等价于对每个 $\ell \in \{0, \ldots, h\}$,
\[
a_{\ell} = \frac{(n - m - h + 1)_{h-\ell}}{(-m)_{h-\ell}}a_h.
\]
这蕴含所要的结果.
(f). 由 (e) (和 $d^*$ 的 $S_n$-等变性), 函数 $\phi_h$ 是
\[
\psi := \frac{1}{(m-h)!}(d^*)^{m-h}\Bigl(\sum_{i=0}^h\frac{(n - m - h + 1)_{h-i}}{(-m)_{h-i}}\sigma_{i,h-i}(A)\Bigr)
\]
的倍数. 用 (c) 的公式计算 $\psi$. 得到
\begin{align*}
\psi &= \sum_{i=0}^h\frac{(n - m - h + 1)_{h-i}}{(-m)_{h-i}}\sum_{\ell=i}^{\min(m-h+i,n-m)}\binom{\ell}{i}\binom{m-\ell}{h-i}\sigma_{\ell,m-\ell}(A) \\
&= \sum_{\ell=0}^{\min(m,n-m)}\ \sum_{i=\max(0,\ell-m+h)}^{\min(\ell,h)}\binom{\ell}{i}\binom{m-\ell}{h-i}\frac{(n-m-h+1)_{h-i}}{(-m)_{h-i}}\sigma_{\ell,m-\ell}(A).
\end{align*}
这几乎就是我们要的公式, 除了开头的常数 $(-1)^h\binom{n-m}{h}^{-1}$. 球函数 $\phi_h$ 由 $\phi_h(A_0) = 1$ 规范化, 所以有 $\phi_h = \frac{1}{\psi(A_0)}\psi$. 所以为完成证明, 只需证明 $\psi(A_0) = (-1)^h\binom{n-m}{h}$.
注意 $\sigma_{\ell,m-\ell}(A)(A_0)$ 除非 $\ell = 0$ 为 $0$, 且 $\sigma_{0,m}(A)(A_0) = 1$, 所以
\begin{align*}
\psi(A_0) &= \binom{m}{h}\frac{(n - m - h + 1)_h}{(-m)_h} \\
&= \frac{m!}{h!(m-h)!}\cdot\frac{(n-m)!}{(n-m-h)!}\cdot\frac{(-1)^h(m-h)!}{m!} \\
&= (-1)^h\binom{n-m}{h}.
\end{align*}
(g). 试着忽略 (c), (d), (e) 和 (f).
函数 $\phi_h$ 球, 所以它在 $\Omega_m$ 中的 $S_m \times S_{n-m}$-轨道上为常数, 而这些轨道由习题 \ref{ex:V.8.2.1}(d) 是以 $A$ 为中心的球面. 由 (a) 和习题 \ref{ex:V.8.2.2}(a), 这意味着 $\phi_h$ 是函数 $\sigma_{\ell,m-\ell}(A)$ ($0 \leq \ell \leq \min(m, n-m)$) 的线性组合. 又, 由习题 \ref{ex:V.8.2.3}(g), 函数 $\phi_h$ 是 $\Delta$ 的特征值 $m(n-m) - h(n-h+1)$ 的特征向量. 设 $S = \{B \in \Omega_m \mid d(A, B) = 1\}$. 则 $\sigma_{1,m-1}(A)$ 是 $S$ 的特征函数, 且看到 $\phi_h$ 在 $S$ 上为常数, $\phi_h$ 中 $\sigma_{1,m-1}(A)$ 的系数是 $\frac{1}{|S|}\sum_{B \in S}\phi_h(B)$.
另一方面, $\sum_{B \in S}\phi_h(B)$ 由 $\Delta$ 的定义等于 $\Delta\phi_h(A)$, 而这等于 $(m(n-m) - h(n-h+1))\phi_h(A)$. 又, 由于 $\phi_h$ 球, 必有 $\phi_h(A) = 1$. 最后, $\phi_h$ 中 $\sigma_{1,m-1}(A)$ 的系数是 $\frac{1}{|S|}(m(n-m) - h(n-h+1))$. 为完成计算, 只需证明 $|S| = m(n-m)$. 这由从 $A$ 中移除一个元素 ($m$ 种选择) 并添加 $\{1, \ldots, n\} - A$ 的一个元素 ($n - m$ 种选择) 得到 $S$ 的每个元素 $B$ 这一事实得出.
\end{solution}

\subsection{问题: \index{Satake 同构}Satake 同构}
\label{sec:V.8.3}

此问题摘录自 Cartier 在 Corvallis 的综述 {[}6{]}, 尤其是第四章, 以及 Kottwitz 未发表的笔记.
我们使用下列约定:
\begin{itemize}
\item 若 $G$ 是局部紧群, 记 $\mu_G$ 为 $G$ 上的左 Haar 测度 (有时按特定方式规范化). 记对应的 $L^p$ 空间为 $L^p(G)$.
\item 若 $G$ 是局部紧群且 $H \subset G$ 是闭子群, 记 $\mu_{G/H}$ (相应地 $\mu_{H\backslash G}$) 为由测度 $\mu_G$ 和 $\mu_H$ 诱导的 $G/H$ (相应地 $H\backslash G$) 上左 (相应地右) $G$-不变测度. 见命题 \ref{prop:V.1.3}. 为构造测度 $\mu_{H\backslash G}$, 需要 $\mu_G$ 右不变.
\end{itemize}
为使事情更具体, 我为 $\GL_3(\Q_p)$ 写了此问题. 一切都推广到 $\GL_n(\Q_p)$ (对任意 $n \geq 1$), 事实上大多数问题都有只做最小修改就适用于 $\GL_n(\Q_p)$ 的解法.
设 $n = 3$. 写 $G = \GL_3(\Q_p)$, $K = \GL_3(\Z_p)$. 记 $B$ 为 $G$ 中上三角矩阵的子群, $T$ 为 $G$ 中对角矩阵的子群, $N$ 为 $G$ 中幂幺上三角矩阵 (即所有对角线元素等于 $1$ 的上三角矩阵) 的子群. 写 $X = \Z^3$ 和
\[
X^+ = \{(\lambda_1, \lambda_2, \lambda_3) \in \Z^3 \mid \lambda_1 \geq \lambda_2 \geq \lambda_3\}.
\]
若 $\lambda = (\lambda_1, \lambda_2, \lambda_3) \in X$, 令
\[
p^{\lambda} = \begin{pmatrix} p^{\lambda_1} & 0 & 0 \\ 0 & p^{\lambda_2} & 0 \\ 0 & 0 & p^{\lambda_3} \end{pmatrix} \in T.
\]
固定 $G$, $B$, $T$, $N$ 和 $K$ 上的左 Haar 测度 $\mu_G$, $\mu_B$, $\mu_T$, $\mu_N$ 和 $\mu_K$; 若没有混淆的危险, 在积分中只写 $dx$ 代替 $d\mu_G(x)$ 等. 我们要求 $H \in \{G, B, T, N, K\}$ 时 $\operatorname{vol}(H \cap K) = 1$ 来规范化所有 Haar 测度. 在此问题中, 记 $\mathcal{C}_c(X)$ 为具紧支撑的局部常值函数 $f \colon X \to \CC$ 的空间. (其中``局部常值''就是你想象的意思: $f \colon X \to \CC$ 局部常值, 若对每个 $x \in X$, 存在 $X$ 中 $x$ 的邻域 $U$ 使得 $f|_U$ 为常数. 注意这蕴含 $f$ 连续.)
设 $\mathcal{H} = \mathcal{C}_c(K\backslash G/K)$, $\mathcal{H}_T = \mathcal{C}_c(T/(T \cap K))$, 带由卷积乘积给出的代数结构. (很容易验证卷积乘积保持``局部常值''条件.) 此问题的目标是理解代数 $\mathcal{H}$ 的结构.
\begin{enumerate}
\item[(1)] 定义 $\delta \colon B \to \R_{>0}$:
\[
\delta\begin{pmatrix} a_1 & * & * \\ 0 & a_2 & * \\ 0 & 0 & a_3 \end{pmatrix} = \prod_{i=1}^n|a_i|_p^{n+1-2i}.
\]
证明 $\delta^{-1}$ 是 $B$ 的模函数, 且对每个 $f \in \mathcal{C}_c(B)$ 和每个 $b' \in B$,
\[
\int_B f(b'b)\,\,\dd b = \delta^{-1}(b')\int_B f(b)\,\,\dd b.
\]
\item[(2)] 证明 $G = BK$, 且对每个 $f \in \mathcal{C}_c(G)$,
\[
\int_G f(g)\,\dd g = \int_B\int_K f(bk)\,\dd b\,\dd k = \int_T\int_N\int_K f(tnk)\,\dd t\,\dd n\,\dd k.
\]
\item[(3)] 设 $(e_1, e_2, e_3)$ 是 $\Q_p^3$ 的典范基. 证明映射 $g \longmapsto \Z_pg(e_1) + \Z_pg(e_2) + \Z_pg(e_3)$ 诱导 $G/K$ 与 $\Q_p^3$ 的秩 $3$ 自由 $\Z_p$-子模的集合之间的双射.
\item[(4)] 证明
\[
G = \bigsqcup_{\lambda \in X^+}Kp^{\lambda}K.
\]
(这称为 Cartan 分解.)
\item[(5)] 证明 $(G, K)$ 是 Gelfand 对. (提示: $\theta(x) = (x^T)^{-1}$.)
\end{enumerate}
设 $z = (z_1, z_2, z_3) \in (\CC^{\times})^3$. 定义态射 $\chi_z \colon B \to \CC^{\times}$:
\[
\chi_z\begin{pmatrix} a_1 & * & * \\ 0 & a_2 & * \\ 0 & 0 & a_3 \end{pmatrix} = \prod_{i=1}^n z_i^{v_p(a_i)},
\]
其中对每个 $a \in \Q_p^{\times}$ 写 $|a|_p = p^{-v_p(a)}$ (所以 $v_p(a) \in \Z$). 记 $V_z$ 为满足下列条件的函数 $f \colon G \to \CC$ 的空间:
\begin{enumerate}
\item[(i)] 对每个 $b \in B$ 和每个 $g \in G$, $f(bg) = \chi_z(b)\delta(b)^{1/2}f(g)$;
\item[(ii)] 存在 $G$ 的开紧子群 $K_1$ 使得对每个 $g \in G$ 和每个 $k \in K_1$, $f(gk) = f(g)$.
\end{enumerate}
让 $G$ 按 $g \cdot f = R_g(f)$ 作用于 $V_z$.
\begin{enumerate}
\item[(6)] 证明 $V_z$ 是 $G$ 的表示 (注意我没有说关于连续性的任何事) 且 $\dim(V_z^K) = 1$.
\end{enumerate}
注记. (不必做此题.) $V_z$ 的元素局部常值但一般没有紧支撑, 因为条件 (i). 又, 我们没有在 $V_z$ 上定义拓扑, 所以不能说 $V_z$ 是 $G$ 的连续表示. 事实上, 它是所谓的光滑表示 (smooth representation), 这意味着 $V_z$ 的每个元素都被 $G$ 的一个开紧子群固定 (这由条件 (ii) 自动). 这意味着若在 $V_z$ 上放离散拓扑, 表示连续, 而且对完全不连通群的表示这是更自然的条件. 若想要希尔伯特空间上的连续表示, 会在 $V_z$ 上放内积 $\langle f_1, f_2 \rangle = \int_K f_1(k)\overline{f_1(k)}\,\dd k$ 并对这个内积完备化. 注意即使如此, 表示也不会酉, 除非 $|z_1| = |z_2| = |z_3| = 1$.
记 $\phi_z$ 为满足 $\phi_z(1) = 1$ 的 $V_z^K$ 的唯一元素.\footnote{它是球函数 (除非一般地无界), 因此有记号 $\phi$. 但除非你先证明它, 不能使用这一事实!} 对每个 $f \in \mathcal{H}$, 记 $f^{\vee} \colon (\CC^{\times})^3 \to \CC$ 为函数:\footnote{这道德上是球傅里叶变换, 但不幸的是 $p$-进群的通常约定不同于我们在课堂上使用的约定.}
\[
z \longmapsto \int_G f(g)\phi_z(g)\,\dd g.
\]
若 $f \in \mathcal{H}_T$, 定义 $f^{\vee} \colon (\CC^{\times})^3 \to \CC$:
\[
f^{\vee}(z) = \int_T f(t)\chi_z(t)\,\dd t.
\]
\begin{enumerate}
\item[(7)] 证明对所有 $f_1, f_2 \in \mathcal{H}$, $(f_1 * f_2)^{\vee} = f_1^{\vee}f_2^{\vee}$.
\item[(8)] 证明
\[
T = \bigsqcup_{\lambda \in X}(T \cap K)p^{\lambda}.
\]
\item[(9)] 证明 $f \longmapsto f^{\vee}$ 是从 $\mathcal{H}_T$ 到 $(\CC^{\times})^3$ 上在 $z_1^{\pm 1}, z_2^{\pm 1}, z_3^{\pm 1}$ 中多项式的函数的代数的 $\CC$-代数同构; 把这个代数等同于 $\CC[z_1^{\pm 1}, z_2^{\pm 1}, z_3^{\pm 1}]$.
\end{enumerate}
对每个 $f \in \mathcal{C}_c(G)$, 定义函数 $f^{(B)} \colon T \to \CC$:
\[
f^{(B)}(t) = \delta(t)^{1/2}\int_N f(tn)\,\dd n.
\]
\begin{enumerate}
\item[(10)] 证明对每个 $f \in \mathcal{C}_c(G)$, $f^{(B)} \in \mathcal{C}_c(T)$, 且若 $f \in \mathcal{H}$ 则 $f^{(B)} \in \mathcal{H}_T$.
\item[(11)] 证明对每个 $f \in \mathcal{H}$ 和每个 $z \in (\CC^{\times})^3$, $f^{\vee}(z) = (f^{(B)})^{\vee}(z)$.
\end{enumerate}
定义 $D \colon T \to \R_{>0}$:
\[
D\begin{pmatrix} a_1 & 0 & 0 \\ 0 & a_2 & 0 \\ 0 & 0 & a_3 \end{pmatrix} = \prod_{i,j \in \{1,\ldots,n\}, i \neq j}|1 - a_ia_j^{-1}|_p.
\]
若 $f \in \mathcal{C}_c(G)$, $t \in T$ 使得 $D(t) \neq 0$, 令
\[
\mathcal{O}_t(f) = \int_{T\backslash G}f(g^{-1}tg)\,\dd\mu_{T\backslash G}.
\]
\begin{enumerate}
\item[(12)] 证明对每个 $f \in \mathcal{C}_c(N)$ 和每个满足 $D(t) \neq 0$ 的 $t \in T$,
\[
\int_N f(n)\,\dd n = D(t)^{1/2}\delta(t)^{1/2}\int_N f(ntn^{-1}t^{-1})\,\dd n.
\]
(提示: 考虑由
\[
U = \left\{\begin{pmatrix} 1 & 0 & * \\ 0 & 1 & 0 \\ 0 & 0 & 1 \end{pmatrix}\right\}
\]
定义的子群 $U \subset N$, 并用公式 $\int_N f(n)\,\dd n = \int_{N/U}\int_U f(nu)\,\dd\mu_{N/U}(n)\,\dd\mu_U(u)$.)
\item[(13)] 证明对每个 $f \in \mathcal{H}$ 和每个满足 $D(t) \neq 0$ 的 $t \in T$,
\[
\mathcal{O}_t(f) = D(t)^{-1/2}f^{(B)}(t).
\]
\item[(14)] 证明对每个 $f \in \mathcal{H}$, 函数 $f^{\vee} \colon (\CC^{\times})^3 \to \CC$ 是 $\CC[z_1^{\pm 1}, z_2^{\pm 1}, z_3^{\pm 1}]^{S_3}$ 的元素, 即 $z_1^{\pm 1}, z_2^{\pm 1}, z_3^{\pm 1}$ 的对称多项式函数的代数:
\[
\{p(z_1, z_2, z_3) \in \CC[z_1^{\pm 1}, z_2^{\pm 1}, z_3^{\pm 1}] \mid \forall \sigma \in S_3,\, p(z_{\sigma(1)}, z_{\sigma(2)}, z_{\sigma(3)}) = p(z_1, z_2, z_3)\}.
\]
\end{enumerate}
在 $X = \Z^3$ 上定义偏序, 说 $(\mu_1, \mu_2, \mu_3) \leq (\lambda_1, \lambda_2, \lambda_3)$ 当且仅当 $\mu_1 \leq \lambda_1$, $\mu_1 + \mu_2 \leq \lambda_1 + \lambda_2$ 且 $\mu_1 + \mu_2 + \mu_3 = \lambda_1 + \lambda_2 + \lambda_3$. (注意最后一个关系是等式!) 对每个 $\lambda = (\lambda_1, \lambda_2, \lambda_3) \in X$, 定义 $z^{\lambda} \in \CC[z_1^{\pm 1}, z_2^{\pm 1}, z_3^{\pm 1}]$:
\[
z^{\lambda} = z_1^{\lambda_1}z_2^{\lambda_2}z_3^{\lambda_3}
\]
和 $f_{\lambda} \in \mathcal{H}$:
\[
f_{\lambda} = \mathbb{1}_{Kp^{\lambda}K}.
\]
\begin{enumerate}
\item[(15)] 证明 $(f_{\lambda})_{\lambda \in X^+}$ 是 $\mathcal{H}$ 的基.
\item[(16)] 对每个 $\lambda = (\lambda_1, \lambda_2, \lambda_3) \in X^+$, 设
\[
c_{\lambda} = \sum_{\sigma \in S_3}z^{(\lambda_{\sigma(1)}, \lambda_{\sigma(2)}, \lambda_{\sigma(3)})}.
\]
证明 $(c_{\lambda})_{\lambda \in X^+}$ 是 $\CC[z_1^{\pm 1}, z_2^{\pm 1}, z_3^{\pm 1}]^{S_3}$ 的基.
\item[(17)] 对每个 $\lambda \in X$, 写
\[
(f_{\lambda})^{\vee} = \sum_{\mu \in X}c_{\lambda}(\mu)z^{\mu},
\]
其中 $c_{\lambda}(\mu) \in \CC$. 此题的目标是证明下面两个事实:
\begin{enumerate}
\item[(i)] 若 $\lambda, \mu \in X^+$ 且 $\mu \nleq \lambda$, 则 $c_{\lambda}(\mu) = 0$.
\item[(ii)] 对每个 $\lambda \in X^+$,
\[
c_{\lambda}(\lambda) = p^{\frac{1}{2}\sum_{i=1}^n(n+1-2i)\lambda_i}.
\]
\end{enumerate}
若 $\mu = (\mu_1, \mu_2, \mu_3) \in X$, 写 $\check{\mu} = (\mu_3, \mu_2, \mu_1)$.
\begin{enumerate}
\item[(a)] 对每个 $g = (g_{ij}) \in M_3(\Q_p)$, 设 $\|g\| = \sup_{i,j}|g_{ij}|_p$. 证明对每个 $g \in M_3(\Q_p)$ 和 $k, k' \in K$, $\|g\| = \|kgk'\|$.
\item[(b)] 设 $r \in \{1, \ldots, n\}$, $\Omega_r$ 是 $\{1, \ldots, n\}$ 的基数为 $r$ 的子集的集合. 若 $g = (g_{ij}) \in M_3(\Q_p)$, 定义 $\Lambda_rg \colon \Omega_r \times \Omega_r \to \Q_p$:
\[
\Lambda_r(g)(A, A') = \det((g_{ij})_{i \in A, j \in A'}),
\]
且 $\|\Lambda_rg\|$:
\[
\|\Lambda_rg\| = \sup_{A,A' \in \Omega_r}|\Lambda_rg(A, A')|_p.
\]
证明对每个 $g \in M_3(\Q_p)$ 和 $r \in \{1, \ldots, n\}$:
\begin{enumerate}
\item[$(\alpha)$] 若 $k, k' \in K$, 则 $\|\Lambda_r(kgk')\| = \|\Lambda_rg\|$.
\item[$(\beta)$] 若 $t \in T$, $n \in N$, 则 $\|\Lambda_r(tn)\| \geq \|\Lambda_r(t)\|$.
\item[$(\gamma)$] 若 $g = p^{\lambda}$, $\lambda = (\lambda_1, \lambda_2, \lambda_3) \in X^+$, 则
\[
\|\Lambda_rg\| = p^{-(\lambda_{3-r+1} + \cdots + \lambda_3)}.
\]
\end{enumerate}
\item[(c)] 对所有 $\lambda, \mu \in X$, 证明
\[
c_{\lambda}(\mu) = f_{\lambda}^{(B)}(p^{\mu}) = f_{\lambda}^{(B)}(p^{\check{\mu}}).
\]
\item[(d)] 设 $\lambda, \mu \in X^+$. 若存在 $n \in N$ 使得 $p^{\mu}n \in Kp^{\lambda}K$, 证明 $\mu \leq \lambda$.
\item[(e)] 证明 (i).
\item[(f)] 设 $\lambda \in X^+$, $t = p^{\lambda}$. 对每个 $n \in N$, 证明 $tn \in KtK$ 当且仅当 $n \in N \cap K$.
\item[(g)] 证明 (ii).
\end{enumerate}
\item[(18)] 证明 $f \longmapsto f^{\vee}$ 诱导 $\CC$-代数同构 $\mathcal{H} \xrightarrow{\sim} \CC[z_1^{\pm 1}, z_2^{\pm 1}, z_3^{\pm 1}]^{S_3}$. (此同构称为 Satake 同构.)
\end{enumerate}

\begin{solution}
(1) 我们在习题 \ref{ex:I.5.3.8}(d) 中看到 $\delta^{-1}$ 是 $B$ 的模函数. 第二个结论由 $db$ 的左不变性和命题 \ref{prop:I.2.8} 得出.
(2) 注意, 由于 $B$ 和 $K$ 是 $G$ 的子群, 证明 $G = BK$ 与证明 $G = KB$ 等价. (此分解称为 Iwasawa 分解.) 我们用高斯消去算法 (a.k.a. 行化简) 的变体证明第二个结论. 在此算法中, 对 $\GL_3(\Q_p)$ 的矩阵应用三种操作:
\begin{itemize}
\item[(a)] 交换两行;
\item[(b)] 把一行乘以 $\Q_p^{\times}$ 的元素 $\alpha$;
\item[(c)] 把一行的 $\Q_p$ 的元素 $\beta$ 的倍数加到另一行.
\end{itemize}
每个这样的操作对应于左乘一个初等矩阵. 而且, 此矩阵在情形 (a) 属于 $\GL_3(\Z_p)$, 在情形 (b) 中若 $\alpha \in \Z_p^{\times}$ 属于 $\GL_3(\Z_p)$, 在情形 (c) 中若 $\beta \in \Z_p$ 属于 $\GL_3(\Z_p)$. 所以只需证明: 若对操作应用刚才对操作注意到的限制的操作序列 (a), (b), (c), 总可以得到上三角矩阵. 所以设 $g = (g_{ij})$ 是 $\GL_3(\Q_p)$ 的元素. 由于 $g$ 可逆, 它第一列有非零元素. 交换 $g$ 的行 (操作 (a)) 后, 可以假设 $g_{11} \neq 0$ 且对 $i = 2, 3$ 有 $|g_{11}|_p \geq |g_{i1}|_p$. 把 $g$ 的第一行的 $-g_{11}^{-1}g_{21}$ (相应地 $-g_{11}^{-1}g_{31}$) 倍加到 $g$ 的第二 (相应地第三) 行 (操作 (b); 注意由第一步我们是用 $\Z_p$ 的元素乘以第一行) 后, 可以假设 $g_{11}$ 是 $g$ 第一列的唯一非零元素. 则 $g_{22}$ 和 $g_{23}$ 至少有一个非零 (否则 $g$ 的秩 $\leq 2$), 所以交换 $g$ 的后两行后, 可以假设 $g_{22} \neq 0$ 且 $|g_{22}|_p \geq |g_{23}|_p$. 然后, 把第二行的 $-g_{22}^{-1}g_{23}$ 倍加到第三行后, 可以假设 $g_{23} = 0$. 此时矩阵 $g$ 上三角, 完成了.
现在证明第二个结论. 考虑群 $G_0 = B \times K$, 及其子群 $H = \{(x, x) \mid x \in B \cap K\}$. 考虑 $G_0$ 上的测度 $\mu_{G_0} = \mu_B \times \mu_K$; 这显然是左 Haar 测度. 还在紧群 $H$ 上放规范化 Haar 测度 $\mu_H$. 考虑 $G_0$ 在 $G$ 上的连续作用 $(b, k) \cdot g = bgk^{-1}$. 由此题的第一部分, 这个作用传递. 又, $1 \in G$ 的稳定化子显然是 $H$, 所以得到满连续映射 $\phi \colon G_0/H \to G$. 证明 $\phi$ 也开. 若 $x \in G_0/H$, 则 $((B \cap K) \times K)x$ 是 $x$ 的开邻域 (因为 $B \cap K$ 在 $B$ 中开), 且它在 $\phi$ 下的像是 $(B \cap K)\phi(x)K$, 这在 $G$ 中开, 因为 $K$ 是 $G$ 的开子群. 所以 $\phi$ 开, 从而是同胚. 设 $\mu$ 是 $G$ 上固定 Haar 测度在 $\phi^{-1}$ 下的原像. 由于 $G$ 幺模 (由习题 \ref{ex:I.5.3.8}(c)), 测度 $\mu$ 被 $G_0$ 左不变, 所以正线性泛函 $\mathcal{C}_c(G_0) \to \CC$, $h \longmapsto \int_{G_0/H}h^H(x)\,\dd\mu(x)$ 左不变, 因此它是 $\mu_{G_0}$ 的正倍数, 设为 $c\mu_{G_0}$ ($c > 0$). 设 $f \in \mathcal{C}_c(G)$. 定义函数 $h \in \mathcal{C}(G_0)$: $h(b, k) = f(bk^{-1})$; 有 $h(b, k) = 0$ 若 $b \notin (\operatorname{supp} f)K$, 所以 $h$ 有紧支撑. 应用习题 \ref{ex:I.5.3.6}, 得到
\[
\int_{G_0}h(b, k)\,\dd\mu_{G_0}(b, k) = \int_B\int_K f(bk^{-1})\,\dd b\,\dd k = c\int_{G_0/H}h^H(x)\,\dd\mu(x).
\]
设 $x \in G_0/H$. 选取 $x$ 的原像 $(b, k) \in G_0$, 使得 $\phi(x) = bk$. 则
\[
h^H(x) = \int_H h((b, k)h)\,\dd h = \int_{B \cap K}f(bhh^{-1}k)\,\dd h = f(bk).
\]
所以
\[
\int_B\int_K f(bk)\,\dd b\,\dd k = c\int_G f(g)\,\dd g,
\]
对每个 $f \in \mathcal{C}_c(G)$. 为计算 $c$, 把它应用于 $f = \mathbb{1}_K$ (这是具紧支撑的连续函数, 因为由习题 \ref{ex:I.5.1.4}(m), $K$ 是 $G$ 的开紧子群). 左端是 $\mu_B(B)\mu_K(K) = 1$, 右端是 $\mu_G(K) = 1$ (由 $\mu_G$ 的选取). 所以 $c = 1$.
最后, 把习题 \ref{ex:I.5.3.6} 应用于 $B$ 的子群 $T$ 和 $N$, 得到正常数 $c'$ 使得对每个 $f \in \mathcal{C}_c(B)$,
\[
\int_B f\,\dd\mu_B = c'\int_T\int_N f(tn)\,\dd\mu_T(t)\,\dd\mu_N(n).
\]
为证明 $c' = 1$, 只需对 $f = \mathbb{1}_{B \cap K}$ 计算两边, 并用 (容易的) 事实 $B \cap K = (T \cap K)(N \cap K)$.
(3) 首先注意如下容易的事实: 设 $(x_1, x_2, x_3)$ 是 $\Q_p^3$ 的元素族; 则此族在 $\Z_p$ 上线性无关当且仅当它在 $\Q_p$ 上线性无关. 事实上, 在 $\Q_p$ 上线性无关显然蕴含在 $\Z_p$ 上线性无关 (因为 $\Z_p \subset \Q_p$). 反之, 假设 $(x_1, x_2, x_3)$ 在 $\Z_p$ 上线性无关; 若有 $a_1x_1 + a_2x_2 + a_3x_3 = 0$, $a_1, a_2, a_3 \in \Q_p$, 则存在 $m \in \Z$ 使得 $p^ma_1, p^ma_2, p^ma_3 \in \Z_p$, 且关系 $p^ma_1x_1 + p^ma_2x_2 + p^ma_3x_3 = 0$ 蕴含 $p^ma_1 = p^ma_2 = p^ma_3 = 0$, 所以 $a_1 = a_2 = a_3 = 0$.
设 $\mathcal{L}$ 是 $\Q_p^3$ 的秩 $3$ 自由 $\Z_p$-子模的集合, $L_0 = \Z_pe_1 + \Z_pe_2 + \Z_pe_3 = \Z_p^3 \in \mathcal{L}$. 若 $L \in \mathcal{L}$, $g \in \GL_3(\Q_p)$, 则 $g(L)$ 也是 $\mathcal{L}$ 的元素. 事实上, 若 $(x_1, x_2, x_3)$ 是 $L$ 的 $\Z_p$-基, 则族 $(x_1, x_2, x_3)$ 也在 $\Q_p$ 上线性无关, 所以族 $(g(x_1), g(x_2), g(x_3))$ 在 $\Q_p^3$ 上线性无关, 从而在 $\Z_p$ 上线性无关, 所以它生成的 $\Q_p^3$ 的 $\Z_p$-子模是 $\mathcal{L}$ 的元素. 这证明 $(g, L) \longmapsto g(L)$ 定义 $G$ 在 $\mathcal{L}$ 上的左作用. 又, $L_0$ 的稳定化子显然是 $\GL_3(\Z_p)$.
证明这个作用传递, 这将给出所要的双射 $G/K \xrightarrow{\sim} \mathcal{L}$. 设 $L \in \mathcal{L}$, $(x_1, x_2, x_3)$ 是 $L$ 在 $\Z_p$ 上的基. 则看到族 $(x_1, x_2, x_3)$ 也在 $\Q_p$ 上线性无关. 所以 $(x_1, x_2, x_3)$ 是 $\Q_p$-向量空间 $\Q_p^3$ 的基, 因此存在 $g \in \GL_3(\Q_p)$ 使得对 $i = 1, 2, 3$ 有 $x_i = g(e_i)$, 然后 $L = g(L_0)$.
(4) 首先证明 $G = \bigcup_{\lambda \in X^+}Kp^{\lambda}K$. 设 $g \in G$. 设 $x_1, x_2, x_3 \in \Q_p^3$ 是 $g$ 的列, $L$ 是 $(x_1, x_2, x_3)$ 生成的 $\Q_p^3$ 的 $\Z_p$-子模; 在 (3) 的解答的记号中, 这是 $g(L_0)$. 由于 $\Q_p = \bigcup_{m \in \Z}p^m\Z_p$, 可以选择 $m \in \Z$ 使得矩阵 $p^mg$ 的元素都在 $\Z_p$ 中. 则 $(p^mx_1, p^mx_2, p^mx_3)$ 生成的 $\Q_p^3$ 的 $\Z_p$-子模 $L'$ 满足 $L' \supset L_0 = p^mL$. 由 Dummit 和 Foote 第 12 章定理 4, 存在 $L'$ 的基 $(y_1, y_2, y_3)$ 和元素 $a_1, a_2, a_3 \in \Z_p - \{0\}$ 使得 $a_3|a_2|a_1$ 且 $(a_1y_1, a_2y_2, a_3y_3)$ 是 $L_0$ 的基. 由于 $\Z_p - \{0\}$ 的每个元素 $a$ 都可以 (唯一地) 写成 $a = p^nu$, $n \in \Z_{\geq 0}$, $u \in \Z_p^{\times}$, 可以假设 $a_i = p^{n_i}$, 其中 $n_1, n_2, n_3 \in \Z_{\geq 0}$, $n_1 \geq n_2 \geq n_3$. 设 $g'$ (相应地 $k$) 是列 $a_1y_1, a_2y_2, a_3y_3$ (相应地 $y_1, y_2, y_3$) 的矩阵, $t$ 是对角线元素 $a_1, a_2, a_3$ 的对角矩阵; 换句话说, $t = p^{(n_1,n_2,n_3)}$. 则 $k \in K$, $g' = kt$. 又, 由于 $L_0 = g'(L_0) = p^mL = p^mg(L_0)$, 问题 (3) 蕴含存在 $k' \in K$ 使得 $p^mg = g'k'$. 所以得到
\[
g = p^{-m}g'k' = kp^{(n_1-m,n_2-m,n_3-m)}k' \in Kp^{(n_1-m,n_2-m,n_3-m)}K.
\]
注意还证明了: 若 $m = 0$ (例如若 $g$ 的所有元素都已在 $\Z_p$ 中), 使得 $L \subset L_0$, 则 $L_0/L \simeq (\Z/p^{n_1}\Z) \times (\Z/p^{n_2}\Z) \times (\Z/p^{n_3}\Z)$.
现在证明这个并是不相交并. 设 $\lambda, \lambda' \in X^+$, 假设 $Kp^{\lambda}K \cap Kp^{\lambda'}K \neq \emptyset$. 把 $p^{\lambda}$ 和 $p^{\lambda'}$ 乘以 $m$ 足够大的 $p^m$ 后, 可以假设 $\lambda$ 和 $\lambda'$ 在 $\Z_{\geq 0}^3$ 中. 设 $g \in Kp^{\lambda}K \cap Kp^{\lambda'}K$, $L = g(L_0)$. 由于 $g$ 的元素是 $\Z_p$ 的元素 (由对 $\lambda$ 的假设), 有 $L \subset L_0$. 写 $\lambda = (n_1, n_2, n_3)$, $\lambda' = (n'_1, n'_2, n'_3)$, 其中 $n_1 \geq n_2 \geq n_3$, $n'_1 \geq n'_2 \geq n'_3$. 由上一段,
\[
L_0/L \simeq (\Z/p^{n_1}\Z) \times (\Z/p^{n_2}\Z) \times (\Z/p^{n_3}\Z) \simeq (\Z/p^{n'_1}\Z) \times (\Z/p^{n'_2}\Z) \times (\Z/p^{n'_3}\Z).
\]
由 Dummit 和 Foote 第 12 章定理 9, 这蕴含对 $i = 1, 2, 3$ 有 $n_i = n'_i$, 即 $\lambda = \lambda'$.
(5) 考虑映射 $\theta \colon G \to G$, $x \longmapsto (x^T)^{-1}$. 这是群的连续同构, 因为 $\theta^{-1} = \theta$ 而是同胚. 也有 $\theta(K) = K$. 设 $g \in G$. 由 (4), 可以把 $g$ 写成 $g = kp^{\lambda}k'$, 其中 $k, k' \in K$, $\lambda \in \Z^3$. 则
\[
\theta(g) = \theta(k)p^{-\lambda}\theta(k') = \theta(k)k'g^{-1}k\theta(k') \in Kg^{-1}K.
\]
所以由命题 \ref{prop:V.2.5}, $(G, K)$ 是 Gelfand 对.
(6) 所有函数 $f \colon G \to \CC$ 的空间 (带 $G$ 的作用 $x \cdot f = R_xf$) 是 $G$ 的表示. 验证 $V_z$ 是这个空间的子表示. 首先, 它是向量子空间:
\begin{itemize}
\item 函数 $0$ 在 $V_z$ 中;
\item 条件 (i) 显然被线性组合保持;
\item 设 $f, f' \in V_z$, $K_1$ 和 $K'_1$ 是 $G$ 的开紧子群, 使得 $f$ (相应地 $f'$) 被 $K_1$ (相应地 $K'_1$) 右不变; 则 $K_1 \cap K'_1$ 是 $G$ 的紧子群, 且 $f$ 和 $f'$ 的任何线性组合被 $K_1 \cap K'_1$ 右不变, 所以满足条件 (ii).
\end{itemize}
现在验证 $V_z$ 是 $G$-不变的. 设 $f \in V_z$, $x \in G$. 对每个 $b \in B$ 和每个 $g \in G$,
\[
(R_xf)(bg) = f(bgx) = \chi_z(b)\delta(b)^{1/2}f(gx) = \chi_z(b)\delta(b)^{1/2}R_xf(g),
\]
所以 $R_xf$ 满足条件 (i). 设 $K_1$ 是 $G$ 的开紧子群, $f$ 被 $K_1$ 右不变. 则对每个 $g \in G$ 和 $k \in K_1$,
\[
R_xf(gxkx^{-1}) = f(gxk) = f(gx) = R_xf(g),
\]
所以 $R_xf$ 被 $G$ 的开紧子群 $xK_1x^{-1}$ 右不变, 因此满足条件 (ii).
最后计算 $\dim(V_z^K)$. 回忆映射 $G_0 := B \times K \to G$, $(b, k) \longmapsto bk^{-1}$ 把 $G$ 等同于商 $G_0/H$, 其中 $H = \{(x, x) \mid x \in B \cap K\}$. 考虑函数 $\varphi_z^0 \colon G_0 \to \R_{>0}$, $(b, k) \longmapsto \chi_z(b)\delta(b)^{1/2}$. 这是连续的群态射, 所以它在紧子群 $H$ 上平凡, 因此它下降到商 $G_0/H$, 并定义从 $G$ 到 $\R_{>0}$ 的连续函数, 记为 $\phi_z$. 由构造, 函数 $\phi_z$ 满足
\[
\phi_z(bk) = \chi_z(b)\delta(b)^{1/2}
\]
对每个 $b \in B$ 和每个 $k \in K$, 所以 $\phi_z(1) = 1$ 且 $\phi_z \in V_z^K$. 为完成证明, 证明 $V_z^K$ 由 $\phi_z$ 生成. 设 $f \in V_z^K$. 若 $g \in G$, 由问题 (2), 可以找到 $b \in B$, $k \in K$ 使得 $g = bk$, 然后
\[
f(g) = f(b) = \chi_z(b)\delta(b)^{1/2}f(1) = \phi_z(g)f(1).
\]
所以 $f$ 是 $\phi_z$ 的倍数.
(7) 固定 $z \in (\CC^{\times})^3$.
设 $f \in \mathcal{H}$. 定义 $\tilde{f} \in \mathcal{H}$: $\tilde{f}(x) = f(x^{-1})$. 证明 $\phi_z * \tilde{f} = f^{\vee}(z)\phi_z$. 首先注意 $\phi_z * \tilde{f}$ 被 $K$ 右不变, 因为 $\tilde{f}$ 被 $K$ 右不变 (这由命题 \ref{prop:V.1.8} 得出). 设 $x \in G$, $b \in B$. 则
\[
\phi_z * \tilde{f}(bx) = \int_G\phi_z(bxy)f(y)\,\dd y = \int_G\chi_z(b)\delta(b)^{1/2}\phi_z(xy)f(y)\,\dd y = \chi_z(b)\delta(b)^{1/2}\phi_z * \tilde{f}(x)
\]
(因为 $\phi_z \in V_z$). 所以 $\phi_z * \tilde{f} \in V_z^K$, 这由问题 (6) 蕴含 $\phi_z * \tilde{f}$ 是 $\phi_z$ 的纯量倍数. 为找到纯量, 在两个函数处取值 $1$: 有 $\phi_z(1) = 1$, 且
\[
\phi_z * \tilde{f}(1) = \int_G\phi_z(y)f(y)\,\dd y = f^{\vee}(z).
\]
现在设 $f_1, f_2 \in \mathcal{H}$. 则对每个 $x \in G$,
\begin{align*}
\widehat{f_1 * f_2}(x) &= \widetilde{f_1 * f_2}(x^{-1}) \\
&= \int_G f_1(x^{-1}y^{-1})f_2(y)\,\dd y \\
&= \int_G \tilde{f}_1(yx)\tilde{f}_2(y^{-1})\,\dd y = (\tilde{f}_2 * \tilde{f}_1)(x) \quad \text{(因为 } G \text{ 幺模)} \\
&= (\tilde{f}_1 * \tilde{f}_2)(x) \quad \text{(因为 } (G, K) \text{ 是 Gelfand 对)},
\end{align*}
所以 $\widehat{f_1 * f_2} = \tilde{f}_1 * \tilde{f}_2$. (实际上, 事实 $\widehat{f_1 * f_2} = \tilde{f}_2 * \tilde{f}_1$ 对我们的目的已足够好.) 用上一段的计算, 得到
\[
(f_1 * f_2)^{\vee}(z) = (\phi_z * \widehat{f_1 * f_2})(1) = (\phi_z * \tilde{f}_1 * \tilde{f}_2)(1) = ((f_1^{\vee}(z)\phi_z) * \tilde{f}_2)(1) = f_1^{\vee}(z)f_2^{\vee}(z).
\]
(8) 注意 $T \cap K$ 是使得 $t$ 和 $t^{-1}$ 的所有元素都在 $\Z_p$ 中的可逆对角矩阵的群, 即所有元素都在 $\Z_p^{\times}$ 中的群. 所以结论由 $\Q_p^{\times} = \bigsqcup_{m \in \Z}p^m\Z_p^{\times}$ (应用于矩阵的每个对角线元素) 这一事实得出, 而这个事实由习题 \ref{ex:I.5.1.4}(i) 得出.
(9) 映射 $f \longmapsto f^{\vee}$ 显然对 $f \in \mathcal{H}_T$ 线性. 证明它是代数态射. 设 $f_1, f_2 \in \mathcal{H}_T$, $z \in (\CC^{\times})^3$. 则
\begin{align*}
(f_1 * f_2)^{\vee}(z) &= \int_T (f_1 * f_2)(t)\chi_z(t)\,\dd t \\
&= \int_{T \times T}f_1(t')f_2((t')^{-1}t)\chi_z(t)\,\dd t\,\dd t' \\
&= \int_{T \times T}f_1(t')f_2(t)\chi_z(t't)\,\dd t\,\dd t' \\
&= \Bigl(\int_T f_1(t')\chi_z(t')\,\dd t'\Bigr)\Bigl(\int_T f_2(t)\chi_z(t)\,\dd t\Bigr) \quad \text{(因为 } \chi_z \text{ 乘法)} \\
&= f_1^{\vee}(z)f_2^{\vee}(z),
\end{align*}
如所愿.
由问题 (8), 族 $(\mathbb{1}_{(T \cap K)p^{\lambda}})_{\lambda \in X}$ 是 $\mathcal{H}_T$ 的基. 设 $\lambda = (\lambda_1, \lambda_2, \lambda_3) \in X$, $g_{\lambda} = \mathbb{1}_{(T \cap K)p^{\lambda}}$. 则对每个 $z = (z_1, z_2, z_3) \in (\CC^{\times})^3$,
\[
g_{\lambda}^{\vee}(z) = \int_{(T \cap K)p^{\lambda}}\chi_z(t)\,\dd t = \operatorname{vol}(T \cap K)\chi_z(p^{\lambda}) = z_1^{\lambda_1}z_2^{\lambda_2}z_3^{\lambda_3}.
\]
由于形如 $(z_1, z_2, z_3) \longmapsto z_1^{\lambda_1}z_2^{\lambda_2}z_3^{\lambda_3}$ ($(\lambda_1, \lambda_2, \lambda_3) \in X$) 的函数构成 $\CC[z_1^{\pm 1}, z_2^{\pm 1}, z_3^{\pm 1}]$ 的基, 这证明线性变换 $f \longmapsto f^{\vee}$ 把 $\mathcal{H}_T$ 的基送到 $\CC[z_1^{\pm 1}, z_2^{\pm 1}, z_3^{\pm 1}]$ 的基, 所以它是从 $\mathcal{H}_T$ 到 $\CC[z_1^{\pm 1}, z_2^{\pm 1}, z_3^{\pm 1}]$ 的同构.
(10) 由习题 \ref{ex:I.5.2.2}, $G$ 的开紧子群构成 $1$ 的邻域基. 设 $f \in \mathcal{C}_c(G)$. 由 $\mathcal{C}_c(G)$ 的定义, 对每个 $x \in G$, 存在 $G$ 的开紧子群 $K_x$ 使得 $f$ 在 $K_xxK_x$ 上为常数. 有 $\operatorname{supp} f \subset \bigcup_{x \in \operatorname{supp} f}K_xxK_x$, 且 $\operatorname{supp} f$ 紧, 所以存在 $x_1, \ldots, x_m \in \operatorname{supp} f$ 使得 $\operatorname{supp} f \subset \bigcup_{i=1}^m K_{x_i}x_iK_{x_i}$. 设 $L = \bigcap_{i=1}^m K_{x_i}$; 这是 $G$ 的开紧子群. 证明 $f$ 是双 $L$-不变的. 设 $x \in G$. 若 $x \in \operatorname{supp} f$, 则存在 $i \in \{1, \ldots, m\}$ 使得 $x \in K_{x_i}x_iK_{x_i}$. 则 $LxL \subset K_{x_i}^2x_iK_{x_i}^2$, 所以 $f$ 在 $LxL$ 上为常数. 若 $x \notin \operatorname{supp} f$, 要证明 $LxL \cap (\operatorname{supp} f) = \emptyset$. 但假设存在 $l, l' \in L$ 使得 $lxl' \in \operatorname{supp} f$, 则由刚才看到的, 知道 $f$ 在 $L(lxl')L = LxL$ 上为常数, 特别地 $f$ 在 $x$ 处不可能为零. 所以 $LxL \cap (\operatorname{supp} f) = \emptyset$, 特别地 $f$ 在 $LxL$ 上为常数 (等于 $0$).
现在证明 $f^{(B)}$ 被 $L \cap T$ 双不变; 这蕴含 $f^{(B)}$ 局部常值, 因为 $L \cap T$ 是 $T$ 的开紧子群. 由于 $T$ 交换, 只需证明 $f^{(B)}$ 被 $K \cap T$ 左不变, 但这立即由定义和 $\delta \colon T \to \R_{>0}$ 乘法 (因为 $\delta$ 必须在 $T$ 的紧子群 $K \cap T$ 上平凡) 这一事实得出.
最后, 需要证明 $f^{(B)}$ 有紧支撑. 设 $t \in T$ 使得 $f^{(B)}(t) \neq 0$. 则存在 $n \in N$ 使得 $tn \in B \cap (\operatorname{supp} f)$. 特别地, $t$ 在紧集 $B \cap (\operatorname{supp} f)$ 在连续投影映射 $B \to T \simeq B/N$ 下的像中. 这个像是紧的且包含 $f^{(B)}$ 的支撑, 所以 $f^{(B)}$ 有紧支撑.
若 $f \in \mathcal{H}$, 则可以在第一个断言的证明中取 $L = K$, 得到 $f^{(B)}$ 被 $K \cap T$ 双不变, 即 $f^{(B)} \in \mathcal{H}_T$.
(11) 设 $z \in (\CC^{\times})^3$. 注意乘法函数 $\delta$ 和 $\chi_z$ 都在 $N$ 上平凡. 所以若 $t \in T$, $n \in N$, 则 $\phi_z(tn) = \chi_z(t)\delta(t)^{1/2}$.
设 $f \in \mathcal{H}$. 则用问题 (2) 的积分公式, 得到
\begin{align*}
f^{\vee}(z) &= \int_T\int_N\int_K f(tnk)\phi_z(tnk)\,\dd t\,\dd n\,\dd k \\
&= \int_T\int_N f(tn)\phi_z(tn)\,\dd t\,\dd n \quad \text{(因为 } f \text{ 和 } \phi_z \text{ 都被 } K \text{ 右不变)} \\
&= \int_T\int_N f(tn)\chi_z(t)\delta(t)^{1/2}\,\dd t\,\dd n \\
&= \int_T f^{(B)}(t)\chi_z(t)\,\dd t \\
&= (f^{(B)})^{\vee}(z).
\end{align*}
(12) 实际上可以直接计算 (通过写出 $ntn^{-1}t^{-1}$ 的公式并用换元), 但采用提示, 因为它更好地推广到任意 $n$ 的 $\GL_n(\Q_p)$.
首先注意 $N$ 的子群 $U$ 正规 (甚至中心) 且同构于加法群 $\Q_p$. 特别地, 有由习题 \ref{ex:I.5.3.4} 给出的 $U$ 上的 Haar 测度 $\mu_U = \lambda$ (它由要求 $\Z_p$ 体积为 $1$ 规范化).
其次, 注意映射 $N \to \Q_p^2$, $\begin{pmatrix} 1 & a & b \\ 0 & 1 & c \\ 0 & 0 & 1 \end{pmatrix} \longmapsto (a, c)$ 诱导拓扑群同构 $N/U \xrightarrow{\sim} \Q_p^2$, 且 $N$ 在 $N/U$ 上通过左平移的作用在此同构下对应于 $N$ 在 $\Q_p^2$ 上的作用 $\begin{pmatrix} 1 & a & b \\ 0 & 1 & c \\ 0 & 0 & 1 \end{pmatrix} \cdot (x, y) = (a + x, c + y)$. 特别地, $\Q_p^2$ 上的测度 $\lambda^2$ 给出 $N/U$ 上被 $N$ 左不变的规则 Borel 测度 $d\mu_{N/U}$. 则正线性函数 $\mathcal{C}_c(N) \to \CC$, $f \longmapsto \int_{N/U}\int_U f(nu)\,\dd\mu_{N/U}(n)\,\dd\mu_U(u)$ 左不变, 所以它是 Haar 测度 $\mu_N$ 的倍数. 又, 取 $f = \mathbb{1}_{K \cap N}$ (并用 $K \cap U = \Z_p$ 以及 $(K \cap N)/(K \cap U)$ 在上述同构下对应于 $\Q_p^2$ 的子群 $\Z_p^2$ 这一事实), 看到这个正线性函数实际上等于 $\mu_N$, 这证明提示中的积分公式.
在下面的计算中, 把 $U$ 等同于 $\Q_p$, 把 $N/U$ 等同于 $\Q_p^2$. 设 $t = \begin{pmatrix} x_1 & 0 & 0 \\ 0 & x_2 & 0 \\ 0 & 0 & x_3 \end{pmatrix} \in T$. 注意映射 $N \to N$, $n \longmapsto ntn^{-1}t^{-1}$ 保持 $U$, 所以它诱导映射 $U \to U$ 和 $N/U \to N/U$. 若 $u \in U = \Q_p$, 有 $utu^{-1}t^{-1} = (1 - x_1x_3^{-1})u$. 另一方面, 若 $n \in N$, 且 $n$ 在 $N/U = \Q_p^2$ 中的像是 $(a, b)$, 则 $ntn^{-1}t^{-1}$ 在 $N/U$ 中的像是 $(a(1 - x_1x_2^{-1}), b(1 - x_2x_3^{-1}))$.
假设 $D(t) \neq 0$. 由习题 \ref{ex:I.5.3.8} 证明的换元公式, 有:
\begin{itemize}
\item 对每个 $f \in \mathcal{C}_c(U)$,
\[
\int_U f(u)\,\dd u = |1 - x_1x_3^{-1}|_p\int_U f(utu^{-1}t^{-1})\,\dd u;
\]
\item 对每个 $f \in \mathcal{C}_c(N/U)$,
\[
\int_{N/U}f(n)\,\dd n = |1 - x_1x_2^{-1}|_p|1 - x_2x_3^{-1}|_p\int_{N/U}f(ntn^{-1}t^{-1})\,\dd n.
\]
\end{itemize}
又, 用 $U$ 在 $N$ 中中心且 $t^{-1}Ut = U$ 这一事实, 看到对每个 $n \in N$ 和每个 $u \in U$,
\[
(ntn^{-1}t^{-1})(utu^{-1}t^{-1}) = ntn^{-1}(t^{-1}ut)u^{-1}t^{-1} = nt(t^{-1}ut)u^{-1}n^{-1}t^{-1} = nut(nu)^{-1}t^{-1}.
\]
把积分公式和上面的计算放在一起, 看到对每个 $f \in \mathcal{C}_c(N)$,
\begin{align*}
\int_N f(n)\,\dd n &= \int_{N/U}\int_U f(nu)\,\dd n\,\dd u \\
&= |1 - x_1x_2^{-1}|_p|1 - x_2x_3^{-1}|_p\int_{N/U}\int_U f(ntn^{-1}t^{-1}utu^{-1}t^{-1})\,\dd n\,\dd u \\
&= |1 - x_1x_2^{-1}|_p|1 - x_2x_3^{-1}|_p\int_{N/U}\int_U f(nut(nu)^{-1}t^{-1})\,\dd n\,\dd u.
\end{align*}
现在做代换 $(n, u) \longmapsto nu$ (这给出从 $N/U \times U$ 到 $N$ 的同胚), 得到
\[
\int_N f(n)\,\dd n = |1 - x_1x_2^{-1}|_p|1 - x_2x_3^{-1}|_p\int_N f(ntn^{-1}t^{-1})\,\dd n.
\]
这正是所要的, 因为
\[
D(t) = |1 - x_1x_2^{-1}|_p|1 - x_1x_3^{-1}|_p|1 - x_2x_3^{-1}|_p|1 - x_2x_1^{-1}|_p|1 - x_3x_1^{-1}|_p|1 - x_3x_2^{-1}|_p
\]
且 $|1 - x_ix_j^{-1}|_p = |1 - x_jx_i^{-1}|_p$ (因为 $|1 - x_ix_j^{-1}|_p = |x_jx_i^{-1}|_p|1 - x_jx_i^{-1}|_p = |1 - x_jx_i^{-1}|_p$), 所以 $D(t) = (|1 - x_1x_2^{-1}|_p|1 - x_2x_3^{-1}|_p)^2$, 且 $D(t)^{1/2} = |1 - x_1x_2^{-1}|_p|1 - x_2x_3^{-1}|_p$. 而 $\delta(t) = |x_1|_p^2|x_2|_p^0|x_3|_p^{-2}$, 所以 $\delta(t)^{1/2} = |x_1|_p|x_3|_p^{-1} = |1 - x_1x_3^{-1}|_p$ 在 $D(t) \neq 0$ 时由上一段最后一个等式等于 $|1 - x_1x_3^{-1}|_p$.
实际上, 上面的计算给出
\[
\int_N f(n)\,\dd n = D(t)^{1/2}\delta(t)^{1/2}\int_N f(ntn^{-1}t^{-1})\,\dd n
\]
对每个 $f \in \mathcal{C}_c(N)$, 如所愿.
(13) 设 $f \in \mathcal{H}$, $t \in T$ 满足 $D(t) \neq 0$. 回忆 $f^{(B)}(t) = \delta(t)^{1/2}\int_N f(tn)\,\dd n$. 用 (12) (把 $f$ 换成 $n \longmapsto f(tn)$ 并做代换), 得到
\begin{align*}
f^{(B)}(t) &= \delta(t)^{1/2}D(t)^{-1/2}\delta(t)^{-1/2}\int_N f(tntn^{-1}t^{-1})\,\dd n \\
&= D(t)^{-1/2}\int_N f(tntn^{-1}t^{-1})\,\dd n.
\end{align*}
另一方面, 需要把 $\mathcal{O}_t(f) = \int_{T\backslash G}f(g^{-1}tg)\,\dd\mu_{T\backslash G}$ 用积分公式 $\int_G = \int_{T\backslash G}\int_N\int_T$ 重写. 用 (2) 的分解 $G = TNK$ 和双不变性, 得到
\[
\mathcal{O}_t(f) = \int_N f(tntn^{-1}t^{-1})\,\dd n
\]
(这里我们用了 $\int_G f = \int_{T\backslash G}\int_N\int_T f(tnk)$, 而 $\mu_{T\backslash G}$ 与 $\mu_T \times \mu_N$ 的配合给出测量分解). 组合两式给出 $\mathcal{O}_t(f) = D(t)^{-1/2}f^{(B)}(t)$.
(14) 设 $f \in \mathcal{H}$. 由 (5), $(G, K)$ 是 Gelfand 对, 所以由定理 \ref{thm:V.4.4}, $\chi_{\phi_z} \colon \mathcal{H} \to \CC$ 是乘法泛函. 但由 (7), $\chi_{\phi_z}(f) = f^{\vee}(z)$, 所以对 $f_1, f_2 \in \mathcal{H}$, $(f_1 * f_2)^{\vee}(z) = f_1^{\vee}(z)f_2^{\vee}(z)$, 且 $\mathcal{H}$ 交换 (因为 $(G, K)$ 是 Gelfand 对). 现在 $\mathcal{H} = \mathcal{C}_c(K\backslash G/K)$, 且由 (10), 映射 $f \longmapsto f^{(B)}$ 是 $\mathcal{H}$ 到 $\mathcal{H}_T$ 的代数同态 (对卷积). 所以 $f \mapsto f^{(B)}$ 诱导 $\mathcal{H}$ 在 $\CC[z_1^{\pm 1}, z_2^{\pm 1}, z_3^{\pm 1}]$ 的表示, 而 (11) 说 $f^{\vee}(z) = (f^{(B)})^{\vee}(z)$. 由 (9), $(f^{(B)})^{\vee}$ 在 $\CC[z_1^{\pm 1}, z_2^{\pm 1}, z_3^{\pm 1}]$ 中. 为证明对称性: 由 (13), 对满足 $D(t) \neq 0$ 的 $t \in T$, $\mathcal{O}_t(f) = D(t)^{-1/2}f^{(B)}(t)$. 但 $\mathcal{O}_t(f) = \int_{T\backslash G}f(g^{-1}tg)$ 是对 $t$ 的共轭类积分, 且显然 $\mathcal{O}_t(f) = \mathcal{O}_{t'}(f)$ 若 $t$ 与 $t'$ 共轭. 因为 $f$ 双不变, 函数 $t \longmapsto f^{(B)}(t)$ 在 $T$ 的 $W = S_3$ (Weyl 群) 轨道上为常数. 因此 $f^{\vee}(z) = (f^{(B)})^{\vee}(z)$ 是对 $z$ 的坐标置换对称的, 即 $f^{\vee} \in \CC[z_1^{\pm 1}, z_2^{\pm 1}, z_3^{\pm 1}]^{S_3}$.
(15) 由 (4), $\mathcal{H}$ 作为向量空间由 $f_{\lambda} = \mathbb{1}_{Kp^{\lambda}K}$, $\lambda \in X^+$ 张成 (这些是双不变函数的自然基). 为证明线性无关: 若 $\sum_{\lambda}c_{\lambda}f_{\lambda} = 0$, 在 $p^{\mu}$ ($\mu \in X^+$) 处取值, 得到 $c_{\mu} = 0$ (因为 $p^{\mu} \in Kp^{\mu}K$ 且 $p^{\mu} \notin Kp^{\lambda}K$ 对 $\lambda \neq \mu$, 由 (4) 的不相交性).
(16) 由 (8) 和 (9), 更准确地说由 (9) 的证明, $\CC[z_1^{\pm 1}, z_2^{\pm 1}, z_3^{\pm 1}]$ 有基 $(z^{\lambda})_{\lambda \in X}$. $S_3$ 通过置换变量作用于其上, 且由定义 $c_{\lambda} = \sum_{\sigma \in S_3}z^{\sigma \cdot \lambda}$. 由于 $S_3$ 有限, 投影算子 $p \longmapsto \frac{1}{|S_3|}\sum_{\sigma \in S_3}\sigma \cdot p$ 把 $\CC[z_1^{\pm 1}, z_2^{\pm 1}, z_3^{\pm 1}]$ 映到其 $S_3$-不变部分, 且它是满射. 由 (17) 后面的标准论证 (用偏序和三角形论证) 或者更简单地: 由对称函数的基本定理 (对称多项式 $\Leftrightarrow$ 基本对称多项式的多项式), 或者直接论证: 把 $z^{\lambda}$ 在 $S_3$-轨道上的和 $c_{\lambda}$ 记为 $\lambda \in X^+$; 它们在偏序 $\leq$ 下形成一个三角形系 (每个 $c_{\lambda}$ 有首项 $z^{\lambda}$), 所以构成不变部分的基.
(17) 参考上面 (14) 和 (15) 以及 Cartier 的论证. 由于篇幅, 我们把细节留给读者并给出核心步骤:
(a) 证明 $\|g\| = \|kgk'\|$ 对 $k, k' \in K$ (因为 $K = \GL_3(\Z_p)$ 保持 sup 范数).
(b) $\Lambda_r$ 是 $r \times r$ 子行列式; 它们在 $K$ 的左、右乘法下不变 (因为子行列式的列/行组合在 $\Z_p$ 中, 范数 $\leq 1$), 给出 $(\alpha)$. 对 $(\beta)$: 对 $t \in T$, $n \in N$, 用行阶梯: $\Lambda_r(tn)$ 的每个子行列式是 $t$ 的某个主子行列式与 $N$-因子之积, 且 $N$-因子的 $p$-进绝对值 $\leq 1$, 所以 $\|\Lambda_r(tn)\| \geq \|\Lambda_r(t)\|$. 对 $(\gamma)$: 对 $g = p^{\lambda}$, $\lambda \in X^+$, 最大的 $r \times r$ 子行列式 (绝对值) 是主子式 $\prod_{i=3-r+1}^3 p^{\lambda_i}$, 给出 $\|\Lambda_r g\| = p^{-(\lambda_{3-r+1} + \cdots + \lambda_3)}$.
(c) 用 (13): $c_{\lambda}(\mu) = (f_{\lambda})^{\vee}$ 中 $z^{\mu}$ 的系数 $= f_{\lambda}^{(B)}(p^{\mu})$ (用 (11) 和 (9) 的基), 且 $f_{\lambda}^{(B)}(p^{\mu}) = f_{\lambda}^{(B)}(p^{\check{\mu}})$, 因为 $f_{\lambda}^{(B)}$ 在 $T$ 的 Weyl 群轨道上为常数 (由 (14) 的论证, 因为 $f_{\lambda}$ 双不变).
(d) 若 $p^{\mu}n \in Kp^{\lambda}K$, 则 $\|\Lambda_r(p^{\mu}n)\| = \|\Lambda_r(p^{\lambda})\|$ 对每个 $r$ (由 (a)-$(\gamma)$), 所以由 $(\beta)$, $\|\Lambda_r(p^{\mu})\| \geq \|\Lambda_r(p^{\lambda})\|$ 对每个 $r$. 由 $(\gamma)$, 这蕴含 $\sum_{i=3-r+1}^3\mu_i \leq \sum_{i=3-r+1}^3\lambda_i$ 对每个 $r$, 即 $\mu \leq \lambda$.
(e) (i) 由 (d) 和 (c) 得出: 若 $\mu \nleq \lambda$, 则不存在 $n \in N$ 使得 $p^{\mu}n \in Kp^{\lambda}K$; 但 $c_{\lambda}(\mu) = f_{\lambda}^{(B)}(p^{\mu}) = \delta(p^{\mu})^{1/2}\int_N f_{\lambda}(p^{\mu}n)\,\dd n = 0$, 因为 $f_{\lambda}$ 是 $Kp^{\lambda}K$ 的特征函数且 $p^{\mu}n \notin Kp^{\lambda}K$.
(f) 若 $tn \in KtK$ 且 $n \in N$, 则 $p^{\mu}n \in Kp^{\lambda}K$ (取 $\mu = \lambda$), 所以用 (e) 的论证的逆: $\|\Lambda_r(tn)\| = \|\Lambda_r(t)\|$ 对每个 $r$; 由 $(\beta)$ 的等号情形 (需要 $n$ 的因子都已在 $\Z_p^{\times}$ 中), 得到 $n \in N \cap K$. 反之显然.
(g) (ii) 由 (c), (f) 和 (2) 的积分公式得出: $c_{\lambda}(\lambda) = f_{\lambda}^{(B)}(p^{\lambda}) = \delta(p^{\lambda})^{1/2}\int_N f_{\lambda}(p^{\lambda}n)\,\dd n = \delta(p^{\lambda})^{1/2}\operatorname{vol}(N \cap K) = \delta(p^{\lambda})^{1/2}$, 因为 $f_{\lambda}$ 是 $Kp^{\lambda}K$ 的特征函数且 $p^{\lambda}n \in Kp^{\lambda}K$ 当且仅当 $n \in N \cap K$ (由 (f)), 而 $\operatorname{vol}(N \cap K) = 1$. 由 (1) 中 $\delta$ 的定义, $\delta(p^{\lambda})^{1/2} = \prod_{i=1}^3|p^{\lambda_i}|_p^{(3+1-2i)/2} = p^{\frac{1}{2}\sum_{i=1}^3(4-2i)\lambda_i}$, 即所要的公式.
(18) 把 $f \longmapsto f^{\vee}$ 记为 $\mathcal{S}$. 在 (14) 中看到它良定义, 在 (7) 中看到它是 $\CC$-代数态射. 所以只需证明它是 $\CC$-向量空间的同构. 在 (15) 中给出 $\mathcal{H}$ 的基 $(f_{\lambda})_{\lambda \in X^+}$, 在 (16) 中给出 $\CC[z_1^{\pm 1}, z_2^{\pm 1}, z_3^{\pm 1}]^{S_3}$ 的基 $(c_{\lambda})_{\lambda \in X^+}$. 又, 集合 $X^+$ 有偏序, $\mathcal{S}$ 在两个基中的矩阵由 (17)(i) 对这个偏序上三角, 由 (17)(ii) 有非零对角线元素. 所以 $\mathcal{S}$ 可逆.
若想把论证做得更显式, 可以这样说: 由 (17)(i) 和 (ii), 对每个 $\lambda \in X^+$,
\[
\mathcal{S}(f_{\lambda}) = \sum_{\mu \in X^+, \mu \leq \lambda}c_{\lambda}(\mu)c_{\mu},
\]
其中 $c_{\lambda}(\lambda) \in \CC^{\times}$. 递归构造 $\mathcal{S}$ 的逆 $T \colon \CC[z_1^{\pm 1}, z_2^{\pm 1}, z_3^{\pm 1}]^{S_3} \to \mathcal{H}$. 只需给出 $T$ 在基元素 $c_{\lambda}$, $\lambda \in X^+$ 上的值. 若不存在满足 $\mu \leq \lambda$ 且 $\mu \neq \lambda$ 的 $\mu \in X^+$, 设 $T(c_{\lambda}) = c_{\lambda}(\lambda)^{-1}f_{\lambda}$. 否则, 递归定义 $T(c_{\lambda})$:
\[
T(c_{\lambda}) = c_{\lambda}(\lambda)^{-1}f_{\lambda} - c_{\lambda}(\lambda)^{-1}\sum_{\mu \in X^+ - \{\lambda\}, \mu \leq \lambda}c_{\lambda}(\mu)T(c_{\mu}).
\]
由于 $\{\mu \in X^+ \mid \mu \leq \lambda\}$ 有限 (事实上, 若 $\mu \in X^+$ 且 $\mu \leq \lambda$, 则 $\mu_3 = (\lambda_1 + \lambda_2 + \lambda_3) - (\mu_1 + \mu_2) \geq \lambda_3$, 所以 $\lambda_3 \leq \mu_3 \leq \mu_2 \leq \mu_1 \leq \lambda_1$), 这个过程总会终止.
\end{solution}

\subsection{问题}
\label{sec:V.8.4}
\footnote{这个结果叫什么名字? 有参考文献吗?}\footnote{归功于 H. T. G. Dospinescu.}
在此问题中, 写 $G = \SL_2(\R)$, $K = \SO(2)$. 记 $B$ 为 $G$ 中上三角矩阵的子群, $N$ 为幂幺上三角矩阵 (即对角线元素全为 $1$ 的上三角矩阵) 的子群.\footnote{本问题的一切都对 $\SL_n(\R)$ 成立, 但推广比问题 V.8.3 更麻烦.}\footnote{参考文献.}
在所有群上固定左 Haar 测度. 回忆 $G$ 幺模.
\begin{enumerate}
\item[(1)] 此题的目标是证明如下事实 $(*)$: 若 $(\pi, V)$ 是 $G$ 的酉表示, $v \in V$ 被 $G$ 的形如 $\begin{pmatrix} 1 & u \\ 0 & 1 \end{pmatrix}$ ($u \neq 0$) 的元素固定, 则 $v$ 被 $G$ 的每个元素固定.
我们固定 $G$ 的酉表示 $(\pi, V)$.
\begin{enumerate}
\item[(a)] 设 $v \in V$, 假设存在连续群态射 $\chi \colon \R \to S^1$ 使得对每个 $u \in \R$ 有 $\pi\begin{pmatrix} 1 & u \\ 0 & 1 \end{pmatrix}(v) = \chi(u)v$. 定义 $\phi \colon G \to \CC$: $\phi(x) = \langle \pi(x)(v), v \rangle$.
  \begin{enumerate}
  \item[(i)] 证明若 $x = \begin{pmatrix} a & b \\ c & d \end{pmatrix}$, $x' = \begin{pmatrix} a' & b' \\ c' & d' \end{pmatrix}$ 满足 $c = c'$, 则 $|\phi(x)| = |\phi(x')|$.
  \item[(ii)] 证明若 $x \in B$, 则 $|\phi(x)| = |\phi(1)|$.
  \item[(iii)] 证明存在连续群态射 $\psi \colon B \to S^1$ 使得对每个 $x \in B$, $\pi(x)(v) = \psi(x)v$.
  \item[(iv)] 证明函数 $x \longmapsto |\phi(x)|$ 在 $G$ 上为常数.
  \item[(v)] 证明对每个 $x \in G$ 有 $\pi(x)(v) = v$.
  \end{enumerate}
\item[(b)] 证明 $(*)$.
\end{enumerate}
\item[(2)] 设
\[
A^+ = \left\{\begin{pmatrix} a & 0 \\ 0 & a^{-1} \end{pmatrix}, a \in \R_{\geq 1}\right\} \subset G.
\]
\begin{enumerate}
\item[(a)] 证明
\[
G = \bigcup_{a \in A^+}KaK.
\]
(作为集合, 不作为拓扑空间.)
\item[(b)] 证明 $(G, K)$ 是 Gelfand 对.
\item[(c)] 把问题 (a) 和 (b) 推广到 $\SL_n(\R)$, 其中 $K = \SO(n)$, $A^+$ 是满足 $a_1 \geq a_2 \geq \cdots \geq a_n > 0$ 的对角线系数 $a_1, \ldots, a_n$ 的对角矩阵的集合.
\end{enumerate}
\item[(3)] 此题的目标是证明如下事实 $(**)$: 对满足 $V^G = \{0\}$ 的每个酉表示 $(\pi, V)$, 对所有 $v, w \in V$, $\lim_{x \to \infty}\langle \pi(x)(v), w \rangle = 0$, 其中 $\infty$ 是 $G$ 的 Alexandroff 紧化的无穷远点. (所以 $\lim_{x \to \infty}\langle \pi(x)(v), w \rangle = 0$ 意味着: 对每个 $\varepsilon > 0$, 存在 $G$ 的紧子集 $X$ 使得对每个 $x \in G - X$, $|\langle \pi(x)(v), w \rangle| \leq \varepsilon$.)
我们固定 $G$ 的酉表示 $(\pi, V)$, 假设它不满足 $(**)$ 的结论.
\begin{enumerate}
\item[(a)] 证明存在 $v, v' \in V$, $\R_{\geq 1}$ 的元素序列 $(a_n)_{n \geq 0}$ 和 $\alpha \in \CC - \{0\}$ 使得:
\begin{itemize}
\item 当 $n \to +\infty$ 时 $a_n \to +\infty$;
\item 当 $n \to +\infty$ 时 $\langle \pi(t_n)(v), v' \rangle \to \alpha$, 其中 $t_n = \begin{pmatrix} a_n & 0 \\ 0 & a_n^{-1} \end{pmatrix}$.
\end{itemize}
\item[(b)] 证明取 $(a_n)_{n \geq 0}$ 的子序列后, 可以假设存在 $v_0 \in V$ 使得对每个 $w \in V$, 当 $n \to +\infty$ 时 $\langle \pi(t_n)(v), w \rangle \to \langle v_0, w \rangle$.
\item[(c)] 证明 $v_0 \neq 0$ 且 $v_0$ 被 $N$ 的每个元素固定. (提示: 若 $x \in N$, $w \in V$, 当 $n \to +\infty$ 时 $\langle \pi(xt_n)(v), w \rangle - \langle \pi(t_n)(v), w \rangle \to 0$ 的行为如何?)
\item[(d)] 得出结论.
\end{enumerate}
\item[(4)] 此题的目标是证明商 $\SL_2(\Z)\backslash G$ 有有限体积.
设 $\mathfrak{h} := \{z \in \CC \mid \Im(z) > 0\}$.
对 $g = \begin{pmatrix} a & b \\ c & d \end{pmatrix} \in G$ 和 $z \in \mathfrak{h}$, 令
\[
g \cdot z = \frac{az + b}{cz + d}.
\]
\begin{enumerate}
\item[(a)] 证明这定义 $G$ 在 $\mathfrak{h}$ 上的左作用.
\item[(b)] 证明 $i \in \mathfrak{h}$ 的稳定化子是 $K$, 且这诱导同胚 $G/K \xrightarrow{\sim} \mathfrak{h}$.
\item[(c)] 证明 (b) 的同胚把测度 $\mu_{G/K}$ 送到 $\frac{dx\,dy}{y^2}$ 的倍数, 其中 $x$ 和 $y$ 是 $z \in \mathfrak{h}$ 的实部和虚部, $dx$ 和 $dy$ 是 $\R$ 上的 Lebesgue 测度.
\item[(d)] 设 $\Omega = \{z \in \mathfrak{h} \mid |\Re(z)| \leq \frac{1}{2} \text{ 且 } |z| \geq 1\}$. 证明对每个 $z \in \mathfrak{h}$, 存在 $g \in \SL_2(\Z)$ 使得 $g \cdot z \in \Omega$.
\item[(e)] 证明 $\int_{\Omega}\frac{dx\,dy}{y^2} < +\infty$.
\item[(f)] 证明 $\operatorname{vol}(\SL_2(\Z)\backslash G) < +\infty$.
\end{enumerate}
\item[(5)] 在此题中, 固定 $G$ 的闭离散子群 $\Gamma$, 使得 $\operatorname{vol}(\Gamma\backslash G) < +\infty$. 注意由习题 \ref{ex:I.5.3.7}, $G$ 幺模, 所以 $G$ 上的 Haar 测度也右不变. 写 $X = \Gamma\backslash G$, 记 $\mu$ 为测度 $\mu_{\Gamma\backslash G}$; 规范化 $G$ 和 $\Gamma$ 上的 Haar 测度使得 $\mu(X) = 1$. 考虑空间 $L^2(X) = L^2(X, \mu)$, 带 $G$ 的右拟正则表示 $\pi$ (即对 $g \in G$, $f \in L^2(X)$, $\pi(g)(f) = R_gf$).
证明对所有 $f_1, f_2 \in L^2(X)$,
\[
\lim_{g \to \infty}\int_X f_1(xg)f_2(x)\,\dd\mu(x) = \int_X f_1(x)\,\dd\mu(x)\int_X f_2(x)\,\dd\mu(x).
\]
(提示: $V := \{f \in L^2(X) \mid \int_X f(x)\,\dd\mu(x) = 0\}$ 中 $G$-不变向量是什么?)
\item[(6)] 用问题 (5) 的记号. 设 $Y = (\Gamma \cap K)\backslash K$. 在 $Y$ 上放测度 $\mu_Y = \mu_{(\Gamma\cap K)\backslash K}$, 规范化 $\Gamma \cap K$ 和 $K$ 上的 Haar 测度使得 $Y$ 体积为 $1$.
此题的目标是证明如下事实 $(***)$: 对每个 $f \in \mathcal{C}_c(X)$,
\[
\lim_{g \to \infty}\int_Y f(yg)\,\dd\mu_Y(y) = \int_X f(x)\,\dd\mu(x).
\]
\footnote{换句话说, 当 $g \to \infty$ ($G$ 中) 时集合 $Yg$ 在 $X$ 中变得等分布. 注意 $g \to \infty$ 当且仅当 $g^{-1} \to \infty$.}
把 $\mu_K$ 看作 $\mathcal{C}_c(G)$ 上的连续线性泛函:
\[
\mu_K(f) = \int_K f(x)\,\dd\mu_K(x),
\]
对 $f \in \mathcal{C}_c(G)$. 如果 $\mu_K$ 可以由 $\mathcal{C}_c(G)$ 的元素表示 (即如果有 $h \in \mathcal{C}_c(G)$ 使得对每个 $f \in \mathcal{C}_c(G)$, $\int_K f(x)\,\dd\mu_K(x) = \int_G f(x)h(x)\,\dd\mu_G(x)$), 那么证明 $(***)$ 会相对容易 (从已经做过的). 但事实并非如此. 不过, 可以尝试用 $\mathcal{C}_c(G)$ 的元素逼近 $\mu_K$. 你可能记得我们已经多次用过这种技巧来逼近 Dirac 测度.
\begin{enumerate}
\item[(a)] 若 $\psi \in \mathcal{C}_c(G)$, 定义 $\mathcal{C}_c(G)$ 上的连续线性泛函 $\psi * \mu_K$:
\[
\psi * \mu_K(f) = \int_{G \times K}f(xy)\psi(x)\,\dd\mu_G(x)\,\dd\mu_K(y).
\]
证明存在 $h \in \mathcal{C}_c(G)$ 使得对每个 $f \in \mathcal{C}_c(G)$,
\[
\psi * \mu_K(f) = \int_G f(x)h(x)\,\dd\mu_G(x).
\]
\item[(b)] 证明存在 $\mathcal{C}_c(G)$ 的元素序列 $(\psi_n)_{n \geq 0}$ 使得 $\mu_K$ 是序列 $(\psi_n * \mu_K)_{n \geq 0}$ 在 $\Hom(\mathcal{C}_c(G), \CC)$ 的弱*拓扑中的极限.
\item[(c)] 证明 $(***)$.
\end{enumerate}
\end{enumerate}

\begin{solution}
(1)(a). 注意若 $n = \begin{pmatrix} 1 & u \\ 0 & 1 \end{pmatrix}$, $n' = \begin{pmatrix} 1 & u' \\ 0 & 1 \end{pmatrix} \in N$, $x \in G$, 则
\[
|\phi(nxn')| = |\langle \pi(x)(\pi(n')(v)), \pi(n)^{-1}(v) \rangle| = |\chi(u')\chi(u)\langle \pi(x)(v), v \rangle| = |\phi(x)|.
\]
(i) 首先假设 $c = c' \neq 0$. 设 $n = \begin{pmatrix} 1 & c^{-1}(d' - d) \\ 0 & 1 \end{pmatrix}$, $n' = \begin{pmatrix} 1 & c^{-1}(a' - a) \\ 0 & 1 \end{pmatrix}$. 则
\[
n'xn = \begin{pmatrix} a' & w \\ c' & d' \end{pmatrix},
\]
对某个 $w \in \R$. 由于 $1 = \det(n'xn) = a'd' - wc' = a'd' - b'c'$ 且 $c' \neq 0$, 必有 $w = b'$, 所以 $n'xn = x'$, 且上面的注记蕴含 $|\phi(x)| = |\phi(x')|$.
现在假设 $c = c' = 0$. 则可以找到 $G$ 中的序列 $(x_n)_{n \geq 1}$ 和 $(x'_n)_{n \geq 1}$ 使得 $x_n \to x$, $x'_n \to x'$, 且对每个 $n \geq 1$, $x_n$ 和 $x'_n$ 的 $(2,1)$ 项都等于 $\frac{1}{n}$. 例如可以取 $x_n = \begin{pmatrix} a + nd & b \\ \frac{1}{n} & d \end{pmatrix}$, 这有意义因为 $\det(x) = ad = 1$ 蕴含 $d \neq 0$, 用类似的公式定义 $x'_n$. 把第一种情形应用于 $x_n$ 和 $x'_n$ 给出对每个 $n \geq 1$ 有 $|\phi(x_n)| = |\phi(x'_n)|$; 由于 $\phi$ 连续, 可以在 $n$ 中取极限, 得到 $|\phi(x)| = |\phi(x')|$.
(ii) 这由问题 (i) 的``$c = c' = 0$''情形得出.
(iii) 设 $x \in B$. 由 (ii), $|\langle \pi(x)(v), v \rangle| = |\langle v, v \rangle|$, 所以由 Cauchy-Schwarz 不等式的等号情形, 必有 $\pi(x)(v) = \psi(x)v$, 其中 $\psi(x) \in S^1$. 由于 $x \longmapsto \pi(x)(v)$ 对 $x$ 连续且 $\psi(x) = \langle \pi(x)(v), v \rangle$, 函数 $x \longmapsto \psi(x)$ 也连续. 它群态射这一事实立即由对所有 $x, y \in B$ 有 $\pi(xy) = \pi(x) \circ \pi(y)$ 得出.
(iv) 重复 (a) 的证明开头所做的, 用 (iii), 看到对每个 $x \in G$ 和所有 $b, b' \in G$, $|\phi(bxb')| = |\phi(x)|$. 设 $x_0 = \begin{pmatrix} 0 & 1 \\ -1 & 0 \end{pmatrix}$. 设 $x = \begin{pmatrix} a & b \\ c & d \end{pmatrix} \in G$. 若 $c \neq 0$, 则
\[
x = \begin{pmatrix} 1 & c^{-1}a \\ 0 & 1 \end{pmatrix}x_0\begin{pmatrix} c & -d \\ 0 & b - c^{-1}ad \end{pmatrix},
\]
所以 $|\phi(x)| = |\phi(x_0)|$. 若 $c = 0$, 则可以用 $(2,1)$ 项非零的 $G$ 的元素逼近 $x$ (如在 (i) 的证明中), 所以由第一种情形和 $\phi$ 的连续性, 也得到 $|\phi(x)| = |\phi(x_0)|$.
(v) 重复 (iii) 的推理, 证明存在连续群态射 $\psi \colon G \to S^1$ (延拓 (iii) 的映射 $\psi$) 使得对每个 $x \in G$ 有 $\pi(x)(v) = \psi(x)v$. 但不存在从 $G$ 到交换群的非平凡群态射 (因为 $G$ 等于其换位子群, 而到交换群的群态射在每个换位子上平凡), 所以 $\psi = 1$, 因此 $v \in V^G$.
(b). 记 $\pi^0$ 为 $\R$ 在 $V$ 上的酉表示 $\pi^0(t) = \pi\begin{pmatrix} 1 & t \\ 0 & 1 \end{pmatrix}$, 对每个 $t \in \R$. 由假设 $\pi^0(u)(v) = v$. 由于 $\R$ 交换, 这蕴含 $\pi^0(t)(w) = w$ 对每个 $w \in \{\pi(t)(v) \mid t \in \R\}$ 成立, 从而 (因为两边对 $w$ 都线性且连续) 对每个 $w \in W := \overline{\operatorname{Span}\{\pi(t)(v) \mid t \in \R\}}$ 成立. 换句话说, $\R$ 在 $W$ 上的酉表示 $\pi^0|_W$ 在子群 $\Z u$ 上平凡, 所以它诱导紧群 $\R/\Z u$ 在 $W$ 上的酉表示. 由定理 \ref{thm:IV.2.1} 和推论 \ref{cor:I.3.4.4}, 空间 $W$ 是 $\R/\Z u$ 的 $1$ 维表示的直和的闭包, 即存在线性无关族 $(v_i)_{i \in I}$ 生成 $W$ 的稠密子空间, 且每个 $v_i$ 是所有 $\pi(n)$, $n \in N$ 的特征向量. 由问题 (a), 这蕴含所有 $v_i$ 都在 $V^G$ 中, 所以 $W \subset V^G$, 特别地 $v \in V^G$.
(2)(a). 直接对 $\SL_n(\R)$ 证明结果, 所以写 $G = \SL_n(\R)$, $K = \SO(n)$.
对 $\GL_n(\R)$ 的元素用极分解 (见 Rudin {[}20{]} 的定理 12.35\footnote{更好的参考文献.}, 虽然这对我们的目的来说是杀鸡用牛刀, 因为我们只需要有限维情形): 对每个 $g \in \GL_n(\R)$, 存在唯一偶对 $(u, p)$ 使得 $u$ 正交, $p$ 对称正定且 $g = up$; 又, $p^2 = g^Tg$. 若 $g \in G$, 则 $\det(p)^2 = \det(g^Tg) = 1$; 由于 $p$ 正定, 有 $\det(p) > 0$, 所以 $\det(p) = 1$, 然后也推出 $\det(u) = 1$, 即 $u \in K$.
证明 $G = \bigcup_{a \in A^+}KaK$. 设 $g \in G$. 设 $g = up$ 是 $g$ 的极分解. 由于 $p$ 对称, 由谱定理它可以在标准正交基中对角化, 所以可以把 $p$ 写成 $p = kak^{-1}$, 其中 $k \in \Or(n)$, $a$ 对角. $a$ 的对角线元素是 $p$ 的特征值, 可以假设它们降序排列, 所以 $a \in A^+$. 又, 若 $k$ 不在 $\SO(n)$ 中, 则改变特征基的一个向量的符号; 这不影响 $a$, 新的 $k$ 在 $\SO(n)$ 中. 最后得到 $g = (uk)ak^{-1} \in KaK$.
还要证明集合 $KaK$, $a \in A^+$, 两两不相交. 设 $a, a' \in A^+$, 假设 $KaK \cap Ka'K \neq \emptyset$. 则存在 $k, l \in K$ 使得 $a' = kal$. 所以
\[
(a')^2 = (a')^Ta' = k^Ta^Tk^Tkak^T = ka^2k^{-1},
\]
这蕴含 $(a')^2$ 和 $a^2$ 有相同的特征值, 即它们有相同的对角线元素 (可能重排). 由于 $(a')^2$ 和 $a^2$ 的对角线元素递减, 这两个矩阵实际上有相同的对角线元素, 即 $(a')^2 = a^2$. 最后, 由于 $a$ 和 $a'$ 是对角矩阵且元素为正, $(a')^2 = a^2$ 这一事实蕴含 $a' = a$.
(b). (仍取 $G = \SL_n(\R)$, $K = \SO(n)$.) 考虑连续群自同构 $\theta \colon G \to G$, $g \longmapsto (g^T)^{-1}$. 有 $\theta^2 = \mathrm{id}_G$, 且 $\theta|_K = \mathrm{id}_K$. 对每个 $g \in G$, 若写 $g = kal$, 其中 $k, l \in K$, $a \in A^+$, 则 $\theta(g) = ka^{-1}l = klg^{-1}kl \in Kg^{-1}K$. 由命题 \ref{prop:V.2.5}, $(G, K)$ 是 Gelfand 对.
(3) 现在回到 $\SL_2(\R)$, 所以释放字母 $n$.
(a). 由于 $V$ 不满足 $(**)$ 的结论, 存在 $v_1, v_2 \in V$ 和 $\varepsilon > 0$ 使得对 $G$ 的每个紧子集 $X$, 存在 $x \in G - X$ 使得 $|\langle \pi(x)(v_1), v_2 \rangle| > \varepsilon$.
对每个 $n \geq 1$, 设 $A_n^+$ 是满足 $a \leq n$ 的 $\begin{pmatrix} a & 0 \\ 0 & a^{-1} \end{pmatrix} \in A^+$ 的集合. 则 $\bigcup_{t \in A_n^+}KtK$ 在 $G$ 中相对紧, 因为若 $g \in KtK$, $t \in A_n^+$, 则 $g$ 和 $g^{-1}$ 的所有元素在绝对值上都以 $4n$ 为界 (注意若 $k \in K$, 则 $k$ 的所有元素在绝对值上都以 $1$ 为界). 所以可以找到 $x_n \in G - \bigcup_{t \in A_n^+}KtK$ 使得 $|\langle \pi(x_n)(v_1), v_2 \rangle| > \varepsilon$. 写 $x_n = k_nt_nl_n$, 其中 $k_n, l_n \in K$, $t_n = \begin{pmatrix} a_n & 0 \\ 0 & a_n^{-1} \end{pmatrix} \in A_n^+$. 由于 $K$ 紧, 取 $(k_n)$ 和 $(l_n)$ 的子序列后, 可以假设 $(k_n)$ (相应地 $(l_n)$) 有极限 $k \in K$ (相应地 $l \in L$). 设 $v = \pi(l)(v_1)$, $v' = \pi(k)^{-1}(v_2)$. 由于 $\CC$ 的序列 $(\langle \pi(t_n)(v), v' \rangle)$ 有界 (被 $\|v\|\|v'\|$), 取子序列后, 可以假设它收敛到某个 $\alpha \in \CC$. 为完成证明, 只需证明 $\alpha \neq 0$. 注意对每个 $n \geq 1$,
\[
\langle \pi(t_n)v, v' \rangle = \langle \pi(kk_n^{-1})\pi(k_nt_nl_n)\pi(l_n^{-1}l)v_1, v_2 \rangle = \langle \pi(x_n)\pi(l_n^{-1}l)v_1, \pi(k_nk^{-1})v_2 \rangle,
\]
所以
\[
|\langle \pi(t_n)v, v' \rangle - \langle \pi(x_n)v_1, v_2 \rangle| \leq \|v_1\|\|v_2\|(\|\pi(l_n^{-1}l) - \mathrm{id}\|_{\mathrm{op}} + \|\pi(k_nk^{-1}) - \mathrm{id}\|_{\mathrm{op}}).
\]
当 $n \to \infty$ 时这趋于 $0$. 由于对每个 $n \geq 1$ 有 $|\langle \pi(x_n)v_1, v_2 \rangle| \geq \varepsilon$, 序列 $(\langle \pi(t_n)(v), v' \rangle)$ 的极限 $\alpha$ 不能是 $0$.
(b). 结论意味着 $V$ 上的有界线性泛函序列 $(\Lambda_n \colon w \longmapsto \langle \pi(t_n)(v), w \rangle)_{n \geq 1}$ 在 $V^*$ 的弱*拓扑中收敛到有界线性泛函 $w \longmapsto \langle v_0, w \rangle$. 由于每个 $\Lambda_n$ 在 $V^*$ 的闭单位球中 (对算子范数), 且由 Banach-Alaoglu 定理这个闭单位球弱*紧, 取子序列后, 确实可以假设序列 $(\Lambda_n)_{n \geq 1}$ 在弱*拓扑中收敛到某个 $\Lambda \in V^*$. 但知道 $\Lambda$ 必具有 $w \longmapsto \langle v_0, w \rangle$ 的形式, 其中 $v_0 \in V$ 唯一确定, 所以我们完成了.
(c). 有 $\langle v_0, v' \rangle = \lim_{n \to +\infty}\langle \pi(t_n)(v), v' \rangle = \alpha \neq 0$, 所以 $v_0 \neq 0$.
设 $x \in N$, 写 $x = \begin{pmatrix} 1 & u \\ 0 & 1 \end{pmatrix}$. 设 $w \in V$. 则
\[
\langle \pi(xt_n)(v), w \rangle - \langle \pi(t_n)(v), w \rangle = \langle \pi(t_n)(\pi(t_n^{-1}xt_n)(v) - v), w \rangle,
\]
所以
\[
|\langle \pi(xt_n)(v), w \rangle - \langle \pi(t_n)(v), w \rangle| \leq \|(\pi(t_n^{-1}xt_n) - \pi(1))(v)\|\|w\|.
\]
由于 $t_n^{-1}xt_n = \begin{pmatrix} 1 & a_n^{-2}u \\ 0 & 1 \end{pmatrix}$ 且 $\lim_{n \to +\infty}a_n = +\infty$, 有 $\lim_{n \to +\infty}t_n^{-1}xt_n = 1$. 用表示 $\pi$ 的连续性, 推出 $\lim_{n \to +\infty}(\pi(t_n^{-1}xt_n)(v) - v) = 0$, 所以 $\lim_{n \to +\infty}|\langle \pi(xt_n)(v), w \rangle - \langle \pi(t_n)(v), w \rangle| = 0$. 但知道 $\lim_{n \to +\infty}\langle \pi(t_n)(v), w \rangle = \langle v_0, w \rangle$, 所以得到 $\lim_{n \to +\infty}\langle \pi(xt_n)(v), w \rangle = \langle v_0, w \rangle$. 另一方面,
\[
\lim_{n \to +\infty}\langle \pi(xt_n)(v), w \rangle = \lim_{n \to +\infty}\langle \pi(t_n)(v), \pi(x)^{-1}(w) \rangle = \langle v_0, \pi(x)^{-1}(w) \rangle = \langle \pi(x)(v_0), w \rangle.
\]
由于这对每个 $w \in V$ 成立, 推出 $\pi(x)(v_0) = v_0$.
(d). 由问题 (1), (c) 蕴含 $v_0$ 被 $G$ 的每个元素固定. 所以 $V^G \neq \{0\}$, 这完成了 $(**)$ 的逆否命题的证明.
(4)(a). 首先验证 $g \cdot z$ 良定义且在 $\mathfrak{h}$ 中 (若 $z \in \mathfrak{h}$). 写 $g = \begin{pmatrix} a & b \\ c & d \end{pmatrix}$. 若 $c \neq 0$, 则 $\Im(cz + d) = c\,\Im(z) \neq 0$, 所以商 $\frac{az+b}{cz+d}$ 有意义. 若 $c = 0$, 则 $ad = \det(g) = 1$, 所以 $d \neq 0$, 故 $cz + d = d \neq 0$, 因此商 $\frac{az+b}{cz+d}$ 有意义. 计算这个商的虚部. 有
\[
\frac{az + b}{cz + d} = \frac{(az + b)(cz + d)}{(cz + d)(cz + d)} = \frac{ac|z|^2 + bd + adz + bc\bar{z}}{(cz + d)(cz + d)},
\]
所以
\[
\Im\Bigl(\frac{az + b}{cz + d}\Bigr) = \frac{1}{(cz + d)(cz + d)}(ac\,\Im(z) - bd\,\Im(z)) = \frac{\Im(z)}{(cz + d)(cz + d)} > 0,
\]
且 $\frac{az+b}{cz+d} \in \mathfrak{h}$.
若 $g = 1$, 显然 $g \cdot z = z$. 设 $h = \begin{pmatrix} a' & b' \\ c' & d' \end{pmatrix}$ 是 $G$ 的另一个元素. 必须验证 $(gh) \cdot z = g \cdot (h \cdot z)$. 这是直接计算.
(b). 设 $g = \begin{pmatrix} a & b \\ c & d \end{pmatrix} \in G$ 使得 $g \cdot i = i$. 这等价于 $ai + b = i(ci + d) = -c + id$, 即 $a = d$ 且 $b = -c$. 若 $a = d$ 且 $b = -c$, 则 $1 = ad - bc = a^2 + b^2 = d^2 + c^2$, 且 $ac + bd = 0$, 所以 $g$ 是正交矩阵, 因为 $\det(g) = 1$ 而在 $K$ 中. 反之, 若 $g \in K$, 则 $g$ 具有 $\begin{pmatrix} \cos\theta & \sin\theta \\ -\sin\theta & \cos\theta \end{pmatrix}$ 的形式, 所以确实有 $a = d$ 和 $b = -c$, 因此 $g \cdot i = i$.
若 $z$ 是 $\mathfrak{h}$ 的任意元素, 写 $z = x + iy$, $x \in \R$, $y \in \R_{>0}$. 则 $z = g(z) \cdot i$, 其中 $g(z) = \begin{pmatrix} \sqrt{y} & \sqrt{y}^{-1}x \\ 0 & \sqrt{y}^{-1} \end{pmatrix} \in G$. 所以 $G$ 在 $\mathfrak{h}$ 上的作用传递, 且诱导连续双射 $\alpha \colon G/K \xrightarrow{\sim} \mathfrak{h}$. 为证明 $\alpha$ 是同胚, 注意刚才构造的映射 $\mathfrak{h} \to G$, $z \longmapsto g(z)$ 连续, 且 $z \longmapsto g(z)K$ 是 $\alpha$ 的逆.
(c). 设 $\mu$ 是 $\frac{dx\,dy}{y^2}$ 在 $\alpha^{-1}$ 下的原像. 像习题 \ref{sec:V.8.3} 的 (2) 的证明那样推理, 看到只需证明 $\mu$ 被 $G$ 左不变. 换句话说, 只需证明测度 $\frac{dx\,dy}{y^2}$ 被 $G$ 在 $\mathfrak{h}$ 上的作用不变.
设 $g = \begin{pmatrix} a & b \\ c & d \end{pmatrix} \in G$, $z = x + iy \in \mathfrak{h}$. 在 (a) 的证明中看到
\[
\Im(g \cdot z) = \frac{y}{(cz + d)(cz + d)} = \frac{y}{c^2(x^2 + y^2) + 2cdx + d^2}.
\]
类似地,
\[
\Re(g \cdot z) = \frac{ac(x^2 + y^2) + (ad + bc)x + bd}{c^2(x^2 + y^2) + 2cdx + d^2}.
\]
若写 $x' = \Re(g \cdot z)$, $y' = \Im(g \cdot z)$, 要证明 $\frac{dx'\,dy'}{(y')^2} = \frac{dx\,dy}{y^2}$.
设 $g_0 = \begin{pmatrix} 0 & 1 \\ -1 & 0 \end{pmatrix}$. 在 (1)(a)(iv) 的证明中看到 $G = B \cup Bg_0B$, 所以只需证明 $\frac{dx\,dy}{y^2}$ 被 $B$ 的元素和 $g_0$ 的作用不变. 两种情形都用 $\R^2$ 上通常的换元计算 $dx'\,dy'$.
若 $g \in B$ (即 $c = 0$), 则 $x' = \frac{1}{d^2}(x + bd)$, $y' = \frac{y}{d^2}$, 所以 $dx'\,dy' = \frac{1}{d^4}dx\,dy$, 且得到 $\frac{dx'\,dy'}{(y')^2} = \frac{dx\,dy}{y^2}$.
若 $g = g_0$, 则 $x' = -\frac{x}{x^2 + y^2}$, $y' = \frac{y}{x^2 + y^2}$, 所以
\[
dx'\,dy' = \det\begin{pmatrix} \frac{-(x^2+y^2)+2x^2}{(x^2+y^2)^2} & \frac{-2xy}{(x^2+y^2)^2} \\ \frac{2xy}{(x^2+y^2)^2} & \frac{(x^2+y^2)-2y^2}{(x^2+y^2)^2} \end{pmatrix}dx\,dy = \frac{dx\,dy}{(x^2 + y^2)^2},
\]
且再次这立即蕴含 $\frac{dx'\,dy'}{(y')^2} = \frac{dx\,dy}{y^2}$.
(d). 若 $g = \begin{pmatrix} a & b \\ c & d \end{pmatrix} \in \SL_2(\Z)$, 则 $ad - bc = 1$, 所以 $c$ 和 $d$ 在 $\Z$ 中互素. 反之, 若 $c$ 和 $d$ 在 $\Z$ 中互素, 则存在 $a, b \in \Z$ 使得 $ad - bc = 1$, 然后 $g := \begin{pmatrix} a & b \\ c & d \end{pmatrix} \in \SL_2(\Z)$.
设 $z \in \mathfrak{h}$. 则群 $\Z + \Z z$ 在 $\CC$ 中离散 (因为 $z \notin \R$), 所以它与每个球相交有限. 特别地, $|cz + d|$ ($c, d \in \Z$ 互素) 的下确界达到. 由上一段和 $\Im(g \cdot z) = \frac{\Im(z)}{|cz+d|^2}$ (若 $g = \begin{pmatrix} a & b \\ c & d \end{pmatrix}$) 这一事实, 看到 $\sup\{|\Im(g \cdot z)| \mid g \in \SL_2(\Z)\}$ 也达到. 要证明存在 $g \in \SL_2(\Z)$ 使得 $g \cdot z \in \Omega$. 用某个 $g \in \SL_2(\Z)$ 替换 $z$ 后, 可以假设 $|\Im(z)| = \sup\{|\Im(g \cdot z)| \mid g \in \SL_2(\Z)\}$. 特别地, 取 $g = \begin{pmatrix} 0 & -1 \\ 1 & 0 \end{pmatrix}$, 得到 $|\Im(z)| \geq \frac{|\Im(z)|}{|z|^2}$, 所以 $|z| \geq 1$. 若 $g_n = \begin{pmatrix} 1 & n \\ 0 & 1 \end{pmatrix}$ ($n \in \Z$), 则 $g_n \cdot z = z + n$, 所以 $\Im(g_n \cdot z) = \Im(z)$, 且仍有 $|\Im(g_n \cdot z)| = \sup\{|\Im((hg_n) \cdot z)| \mid h \in \SL_2(\Z)\}$, 特别地仍有 $|g_n \cdot z| \geq 1$. 可以选择 $n \in \Z$ 使得 $-\frac{1}{2} \leq \Re(z) + n \leq \frac{1}{2}$, 然后有 $g_n \cdot z \in \Omega$, 完成了.
(e). 有:
\[
\int_{\Omega}\frac{dx\,dy}{y^2} = \int_{-1/2}^{1/2}\int_{\sqrt{1-x^2}}^{+\infty}\frac{dx\,dy}{y^2} = \int_{-1/2}^{1/2}\frac{1}{\sqrt{1-x^2}}\,dx < +\infty.
\]
(f). 对每个 $\varepsilon \in (0, 1/4)$, 设
\[
\Omega_{\varepsilon} = \{z \in \mathfrak{h} \mid |z| > 1 - \varepsilon \text{ 且 } |\Im(z)| < \frac{1}{2} + \varepsilon\}.
\]
如在 (e) 中, 看到
\[
I_{\varepsilon} := \int_{\Omega_{\varepsilon}}\frac{dx\,dy}{y^2} = \int_{-1/2-\varepsilon}^{1/2+\varepsilon}\frac{1}{\sqrt{1-x^2-\varepsilon}}\,dx < +\infty.
\]
设
\[
\Omega'_{\varepsilon} = \{g \in G \mid g \cdot i \in \Omega_{\varepsilon}\}.
\]
这是 $\Omega_{\varepsilon}$ 在连续映射 $G \to G/K$ 下的原像. 特别地, 它是 $G$ 的开子集, 所以由正则 Borel 测度的定义 (和 Urysohn 引理),
\[
\mu_G(\Omega'_{\varepsilon}) = \sup\Bigl\{\int_G f\,\dd\mu_G \mid f \in \mathcal{C}_c(G),\, 0 \leq f \leq 1 \text{ 且 } \operatorname{supp}(f) \subset \Omega'_{\varepsilon}\Bigr\}.
\]
若 $f \in \mathcal{C}_c(G)$ 满足 $0 \leq f \leq 1$ 且 $\operatorname{supp}(f) \subset \Omega'_{\varepsilon}$, 则对每个 $g \in G$, $f^K(gK) = \int_K f(gk)\,\dd\mu_K(k) \in [0,1]$ (因为 Haar 测度 $\mu_K$ 的总体积为 $1$), 且 $f^K(gK) = 0$ 若 $g \notin \Omega'_{\varepsilon}$, 即 $gK \notin \Omega_{\varepsilon}$; 所以 $\int_{G/K}f^K\,\dd\mu_{G/K} \leq \mu_{G/K}(\Omega_{\varepsilon})$. 由于对每个 $f \in \mathcal{C}_c(G)$ 有 $\int_G f\,\dd\mu_G = \int_{G/K}f^K\,\dd\mu_{G/K}$, 且 $\mu_{G/K}$ 由 (c) 是 $\frac{dx\,dy}{y^2}$ 的倍数, 看到 $\mu_G(\Omega'_{\varepsilon})$ 被 $I_{\varepsilon}$ 的倍数控制, 所以有限.
另一方面, 由于 $\Omega_{\varepsilon} \supset \Omega$, 投影 $\pi \colon G \to \SL_2(\Z)\backslash G$ 到 $\Omega'_{\varepsilon}$ 的限制满射; 事实上, 若 $g \in G$, 由 (d), 存在 $h \in \SL_2(\Z)$ 使得 $hg \cdot i \in \Omega$, 然后 $hg \in \Omega'_{\varepsilon}$ 且 $\pi(hg) = \pi(g)$. 再次由正则 Borel 测度的定义和 Urysohn 引理, 存在 $f \in \mathcal{C}_c^+(G)$ 使得 $f|_{\Omega'_{\varepsilon}} = 1$ 且 $\int_G f\,\dd\mu_G < +\infty$. 设 $h = \sum_{\gamma \in \SL_2(\Z)}f \in \mathcal{C}_c(\SL_2(\Z)\backslash G)$. 有 $\int_{\SL_2(\Z)\backslash G}h\,\dd\mu_{\SL_2(\Z)\backslash G} = \int_G f\,\dd\mu_G < +\infty$. 但若 $x \in \SL_2(\Z)\backslash G$, $g$ 是满足 $\pi(g) = x$ 的 $\Omega'_{\varepsilon}$ 的元素, 则
\[
h(x) = \sum_{\gamma \in \SL_2(\Z)}f(\gamma g) \geq f(g) = 1
\]
(用 $f \geq 0$ 这一事实). 所以
\[
\operatorname{vol}(\SL_2(\Z)\backslash G) \leq \int_{\SL_2(\Z)\backslash G}h\,\dd\mu_{\SL_2(\Z)\backslash G} < +\infty.
\]
(5) 设 $V := \{f \in L^2(X) \mid \int_X f(x)\,\dd\mu(x) = 0\}$. 这是 $L^2(X)$ 的 $G$-不变子空间, 因为测度 $\mu$ 被 $G$ 左不变. 证明 $V^G = 0$. 设 $f \in L^2(X)^G$. 则对每个 $g \in G$ 有 $\|R_gf - f\|_{L^2(X)} = 0$. 把 Minkowski 不等式 (见习题 \ref{ex:I.5.6.7}) 应用于函数 $\phi \colon G \times X \to \CC$, $(g, x) \longmapsto |f(xg) - f(x)|$, 看到
\[
\int_X\Bigl(\int_G|f(xg) - f(x)|\,\dd g\Bigr)^2\,\dd x = 0,
\]
所以函数 $x \longmapsto \int_G|f(xg) - f(x)|\,\dd g$ 在 $x$ 上几乎处处为 $0$. 选择满足 $\int_G|f(x_0g) - f(x_0)|\,\dd g = 0$ 的 $x_0 \in X$. 则函数 $g \longmapsto f(x_0g) - f(x_0)$ 在 $G$ 上几乎处处为 $0$, 所以 $f$ 在 $L^2(X)$ 中等于常值函数 $f(x_0)$ 的类. 若此外 $f \in V$, 这迫使 $f = 0$.
现在把问题 (4) 应用于 $V$. 它的结论说对所有 $f_1, f_2 \in V$, $\lim_{g \to +\infty}\langle R_gf_1, f_2 \rangle = 0$.
设 $f_1, f_2 \in L^2(X)$. 写 $c_1 = \int_X f_1\,\dd\mu$, $c_2 = \int_X f_2\,\dd\mu$. 由于 $\operatorname{vol}(X) = 1$, 函数 $f_1 - c_1$ 和 $f_2 - c_2$ 在 $V$ 中, 所以由刚才看到的, $\lim_{g \to +\infty}\langle R_gf_1 - c_1, f_2 - c_2 \rangle = 0$.
对每个 $g \in G$,
\begin{align*}
\langle R_gf_1 - c_1, f_2 - c_2 \rangle &= \int_X (f_1(xg) - c_1)(f_2(x) - c_2)\,\dd x \\
&= \int_X f_1(xg)f_2(x)\,\dd x - c_1\int_X f_2(x)\,\dd x - c_2\int_X f_1(xg)\,\dd x + c_1c_2 \\
&= \int_X f_1(xg)f_2(x)\,\dd x - c_1c_2.
\end{align*}
最后得到
\[
\lim_{g \to +\infty}\int_X f_1(xg)f_2(x)\,\dd x = c_1c_2,
\]
如所愿.
(6)(a). 定义函数 $h \colon G \to \CC$:
\[
h(x) = \int_K\psi(xy^{-1})\,\dd\mu_K(y),
\]
对每个 $x \in G$. 证明 $h$ 连续. 设 $\varepsilon > 0$. 由于 $\psi \in \mathcal{C}_c(G)$, 命题 \ref{prop:I.1.12} 蕴含 $\psi$ 左一致连续, 所以存在 $G$ 中 $1$ 的邻域 $U$ 使得对每个 $x \in G$ 和每个 $x_0 \in U$, $|\psi(x_0x) - \psi(x)| \leq \varepsilon$. 则若 $x \in G$, $x_0 \in U$, 有
\[
|h(x_0x) - h(x)| \leq \int_K|\psi(x_0xy^{-1}) - \psi(xy^{-1})|\,\dd\mu_K(y) \leq \int_K\varepsilon\,\dd\mu_K = \varepsilon,
\]
所以 $h$ 也左一致连续 (特别地连续). 又, 由于 $\psi(xy^{-1}) = 0$ 除非 $x \in (\operatorname{supp}\psi)y \subset (\operatorname{supp}\psi)K$, $\psi$ 的支撑包含于 $(\operatorname{supp}\psi)K$, 所以紧.
最后, 对每个 $f \in \mathcal{C}_c(G)$:
\begin{align*}
\psi * \mu_K(f) &= \int_{G \times K}f(xy)\psi(x)\,\dd\mu_G(x)\,\dd\mu_K(y) \\
&= \int_{G \times K}f(x)\psi(xy^{-1})\,\dd\mu_G(x)\,\dd\mu_K(y) \quad \text{(因为 } \Delta_G|_K = 1\text{)} \\
&= \int_G f(x)h(x)\,\dd\mu_G(x).
\end{align*}
(b). 设 $(U_n)_{n \geq 0}$ 是 $G$ 中 $1$ 的递减邻域序列, 构成邻域基; 例如可以取 $U_n$ 为 $G$ 与 $M_2(\R)$ 中半径 $2^{-n}$ 的球的交, 对 $M_2(\R)$ 上任何范数. 由命题 \ref{prop:I.4.1.8}, 存在近似单位元 $(\psi_n)_{n \geq 0}$ 使得对每个 $n$ 有 $\operatorname{supp}(\psi_n) \subset U_n$.
像命题 \ref{prop:I.4.1.9} 中那样证明序列 $(\psi_n)_{n \geq 0}$ 可行. 设 $f \in \mathcal{C}_c(G)$. 则对每个 $n$, $\mu_K(f) = \int_K f(y)\,\dd\mu_K(y) = \int_{G \times K}f(y)\psi_n(x)\,\dd\mu_G(x)\,\dd\mu_K(y)$, 所以
\[
|\psi_n * \mu_K(f) - \mu_K(f)| \leq \int_{G \times K}|f(xy) - f(y)|\psi(x)\,\dd\mu_G(x)\,\dd\mu_K(y) \leq \sup_{x \in U_n}\|L_{x^{-1}}f - f\|_{\infty}.
\]
由于 $f$ 左一致连续, 当 $n \to +\infty$ 时这趋于 $0$.
(c). 由 (b), 只需对 $f \in \mathcal{C}_c(G)$ 和泛函 $\psi * \mu_K$ ($\psi \in \mathcal{C}_c(G)$) 证明 $(***)$ 的类似物, 即证明
\[
\lim_{g \to \infty}\int_X \psi * \mu_K(\cdot g)f\,\dd x = \int_X f\,\dd\mu
\]
(其中把 $\psi * \mu_K$ 与 $L^2(X)$ 的元素等同). 用 (a) 的显式 $h$, 这由问题 (5) 的等分布结果 ($\psi * \mu_K$ 现在是 $\mathcal{C}_c(G)$ 的元素 $h$) 得出, 因为 $Yg$ 在 $X$ 中的像正是 $y \cdot g$ ($y \in Y$), 而 $h * f$ 的积分极限由 (5) 给出. 详细说来, 由 (5) 应用于 $f_1 = h$ (通过 $\pi$ 下降) 和 $f_2 = f$, 看到 $\lim_{g \to \infty}\int_X h(xg)f(x)\,\dd x = \int_X h\,\dd\mu \int_X f\,\dd\mu$. 而 $\int_X h\,\dd\mu = \mu_K$ 的规范化总质量 $= 1$, 所以得到 $\lim_{g \to \infty}\int_X h(xg)f(x)\,\dd x = \int_X f\,\dd\mu$, 这正是 $(***)$ (把 $f$ 换成 $\tilde{f}$, 用 Fubini 和 $h$ 的定义把 $y$-积分与 $x$-积分交换).
\end{solution}
